% !TEX root = ../../tree_notes.tex
% =====================================================================
%  Chapter 9 (map v2) -- Cartesian Trees and the Monotonic Stack
%  Status: WRITTEN.  (Folder keeps its ch26_cartesian_trees name; it
%   prints as Chapter 9, label prefix c26, build = --only 26.)
%
%  Files, in reading order:
%    00_opener.tex              glance + badges, three reading paths, ladder
%                              staircase, the two running examples (the array
%                              A = 3,1,4,2 and the histogram H = 2,1,5,6,2,3),
%                              "who owns what" comparison, roadmap, recognition
%                              guide, constraint table, notation
%    01_cartesian_build.tex     9.1 The Cartesian Tree in O(n)     sec:c26-cartesian
%                              (solved LC 1793; practice LC 907)
%    02_monotonic_stack.tex     9.2 Monotonic Stack Applications   sec:c26-monostack
%                              (solved LC 84; practices LC 85, 42, 239, 739)
%    03_rmq_via_cartesian.tex   9.3 Range Minimum Query            sec:c26-rmq
%                              (solved CSES 1647)
%    04_random_treap.tex        9.4 The Cartesian Tree as a Treap  sec:c26-treapview
%                              (solved LC 998)
%    99_chapter_end.tex         the 15-part revision contract + Summit (CF 817D)
%
%  This chapter OWNS "read a tree straight out of an array with ONE stack."
%  Chapter 8 ended with balance by construction -- you edit a tree and pay a
%  rotation to keep it short. Here the tree is not edited at all: it is
%  FORCED out of a static array by two simultaneous orders (inorder = index,
%  heap = value), and the O(n) monotonic stack that builds it is the same one
%  move that solves largest-rectangle, subarray minimums, next-greater, sliding
%  extremes, RMQ, and the Summit. One data structure, one stack, a whole family.
%
%  Anti-repetition -- deliberately NOT re-taught here:
%    * treaps, random priorities, expected height, split/merge
%        -> Chapter 8 (\chref{c06}); 9.4 uses the Cartesian<->treap EQUIVALENCE
%           (thm:c06-exptheight) as a black box and never rebuilds treap theory
%    * the BST invariant, search / insert / delete, inorder = sorted
%        -> Chapter 7 (\chref{c05}); we lean on inorder, never re-derive it
%    * heaps / priority_queue and the median-from-stream engines
%        -> Chapter 10 (\chref{c07}); the heap ORDER here is an invariant of a
%           tree built by a stack, not a priority_queue, and is not re-taught
%    * the three DFS orders / traversals
%        -> Chapter 3 (\chref{c03})
%    * LCA machinery (Euler tour, binary lifting, RMQ-on-depth)
%        -> Chapter 16 (\chref{c11}); 9.3 does the RMQ<->LCA REDUCTION and hands
%           the LCA engine forward, never re-derives it
%    * sparse tables and segment trees as query engines
%        -> Chapter 19 (\chref{c13}) / Chapter 18 (\chref{c14}); forward-ref only
%    * sliding-window MAXIMUM via a deque is OWNED here; sliding-window MEDIAN
%      (the multiset / two-heap version) stays in Chapter 8 (\chref{c06}).
%
%  NOTE ON MIGRATION: the Cartesian / monotonic-stack material moved OUT of the
%  old ch06_balanced_trees topic 4 into this new Chapter 9 per
%  structure/topics_by_chapter.txt; it is intentionally absent from Chapter 8.
%  NOTE ON A PRIOR MISLABEL: the old ch06.tex header listed the CF 817D summit
%  as "Showmen". CF 817D is in fact "Imbalanced Array" (sum over all subsegments
%  of max - min). This chapter's Summit uses the correct problem.
% =====================================================================
\bookchapter{Cartesian Trees and the Monotonic Stack}{c26}
\chapterinfo{LC Medium--Hard, CF 1900, CSES}{Written}

% !TEX root = ../../tree_notes.tex
% =====================================================================
%  Chapter 9 opener -- "Chapter at a Glance", three reading paths + the
%  ladder staircase, the two running examples (the array A = 3,1,4,2 and
%  the histogram H = 2,1,5,6,2,3), the "who owns what" comparison
%  (monotonic stack vs segment tree vs treap), roadmap, recognition guide,
%  constraint->tool table, notation.
%  NO numbered boxes here. Exports fig:c26-hist and fig:c26-cartA.
% =====================================================================

\begin{chaptertheorybox}{Chapter at a Glance}
Chapter~8 (\chref{c06}) made a tree short \emph{by editing it}: you insert or delete, then pay a
rotation to repair the height. This chapter never edits anything. It starts from a bare
\term{array} and asks a different question -- \emph{what tree is already hiding inside these
numbers?} Hand it an array $a$ and there is exactly \emph{one} binary tree whose inorder reading is
the array in order and whose heap rule is the values: the \term{Cartesian tree}. And there is one
tiny machine that reads that tree straight out of the array in a single left-to-right pass -- a
\term{monotonic stack}. The astonishing part is that this \emph{same} stack, with no tree built at
all, is the whole solution to largest rectangle in a histogram, sum of subarray minimums, next
greater element, sliding-window maximum, range minimum query, and the chapter's Summit. One
invariant, one stack, an entire problem family:
\begin{itemize}[nosep]
  \item \textbf{Two orders, one tree.} A Cartesian tree obeys \emph{BST order on the index} (inorder
        $=$ the array) and \emph{heap order on the value} at the same time. Under distinct values the
        tree is unique, and its root is just the position of the minimum. \hfill \lv{3}
  \item \textbf{The stack is the build.} Sweep the array, keep the stack values increasing, and pop
        while the new element breaks monotonicity -- the popped links \emph{are} the Cartesian tree's
        parent pointers, produced in $\bigO{n}$. The moment a bar stops being the smallest is exactly
        the moment its rectangle or span is decided. \hfill \lv{3}--\lv{5}
  \item \textbf{Span $=$ subtree.} The subtree of node $i$ is precisely the maximal window in which
        $a[i]$ is the minimum. That one lemma converts a hundred "sum of $a[i]\cdot(\text{how many
        ranges where it wins})$'' problems into two stack passes for the prev/next smaller (or
        greater) bounds. \hfill \lv{4}--\lv{5}
\end{itemize}

\textbf{You will learn to}
\begin{itemize}[nosep]
  \item build a Cartesian tree in $\bigO{n}$ with a stack and read the two invariants off it; \hfill\lv{3}
  \item recognise "the answer is decided when the stack pops" and turn histogram, subarray-minimum
        and rain-water problems into one or two stack passes; \hfill\lv{3}--\lv{5}
  \item keep sliding-window extremes with a monotonic \emph{deque}; \hfill\lv{4}
  \item reduce a static \term{range minimum query} to an LCA on the Cartesian tree, and know exactly
        what you are handing off to the LCA machinery of Chapter~16 (\chref{c11}); \hfill\lv{3}--\lv{4}
  \item see a Cartesian tree as a \emph{treap whose priorities are the values} and a random-permutation
        BST as a random Cartesian tree -- borrowing Chapter~8's height bound, not re-proving it. \hfill\lv{4}
\end{itemize}

\textbf{Prerequisites.} Arrays, and the \term{stack} as a container, from Chapter~1--2; the BST
inorder fact from Chapter~7 (\chref{c05}) is \emph{leaned on}, never re-derived. Chapter~8's treap
(\secref{c06-treaps}) and its expected-height theorem (\thmref{c06-exptheight}) are used as a black
box only in \secref{c26-treapview}. LCA is quoted, not built (Chapter~16, \chref{c11}).
\end{chaptertheorybox}

% ---------------- Reading paths + ladder position ----------------
\begin{chaptertheorybox}{How to Read This Chapter (three paths)}
Level badges $\lv1$--$\lv6$: $\lv3$ is a solid LC-medium stack problem, $\lv4$ is a Hard you can only
crack once the "the span is the subtree" idea clicks, and $\lv5$ is an LC/CF Hard where the whole
difficulty is \emph{realising} the monotonic structure is there at all.

\textbf{Fast path ($\lv3$, ``I just want the stack'').} Read \secref{c26-monostack} directly -- the
monotonic stack is a standalone technique you can use without ever naming a Cartesian tree. Do the
solved histogram (LC 84) and the next-greater / sliding-window practices. Skip the tree build, the
RMQ reduction and the treap view. You leave able to cold-solve largest-rectangle, daily-temperatures
and sliding-window-maximum.

\textbf{Structural path ($\lv3$--$\lv4$).} Start at \secref{c26-cartesian} and \emph{build} the tree
-- it makes \secref{c26-monostack} feel inevitable rather than clever, and pays off in
\secref{c26-rmq} where the tree is the point. Recommended order: 9.1 $\to$ 9.2 $\to$ 9.3.

\textbf{Full path.} Everything in order, finishing at \secref{c26-treapview}, where the Cartesian tree
and the Chapter~8 treap turn out to be the \emph{same object}, and at the Summit (CF 817D), which is
the purest expression of the span-equals-contribution idea: the answer is a signed sum of
$a[i]\cdot(\text{span where it is max}) - a[i]\cdot(\text{span where it is min})$, both from a stack
pass each.

The forward rope is Chapter~10 (\chref{c07}) for genuine \emph{heaps} (a priority queue you edit, not
a heap-order you read off a static array), and Chapter~16 (\chref{c11}) for the LCA engine the RMQ
reduction calls.
\end{chaptertheorybox}

\begin{figure}[H]
\centering
\begin{tikzpicture}[diagram, x=9mm, y=7mm]
  \foreach \st/\lab in {1/lv1,2/lv2,3/lv3,4/lv4,5/lv5,6/lv6} {
    \fill[\lab!22, draw=\lab, line width=0.8pt] (\st-1,0) rectangle (\st,\st*0.52);
    \node[font=\scriptsize\bfseries, text=\lab] at (\st-0.5,\st*0.52-0.42) {L\st};}
  \node[annot, anchor=south west, font=\scriptsize]
    at (0.05,0.62) {read an array as a tree};
  \node[annot, anchor=south west, font=\scriptsize]
    at (1.05,1.14) {9.1 build; 9.2 stack basics};
  \node[annot, anchor=south west, font=\scriptsize]
    at (2.05,1.66) {9.2 histogram, next-greater, deque};
  \node[annot, anchor=south west, font=\scriptsize]
    at (3.05,2.18) {9.1/9.3 span=subtree, RMQ};
  \node[annot, anchor=south west, font=\scriptsize, text=lv5]
    at (4.05,2.70) {9.2 LC 84; 9.3 RMQ lab; treap view};
  \node[annot, anchor=south west, font=\scriptsize, text=lv5]
    at (5.05,3.22) {summit (CF 817D)};
\end{tikzpicture}
\caption*{Chapter~9 on the six-step ladder: seeing the array$\to$tree reading is $\lv3$; the monotonic
stack and histogram sit at $\lv3$--$\lv4$; the ``span $=$ subtree'' contribution trick and the RMQ
reduction climb to $\lv4$; solving LC 84 and the treap-equivalence view reach $\lv5$, where the Summit
(CF 817D \emph{Imbalanced Array}) turns the whole chapter into two signed stack passes on $n\le 10^6$.
Fast path $=$ steps 2--3; full path runs to 6.}
\end{figure}

% ---------------- The running examples ----------------
\subsection*{The Running Examples: One Array, One Histogram}

Two small inputs, reused in every section so the ideas always land on something concrete.

\textbf{Array $A$} $= (3, 1, 4, 2)$, positions $1,2,3,4$. Its values are \emph{distinct}, so its
min-Cartesian tree is unique: the smallest value, $1$ at position $2$, is the root; the elements to
its left form its left subtree and those to its right form its right subtree, recursively. Read the
tree's inorder and you get $A$ back; read any node against its children and the parent is smaller --
both orders at once. $A$ drives the build (\secref{c26-cartesian}), the next-smaller spans, the RMQ
queries (\secref{c26-rmq}) and the treap equivalence (\secref{c26-treapview}).

\textbf{Histogram $H$} $= (2, 1, 5, 6, 2, 3)$ -- the official LeetCode~84 input, whose answer is the
rectangle of height $5$ and width $2$ over bars $3,4$, giving area $\boxed{10}$. $H$ has a
\emph{duplicate} ($2$ appears twice), so it also shows why tie-breaking must be asymmetric when spans
are counted. It drives the monotonic-stack section (\secref{c26-monostack}) and reappears in the
Summit's contribution argument.

\begin{figure}[H]
\centering
\begin{tikzpicture}[diagram, x=11mm, y=6mm]
  % histogram H = 2,1,5,6,2,3
  \def\hs{{2,1,5,6,2,3}}
  \draw[thick, darkgray] (0,0) -- (6.6,0);
  \foreach \i/\h in {0/2,1/1,2/5,3/6,4/2,5/3} {
    \fill[primaryblue!22, draw=primaryblue, thick] (\i,0) rectangle (\i+1,\h);
    \node[font=\scriptsize, text=darkgray] at (\i+0.5,-0.45) {\i+1};
    \node[font=\scriptsize\bfseries, text=deepblue] at (\i+0.5,\h+0.35) {\h};}
  % the winning rectangle height 5 over bars 3,4 (x=2..4)
  \draw[ultra thick, warningred, dashed] (2,0) rectangle (4,5);
  \node[font=\scriptsize\bfseries, text=warningred, align=center] at (3,2.6) {$5\times 2$\\ area $10$};
  \node[annot, align=left, font=\scriptsize] at (7.0,3) {\textbf{Histogram $H$}\\ $2,1,5,6,2,3$\\
    the tallest rectangle is\\ the dashed one, area $10$};
\end{tikzpicture}
\caption{The histogram running example $H=(2,1,5,6,2,3)$. Bars are drawn in index order; the
monotonic stack (\secref{c26-monostack}) finds that the widest rectangle of height $5$ spans bars
$3$--$4$ and has area $10$ -- the LeetCode~84 answer.}\label{fig:c26-hist}
\end{figure}

\begin{figure}[H]
\centering
\begin{tikzpicture}[diagram, scale=0.95, transform shape]
  % min-Cartesian tree of A = 3,1,4,2 (index shown as subscript)
  \node[troot] (a2) at (0,0)   {$a_2{=}1$};
  \node[tleaf] (a1) at (-2,-1.4) {$a_1{=}3$};
  \node[tnode] (a4) at (2,-1.4)  {$a_4{=}2$};
  \node[tleaf] (a3) at (1,-2.8)  {$a_3{=}4$};
  \draw[tedgehl] (a2) -- (a1);
  \draw[tedgehl] (a2) -- (a4);
  \draw[tedgehl] (a4) -- (a3);
  \node[annot, align=left, font=\scriptsize] at (-4.6,-1.4)
    {\textbf{Inorder} (left,root,right):\\ $a_1,a_2,a_3,a_4$\\ $=3,1,4,2$ -- back to $A$};
  \node[annot, align=left, font=\scriptsize] at (5.0,-2.0)
    {\textbf{Min-heap on value:}\\ parent $<$ both children\\ $1<3,1<2<4$};
\end{tikzpicture}
\caption{The min-Cartesian tree of $A=(3,1,4,2)$. The root is the minimum ($1$, at position $2$); the
left part ($a_1$) and right part ($a_3,a_4$) recurse. Its inorder is the array itself, and every
parent is smaller than its children -- the two invariants of \defref{c26-cartesian}, shown on one
object. Node $a_4$'s subtree is the single window $\{3,4\}$ where $2$ is the minimum: the
span-equals-subtree idea of \lemref{c26-span}.}\label{fig:c26-cartA}
\end{figure}

\medskip
\noindent Three engines answer the \emph{same} range-extremum questions; knowing which owns what stops
you re-solving a solved problem:
\begin{center}
\renewcommand{\arraystretch}{1.25}
\setlength{\tabcolsep}{5pt}
\begin{tabular}{@{}p{0.185\linewidth} p{0.42\linewidth} p{0.31\linewidth}@{}}
\toprule
\textbf{Engine} & \textbf{What it is for} & \textbf{Where it lives} \\
\midrule
Monotonic stack / deque & \emph{offline / one-pass} questions where the answer is decided the moment a
  bar is popped: next greater, histogram spans, sliding-window extreme & this chapter (\secref{c26-monostack}) \\
Segment / sparse table  & \emph{arbitrary} online range queries (min, sum, \dots) over a
  \emph{changing} array & Chapter~18--19 (\chref{c14}, \chref{c13}) \\
Treap (heap on random priority) & a \emph{dynamic} ordered set you insert into and delete from while
  staying short & Chapter~8 (\chref{c06}) \\
\bottomrule
\end{tabular}
\end{center}

% ---------------- Roadmap diagram ----------------
\begin{figure}[H]
\centering
\begin{tikzpicture}[diagram, node distance=3.5mm,
    top/.style={boxnode, fill=primaryblue!10, draw=primaryblue, text width=19mm,
                minimum height=11mm, font=\small},
    res/.style={annot, text width=19mm, anchor=north, font=\scriptsize}]
  \node[top] (t1) {\textbf{9.1} Cartesian Build};
  \node[top, right=of t1] (t2) {\textbf{9.2} Monotonic Stack};
  \node[top, right=of t2] (t3) {\textbf{9.3} RMQ via LCA};
  \node[top, right=of t3] (t4) {\textbf{9.4} Random Treap View};
  \foreach \a/\b in {t1/t2, t2/t3, t3/t4} \draw[flow] (\a) -- (\b);
  \node[res] at ([yshift=-2mm]t1.south) {array $\to$ tree; the two orders};
  \node[res] at ([yshift=-2mm]t2.south) {stack \emph{is} the build; histogram};
  \node[res] at ([yshift=-2mm]t3.south) {range min $=$ LCA value};
  \node[res] at ([yshift=-2mm]t4.south) {Cartesian $=$ treap; random BST};
\end{tikzpicture}
\caption*{The four topics are one picture seen four ways. 9.1 \emph{builds} the tree and names its two
invariants; 9.2 reveals the build is just a monotonic stack and runs the whole application family on
that stack; 9.3 uses the tree as an \emph{index} and reduces range minimum to LCA; 9.4 steps back and
shows the tree is secretly the Chapter~8 treap with values-as-priorities -- tying the chapter's ending
to where the balanced-tree story closed.}
\end{figure}

% ---------------- Topic table ----------------
\begin{center}
\renewcommand{\arraystretch}{1.25}
\setlength{\tabcolsep}{4pt}
\begin{tabular}{@{}p{0.235\linewidth} p{0.365\linewidth} c p{0.23\linewidth}@{}}
\toprule
\textbf{Topic} & \textbf{Key results} & \textbf{Levels} & \textbf{Solved problem} \\
\midrule
\secref{c26-cartesian}\ \ Cartesian Tree in $O(n)$
  & two orders (inorder=index, heap=value); unique on distinct values; stack build & $\lv3$--\lv{4}
  & \probref{c26-goodsub} \\
\addlinespace
\secref{c26-monostack}\ \ Monotonic Stack Applications
  & pop $=$ decision; histogram, next-greater, deque sliding max, rain water & $\lv3$--\lv{5}
  & \probref{c26-hist84} \\
\addlinespace
\secref{c26-rmq}\ \ Range Minimum Query via the Tree
  & RMQ of $a[l..r]$ $=$ value at $\lca(l,r)$ in the Cartesian tree & $\lv3$--\lv{4}
  & \probref{c26-rmq1647} \\
\addlinespace
\secref{c26-treapview}\ \ The Cartesian Tree as a Treap
  & values $=$ priorities $\Rightarrow$ treap; permutation $\Rightarrow$ random BST & $\lv4$
  & \probref{c26-mbt2} \\
\bottomrule
\end{tabular}
\end{center}

% ---------------- Problem recognition guide ----------------
\subsection*{Problem Recognition Guide --- When You See \emph{this}, Think \emph{that}}

Monotonic-stack problems are the ones where an element's fate is sealed by the first \emph{smaller}
(or \emph{larger}) neighbour to its left or right. If the phrasing mentions a threshold, a window, or
"extends until", that is the tell; every row points at the section that builds the tool.

\begin{center}
\renewcommand{\arraystretch}{1.3}
\setlength{\tabcolsep}{4pt}
\begin{tabular}{@{}p{0.44\linewidth} p{0.34\linewidth} l@{}}
\toprule
\textbf{When the problem says\ldots} & \textbf{\textcolor{deepblue}{\textbf{Think\ldots}}} & \textbf{Section} \\
\midrule
"largest rectangle in a histogram / maximal binary grid"
  & pop when the bar stops being a minimum; width fixed at pop & \secref{c26-monostack} \\
\addlinespace
"sum over all subarrays of the min (or max)" / "how many ranges is $a[i]$ the winner of"
  & prev/next smaller (or greater) span; $a[i]\cdot(i-L)\cdot(R-i)$ & \secref{c26-cartesian}, Summit \\
\addlinespace
"next greater element", "daily temperatures", "can you see a taller building"
  & decreasing stack; pop when a bigger value arrives & \secref{c26-monostack} \\
\addlinespace
"max (or min) inside a \emph{sliding window} of size $k$"
  & monotonic \emph{deque}; drop from both ends & \secref{c26-monostack} \\
\addlinespace
"minimum of a static range $[l,r]$, many queries"
  & value at $\lca(l,r)$ in the Cartesian tree & \secref{c26-rmq} \\
\addlinespace
"rebuild a max-tree / insert at the far right"
  & the Cartesian / treap view; walk the right spine & \secref{c26-treapview} \\
\bottomrule
\end{tabular}
\end{center}

\begin{center}
\renewcommand{\arraystretch}{1.25}
\setlength{\tabcolsep}{4pt}
\begin{tabular}{@{}p{0.29\linewidth} p{0.42\linewidth} p{0.20\linewidth}@{}}
\toprule
\textbf{Constraint you see} & \textbf{What it forces} & \textbf{Chapter-9 tool} \\
\midrule
$n \le 10^{5}$--$10^{6}$, "sum over all subarrays" & $\bigO{n}$ -- an $\bigO{n^2}$ window scan dies & stack span \\
\addlinespace
answer decided by nearest greater / smaller neighbour & compute all bounds in one pass each & \secref{c26-monostack} \\
\addlinespace
window of fixed width $k$, extreme per window & monotonic deque, amortised $\bigO{1}$ per step & \secref{c26-monostack} \\
\addlinespace
static array, many range-min queries, few updates & Cartesian tree $+$ LCA (or sparse table) & \secref{c26-rmq} \\
\addlinespace
products that overflow $32$ bits ($a[i]\cdot\text{counts}$) & accumulate in \code{long long} & Summit (\probref{c26-summit}) \\
\bottomrule
\end{tabular}
\end{center}

\begin{notebox}[Not this chapter?]
If the structure must \emph{change} -- insert, delete, then ask a range question online -- a one-pass
stack cannot remember your edits; that is the segment tree (\chref{c13}) or a balanced tree
(\chref{c06}). If you only need the \emph{global} largest or smallest over the whole array and nothing
per-window, a single variable suffices -- no stack. If you want the sliding-window \emph{median}, that
multiset / two-heap version is Chapter~8 (\chref{c06}); this chapter owns the sliding
\emph{extreme} via the deque. And if the tree must be \emph{balanced against an adversary}, you need
rotations -- Chapter~8 again. This chapter owns \emph{reading a static array as a tree, and the one
stack that does all the work}.
\end{notebox}

% ---------------- Notation ----------------
\subsection*{Notation Used in This Chapter}

\begin{center}
\renewcommand{\arraystretch}{1.2}
\setlength{\tabcolsep}{4pt}
\begin{tabular}{@{}p{0.20\linewidth} p{0.27\linewidth}@{\quad}p{0.17\linewidth} p{0.27\linewidth}@{}}
\toprule
\textbf{Symbol} & \textbf{Meaning} & \textbf{Symbol} & \textbf{Meaning} \\
\midrule
$A$                 & the running array $(3,1,4,2)$ (\figref{c26-cartA}) & $H$ & the running histogram $(2,1,5,6,2,3)$ (\figref{c26-hist}) \\
$a[i]$              & value at index $i$; indices are $1$-based here & $\mathit{stack}$ & the monotonic stack of indices \\
$L_i,\ R_i$         & nearest index to the left/right of $i$ that \emph{beats} $a[i]$ (prev/next smaller or greater) & $\mathit{span}_i$ & $(L_i, R_i)$, the window where $a[i]$ is the extremum \\
$\lca(u,v)$         & lowest common ancestor in a tree (\chref{c11}); here in the Cartesian tree & $T$ & the Cartesian tree of an array \\
$\text{inorder}(T)$ & the array read off the tree by an inorder walk & $\bigstar$ & marks the Summit \\
\bottomrule
\end{tabular}
\end{center}

% !TEX root = ../../tree_notes.tex
% =====================================================================
%  9.1  The Cartesian Tree in O(n)              sec:c26-cartesian
%  Exports: def:c26-cartesian  thm:c26-unique  lem:c26-span
%           fig:c26-build  lst:c26-stackbuild  ex:c26-cart-run
%           prob:c26-goodsub  lst:c26-goodsub-soln   (practice: LC 907)
% =====================================================================
\section{The Cartesian Tree in $O(n)$}\label{sec:c26-cartesian}

\begin{previouslybox}
\begin{enumerate}[nosep, leftmargin=1.6em]
  \item A binary tree's \emph{inorder} walk visits \text{left}, node, \text{right}. On a BST it lists
        the keys sorted. Which property of the tree makes that true? (\chref{c03}, \thmref{c05-inorder})
  \item A \emph{min-heap} stores its smallest element where, and what is the parent--child rule?
        (at the root; every parent $\le$ its children -- the idea Chapter~10 builds, \chref{c07})
  \item You have a \code{stack}. Pushing and popping a sequence once costs what in total, and why can
        an inner \code{while}-pop loop still be $\bigO{n}$ overall? (each element pushed/popped once)
  \item Given an array, its minimum sits at some position $p$. What two pieces does $p$ split the
        array into? (the strictly-left part $a[1..p-1]$ and strictly-right part $a[p+1..n]$)
\end{enumerate}
\end{previouslybox}

\begin{topicintro}
There is a tree sleeping inside every array. To wake it, impose \emph{two} ordering rules at once: the
tree's inorder reading must be the array left-to-right (so the \emph{index} behaves like a BST key),
and the tree must be a heap on the array's \emph{values} (so a parent beats its children). A structure
that is simultaneously a BST on position and a heap on value is the \term{Cartesian tree}. On distinct
values it is unique (\thmref{c26-unique}), and its root is nothing but the position of the minimum.

Two facts make it worth your attention. First, it is \emph{buildable in }$\bigO{n}$ by a single
monotonic-stack sweep -- you never touch the $O(n\log n)$ recursive-minimum route. Second, and this is
the chapter's beating heart: the subtree hanging under position $i$ is \emph{exactly} the maximal
window over which $a[i]$ is the minimum (\lemref{c26-span}). Every "for each element, how many ranges
does it win?" problem is therefore already solved once you can read a subtree, and every such problem
becomes two stack passes. We build the tree, prove it unique, state the span lemma, and cash it in on
LeetCode 1793 -- find the best subarray through a fixed index $k$.
\textbf{Prerequisites.} Arrays and the $\bigO{n}$ amortised stack discipline; the inorder walk from
Chapter~3 (\chref{c03}) as a black box.

\textbf{After this section you can}
\begin{itemize}[nosep]
  \item write down the two invariants defining a Cartesian tree and build one on a small array by hand; \hfill\lv{3}
  \item prove the tree is unique when values are distinct, and explain what breaks with ties; \hfill\lv{3}--\lv{4}
  \item build it in $\bigO{n}$ with a monotonic stack and read off every node's parent; \hfill\lv{4}
  \item use the span-equals-subtree lemma to turn a "max over ranges containing $k$" into per-index spans. \hfill\lv{4}--\lv{5}
\end{itemize}
\end{topicintro}

\begin{recognitionbox}[Pattern card -- the array \emph{is} a tree]
\reckey{a question about \emph{each element's dominance window} -- where it is the min/max before a
smaller/larger element appears -- or a \emph{fixed index that every candidate range must contain}.}
\recuse{
  \item "largest rectangle through each bar", "sum of subarray minimums", "how many ranges have $a[i]$
        as their smallest" -- per-index spans, \lemref{c26-span};
  \item "subarray with best $\min \cdot \text{length}$ that must contain $k$" -- \probref{c26-goodsub};
  \item "rebuild a max-tree from an array" -- the same object with the heap rule flipped to max,
        \secref{c26-treapview}.
}
\recnotuse{
  \item the array is \emph{changing} between queries -- you need a segment tree (\chref{c13}), not a
        tree read off a frozen snapshot;
  \item you need the \emph{global} min/max only -- one variable, no structure;
  \item you want a \emph{balanced} tree you can insert into for $\bigO{\log n}$ -- that is Chapter~8's
        treap (\chref{c06}); here the shape is dictated by values, not by you.
}
One-line test: \emph{can the array be read as ``a BST by index, a heap by value'' at once? Then its
Cartesian tree, and the span lemma, are waiting to be used.}
\end{recognitionbox}

\subsection{Intuition}

Take the running array $A = (3,1,4,2)$ from the opener. Ask: \emph{which element should be the root of
a tree that is a heap on values?} A heap forces the smallest element to the top, so the root must be
$1$, sitting at position $2$. Now the BST-by-index rule says everything with index $<2$ goes to the
left of the root and everything with index $>2$ goes to its right -- so $a_1=3$ becomes the entire left
subtree, and the pair $a_3=4,a_4=2$ becomes the right subtree. Inside that right subtree the smallest is
$2$ (position $4$), so it is the local root and $4$ (position $3$, index $<4$) lands on its \emph{left}.
That is the whole tree in \figref{c26-cartA}, and it is the same figure \figref{c26-build} builds step
by step: \emph{root $=$ the minimum, then recurse on the two halves the minimum splits off}.

The remarkable coincidence -- and the seed of every application in this chapter -- is what the
subtrees \emph{mean}. Node $2$ (position $4$) owns a subtree covering just positions $3$--$4$, which is
precisely the widest window $(2,5)$ around position $4$ in which the value $2$ is the minimum: to its
left, position $2$ holds $1<2$ and blocks it; to its right there is nothing, so it runs to the end. The
subtree is the dominance window. Read that backwards and building the tree is the same act as finding,
for every position, its nearest smaller neighbours -- which is what a monotonic stack does in one pass.

\begin{figure}[H]
\centering
\begin{tikzpicture}[diagram, scale=0.92, transform shape]
  % the min-Cartesian tree of A = 3,1,4,2 with parent links highlighted
  \node[troot] (n2) at (0,0)  {$1$};
  \node[tleaf] (n1) at (-2.2,-1.5) {$3$};
  \node[tnode] (n4) at (2.2,-1.5) {$2$};
  \node[tleaf] (n3) at (1.1,-3.0) {$4$};
  \draw[tedgehl] (n2) -- (n1) node[midway, above, font=\scriptsize, text=darkgray] {$a_1$};
  \draw[tedgehl] (n2) -- (n4) node[midway, above, font=\scriptsize, text=darkgray] {$a_4$};
  \draw[tedgehl] (n4) -- (n3) node[midway, right, font=\scriptsize, text=darkgray] {$a_3$};
  \node[annot, align=left, font=\scriptsize] at (0,1.3) {root $=$ minimum of $A$, at position $2$};
  % array strip
  \node[font=\scriptsize] at (-3.0,-3.9) {$A=$};
  \foreach \x/\v in {-2.4/3,-1.6/1,-0.8/4,0.0/2}
    \node[cellnode] at (\x,-3.9) {\v};
  \node[annot, align=left, font=\scriptsize] at (3.4,-3.0)
    {\textbf{inorder} gives $3,1,4,2$\\ \textbf{heap} gives $1\le3,\,1\le2\le4$};
\end{tikzpicture}
\caption{The min-Cartesian tree of $A=(3,1,4,2)$, built by "root is the minimum, recurse on the two
halves". Position subscripts are shown on the edges. The inorder traversal reads the array back
($3,1,4,2$); every parent is $\le$ its children. This is \defref{c26-cartesian}; \thmref{c26-unique}
says it is the \emph{only} such tree, and \lemref{c26-span} says each subtree is that node's
dominance window.}\label{fig:c26-build}
\end{figure}

\subsection{Definitions and Notation}

\begin{definition}{Cartesian tree (min-heap form)}{c26-cartesian}
A \term{Cartesian tree} of an array $a[1..n]$ is a binary tree $T$ on the index set $\{1,\dots,n\}$
satisfying two conditions simultaneously:
\begin{itemize}[nosep]
  \item \textbf{BST order on the index.} The \emph{inorder} traversal of $T$ visits the indices in
        increasing order $1,2,\dots,n$ -- equivalently, the indices in any subtree form a contiguous
        interval, and a node's left subtree holds smaller indices, its right subtree larger ones.
  \item \textbf{Heap order on the value.} For every node $i$ and every child $j$ of $i$,
        $a[i] \le a[j]$. (A \emph{max-}Cartesian tree flips this to $a[i]\ge a[j]$.)
\end{itemize}
Equivalently, defined recursively: if $p$ is a position of the minimum of $a$, then $T$ has root $p$,
left subtree $=$ Cartesian tree of $a[1..p-1]$, right subtree $=$ Cartesian tree of $a[p+1..n]$.
\end{definition}

\noindent\textbf{In plain words.} Put the smallest value on top; everything to its left in the array
becomes the left branch, everything to its right the right branch; repeat inside each side. The array
order is the tree's inorder; the value order is the heap.

\begin{theorem}{Uniqueness under distinct values}{c26-unique}
If the $n$ values of $a$ are pairwise distinct, its min-Cartesian tree is unique. If values repeat,
the tree need not be unique: two equal minima at different positions can be ordered either way, and a
tie-break convention (which equal element becomes the ancestor) must be fixed to restore uniqueness.
\end{theorem}
\begin{proofbox}
By induction on $n$. The heap rule forces the root to be a position $p$ with the minimum value, because
the root must be $\le$ every node in its subtree, which is the whole array. With distinct values that
$p$ is unique. The BST (inorder) rule then forces all indices $<p$ into the left subtree (a contiguous
block) and all indices $>p$ into the right, since inorder lists indices ascending and $p$ sits in the
middle. The left and right blocks are independent arrays, so by the induction hypothesis each has a
unique Cartesian tree, and gluing them under root $p$ is forced. Hence the whole tree is unique. With
two equal minima $p,q$ either may serve as root of that combined subtree, giving two valid trees -- the
tie must be broken externally.
\end{proofbox}

\begin{lemma}{Subtree $=$ dominance window}{c26-span}
Let $L_i$ be the nearest index strictly left of $i$ with $a[L_i] < a[i]$ (or $0$ if none), and $R_i$ the
nearest index strictly right of $i$ with $a[R_i] < a[i]$ (or $n{+}1$ if none). Then the subtree of node
$i$ in the min-Cartesian tree spans exactly the open window $(L_i, R_i)$: the set
$\{L_i+1,\dots,R_i-1\}$. Consequently $a[i]$ is the minimum of that window, and the number of
subarrays for which $a[i]$ is \emph{the} minimum element read at position $i$ is
$(i-L_i)\cdot(R_i-i)$.
\end{lemma}
\begin{proofbox}
An ancestor of $i$ must have value $\le a[i]$ (heap rule) and lie outside $i$'s subtree, and the BST
rule makes the nearest such blockers exactly the closest \emph{smaller} values on each side, $L_i$ and
$R_i$: their positions are the highest ancestors of $i$ (they became $i$'s parent-chain because they
are the first values that could sit above it). No index inside $(L_i,R_i)$ has value $<a[i]$ by
definition of $L_i,R_i$, so $i$ is the minimum of that whole window and its subtree is exactly the
window (an ancestor of $i$ has a smaller value and lies at $L_i$ or $R_i$, outside it). The window
$(L_i,R_i)$ contains $i$; a subarray inside it that includes $i$ starts anywhere in $[L_i+1..i]$ (that
many $= i-L_i$ choices) and ends anywhere in $[i..R_i-1]$ ($R_i-i$ choices), giving the product. \end{proofbox}

\begin{insightbox}
\lemref{c26-span} is the hinge of the entire chapter. It says the geometric fact ``size of $i$'s
subtree'' and the combinatorial fact ``number of ranges where $a[i]$ is the champion'' are the
\emph{same thing}. So any sum of the form $\sum_i a[i]\cdot(\text{ranges where } a[i] \text{ wins})$ --
subarray minimums, the histogram, the Summit -- needs only the two arrays $L,R$, and those fall out of
one monotonic stack pass per side. Read the stack pass right, and you have built the tree; read it as
counts, and you have solved the problem. \secref{c26-monostack} makes the stack pass explicit; \secref{c26-rmq} reads the tree as an index.
\end{insightbox}

\subsection{The algorithm: one stack, one pass}

The recursive definition is $\bigO{n^2}$ in the worst case (a sorted array picks the minimum at one end
every time). The linear build instead walks the array left to right maintaining the \emph{right spine}
of the partially built tree -- the chain from the root down through right children -- in a stack. Values
on that spine are increasing; a new element $a[i]$ pops the spine until it finds its parent, and the
last node it popped becomes its left child.

\begin{lstlisting}[caption={The $O(n)$ monotonic-stack build of a min-Cartesian tree; \code{Lc[], Rc[]} are the left and right children}, label={lst:c26-stackbuild}]
// indices 1..n; a[1..n] the array. Builds children of the min-Cartesian tree.
vector<int> Lc(n+1, 0), Rc(n+1, 0), par(n+1, 0);
vector<int> st;                       // right spine: indices with increasing values
for (int i = 1; i <= n; ++i) {
    int last = 0;                     // last subtree popped -- becomes i's LEFT child
    while (!st.empty() && a[st.back()] > a[i]) {
        last = st.back(); st.pop_back();
    }
    if (!st.empty()) {                // current top beats a[i]: i hangs off it
        Rc[st.back()] = i;            //   as that node's right child
        par[i] = st.back();
    }
    if (last) {                       // the popped chain becomes i's left child
        Lc[i] = last;
        par[last] = i;
    }
    st.push_back(i);
}
int root = st.front();                // the bottom of the stack never got a parent: the root
\end{lstlisting}
\begin{complexitybox}
Time $\bigO{n}$: each index is pushed exactly once and popped at most once across the whole sweep, so
the inner \code{while} performs at most $n$ pops in total -- the classic amortised-stack argument.
Memory $\bigO{n}$ for the three child/parent arrays and the stack. Verdict: the $\bigO{n}$ build, versus
the $\bigO{n^2}$ "scan for the minimum" recursion (a linked-list array is its worst case). The build is
also, in disguise, the histogram stack of \secref{c26-monostack} -- here the stack \emph{is} the tree.
\end{complexitybox}

\subsection{Worked Example}

\begin{example}{Building the tree on $A=(3,1,4,2)$}{c26-cart-run}
Process left to right, tracking the stack $\mathit{st}$ and the children set so far.\\
\emph{$i=1$ ($a=3$).} Stack empty; nothing popped; no parent. Push $1$. $\mathit{st}=[1]$.\\
\emph{$i=2$ ($a=1$).} $a_1=3>1$: pop it, $\mathit{last}=1$. Stack empty, so no right-child link.
Set $\mathit{Lc}[2]=1$ (and $\mathit{par}[1]=2$). Push $2$. $\mathit{st}=[2]$.\\
\emph{$i=3$ ($a=4$).} $a_2=1>4$? no. Nothing popped. $\mathit{Rc}[2]=3$, $\mathit{par}[3]=2$. Push
$3$. $\mathit{st}=[2,3]$.\\
\emph{$i=4$ ($a=2$).} $a_3=4>2$: pop $3$, $\mathit{last}=3$. Now $a_2=1>2$? no. $\mathit{Rc}[2]=4$
(this \emph{overwrites} the earlier link, moving node $3$'s parent to $4$), $\mathit{par}[4]=2$; then
$\mathit{Lc}[4]=3$, $\mathit{par}[3]=4$. Push $4$. $\mathit{st}=[2,4]$.\\
Root $=\mathit{st.front()}=2$. Children: $\mathit{Lc}[2]=1,\ \mathit{Rc}[2]=4,\ \mathit{Lc}[4]=3$ --
exactly \figref{c26-build}: root $1$, left $3$, right $2$ with left child $4$. Inorder
$3,1,4,2$; min-heap throughout. \checkmark
\end{example}
\begin{tracebox}[Span of every node from \lemref{c26-span} on $A$]
Prev/next strictly-smaller $(L_i,R_i)$ with $A=(3,1,4,2)$: node $1$ ($a=3$): $L=0,R=2$; node $2$
($a=1$): $L=0,R=5$ (nothing smaller); node $3$ ($a=4$): $L=4$? no -- nearest smaller left is $a_2=1$
at $2$, nearest smaller right is $a_4=2$ at $4$, so $L=2,R=4$; node $4$ ($a=2$): $L=2,R=5$. Read off
the subtree spans: node $2$'s window $(0,5)$ is all four indices (its subtree is everything); node
$4$'s window $(2,5)$ is $\{3,4\}$ -- matching that $a_4{=}2$ is the minimum of exactly $a[3..4]$.
\end{tracebox}

\subsection{Variations and Twists}

\begin{itemize}[nosep]
  \item \textbf{Max-Cartesian tree}: flip the pop test to $a[\mathit{st.back()}] < a[i]$ and the root
        is the maximum -- this is exactly LeetCode's ``maximum binary tree'' (\secref{c26-treapview}). \hfill\lv{3}
  \item \textbf{Building by the divide route}: recursively find the min (via an RMQ index from
        \secref{c26-rmq}) -- $\bigO{n}$ if you already have a linear build to seed the RMQ, else
        $\bigO{n^2}$. The stack build dominates it. \hfill\lv{4}
  \item \textbf{Parent from one neighbour}: the parent of $i$ is whichever of $L_i, R_i$ has the
        \emph{larger} value (the later of the two blockers to have become a spine node) -- handy for
        reconstructing without the stack. \hfill\lv{4}
  \item \textbf{Ties}: with duplicates, decide whether $L/R$ use $<$ or $\le$ on each side; the same
        asymmetry that prevents double counting in the Summit (\probref{c26-summit}). \hfill\lv{5}
\end{itemize}

\subsection{Common Pitfalls}\label{sec:c26-cartesian-pitfalls}

\begin{warningbox}[Trap 1 -- trusting the recursive minimum scan]
Implementing the definition literally (scan the interval for the minimum, recurse) is $\bigO{n^2}$ on
an already-sorted array -- the minimum always sits at one end, the recursion is a chain of $n$ scans.
For $n=10^5$ that is $\sim 10^{10}$ and the build itself \tle s. Use the stack build
\coderef{c26-stackbuild}; its amortised argument is exactly the "push once, pop at most once" fact.
\end{warningbox}
\begin{lstlisting}[caption={\textcolor{warningred}{WRONG}: the naive divide-by-minimum build is $O(n^2)$ on sorted input}, language=C++, basicstyle=\ttfamily\small]
Node* naive(int lo, int hi) {          // worst case: array sorted -> depth n, n per level
    if (lo > hi) return nullptr;
    int p = argmin(a, lo, hi);         // O(hi-lo+1) SCAN at every node  => O(n^2) total
    Node* r = new Node(p);
    r->l = naive(lo, p-1); r->r = naive(p+1, hi);
    return r;
}
\end{lstlisting}
\begin{lstlisting}[caption={\textcolor{successgreen}{RIGHT}: the monotonic stack builds it once, in $O(n)$}, language=C++, basicstyle=\ttfamily\small]
// the sweep of lst:c26-stackbuild: one push and at most one pop per index.
// No scanning, no recursion, no degenerate-depth blow-up.
for (int i=1;i<=n;i++){ while(!st.empty() && a[st.back()]>a[i]){last=st.back();st.pop_back();}
    if(!st.empty()) Rc[st.back()]=i; if(last) Lc[i]=last; st.push_back(i); }
\end{lstlisting}

\begin{warningbox}[Trap 2 -- forgetting the tie rule]
\thmref{c26-unique} guarantees uniqueness \emph{only for distinct values}. If $a$ repeats, popping on
$>$ versus $\ge$ silently picks a different tree and changes which window each element ``owns''. For a
\emph{max} problem (like \probref{c26-goodsub}) any consistent rule works because you take a maximum
over all positions; for a \emph{sum} problem (LC 907, the Summit) you must break ties asymmetrically --
strict on one side, non-strict on the other -- or a range's minimum gets counted twice.
\end{warningbox}

\begin{warningbox}[Trap 3 -- mixing up which order is which]
The BST order is on the \emph{index} and the heap order is on the \emph{value}. Swapping them gives the
``insertion-order tree'' of \secref{c26-treapview}, a different (also useful) object. The inorder that
reconstructs the array uses the index order -- if your inorder prints something other than
$a[1],\dots,a[n]$ you have a heap-shaped BST, not a Cartesian tree.
\end{warningbox}

\subsection{Problem Lab}

The section's anchor is LeetCode 1793 -- among all subarrays that \emph{contain a fixed index} $k$,
maximise $(\text{minimum of the subarray}) \times (\text{its length})$. The naive double loop is
$\bigO{n^2}$; the span lemma collapses it to two stack passes and one scan, because the only subarrays
worth checking are the full dominance windows of \lemref{c26-span}.

% ---------------------------------------------------------------------------
\solvedproblem{c26-goodsub}{Maximum Score of a Good Subarray}{LeetCode 1793}
  {https://leetcode.com/problems/maximum-score-of-a-good-subarray/}{LC Hard · L5}
  {Cartesian span lemma, prev/next smaller via monotonic stack, maximise min times width}

\begin{problemstatement}[Maximum Score of a Good Subarray]
You are given a $0$-indexed integer array $\mathit{nums}$ and an integer $k$. A \textbf{good} subarray
$(i, j)$ is one with $i \le k \le j$. Its \textbf{score} is
$\min(\mathit{nums}[i], \dots, \mathit{nums}[j]) \cdot (j - i + 1)$. Return the \textbf{maximum} score
among all good subarrays. $1 \le \mathit{nums.\text{length}} \le 10^{5}$, $1 \le \mathit{nums}[i] \le 10^{5}$,
and $0 \le k < \mathit{nums.\text{length}}$.
\end{problemstatement}

\begin{attemptgate}[Attempt it yourself first -- \lv{5}, give yourself 30 minutes]
The double loop over all $(i,j)$ through $k$ is $\bigO{n^2}$. The winning move is to stop thinking about
subarrays and think about \emph{which element is the minimum of the chosen subarray}.
\peek{Fix a position $m$. What is the \emph{longest} subarray through $k$ on which $\mathit{nums}[m]$
is the minimum? (Hint: the dominance window of \lemref{c26-span}.)}
\nudge{For $\mathit{nums}[m]$ to be the minimum, the subarray cannot cross a strictly-smaller element,
so it lies inside $(L_m, R_m)$. If that window contains $k$, the best subarray with minimum at $m$ is the
whole window, scoring $\mathit{nums}[m]\cdot(R_m - L_m - 1)$. Compute every $L_m, R_m$ with two stack passes.}
\fullhint{Take $\max$ over $m$ whose window contains $k$ of $\mathit{nums}[m]\cdot(R_m-L_m-1)$, where
$L_m$ ($R_m$) is the nearest index to the left (right) with a strictly smaller value. Any good subarray
has its minimum at some $m$, and its full window through $m$ is at least as good and still contains $k$.
The answer for the sample is $3 \times 5 = 15$.}
\end{attemptgate}

\solstep{Understanding the Problem}
We maximise $(\min \text{ over a window}) \times (\text{window length})$ subject to the window
containing index $k$. The two factors pull opposite ways -- a longer window is smaller in minimum -- so
naively you must try every window through $k$, $\bigO{n^2}$. The reframing: the value $\min$ is attained
at \emph{some} position $m$ in the window, so enumerate $m$ instead of the window. For a fixed $m$,
the minimum over any window containing $m$ is at most $\mathit{nums}[m]$, and equals it as long as the
window stays inside $(L_m, R_m)$ -- the dominance window of \lemref{c26-span}, found by nearest
strictly-smaller neighbours. So the \emph{best window that takes $\mathit{nums}[m]$ as its minimum} is
the whole $(L_m, R_m)$, \emph{if} that window reaches across $k$.

\solstep{Key Observations}
\begin{observation}
For a good subarray with score $s$, let $m$ be a position attaining its minimum. The dominance window
$(L_m, R_m)$ contains that subarray (nothing inside it is strictly smaller than $\mathit{nums}[m]$),
hence contains $k$ too, and using the whole window keeps the minimum $\mathit{nums}[m]$ while never
shrinking the length. So it is never worse to expand to the full window.
\end{observation}
\begin{observation}
Therefore the optimum equals $\max$ over $m$ with $L_m < k < R_m$ of
$\mathit{nums}[m]\cdot(R_m - L_m - 1)$. Only windows that \emph{straddle} $k$ are candidates.
\end{observation}
\begin{observation}
Both $L_m$ and $R_m$ (nearest strictly-smaller on each side) come from monotonic-stack passes -- one
left-to-right, one right-to-left -- in $\bigO{n}$. That is the whole algorithm.
\end{observation}

\solstep{Think It Through With Me}
\begin{thinkbox}
\deadend{Enumerate $i \le k$ and $j \ge k$ and compute the running minimum.}
{$\bigO{n^2}$: for each start you grow the end and update the min -- on $n=10^5$ that is $\sim 10^{10}$,
a \tle. It also hides the structure; every extra $j$ beyond the dominance window contributes nothing new.}
\deadend{Binary-search the answer on the minimum value.}
{Works ($\bigO{n\log V}$: is there a length-$k$-containing window with all values $\ge$ a threshold?),
but it is a heavier hammer than needed and obscures \emph{why} the answer is a per-position span.}
\aha{Assign each good subarray's score to the \emph{position of its minimum}. That position's dominance
window is the best subarray for that minimum; keep only windows that cover $k$. Two stack passes give
every window, one scan gives the max -- $\bigO{n}$ and a direct hit on \lemref{c26-span}.}
\end{thinkbox}

\solstep{Solution}
Compute $L[m]$ = nearest index $< m$ with $\mathit{nums}[L]<\mathit{nums}[m]$ (else $-1$) with an
increasing stack scanning left-to-right; compute $R[m]$ = nearest index $> m$ with
$\mathit{nums}[R]<\mathit{nums}[m]$ (else $n$) by scanning right-to-left the same way. Then the window
in which $\mathit{nums}[m]$ is the minimum is $(L[m], R[m])$, length $R[m]-L[m]-1$. It is a good
subarray iff $L[m] < k < R[m]$. Answer is the maximum of $\mathit{nums}[m]\cdot(R[m]-L[m]-1)$ over such
$m$. Because we only ever take a maximum over all positions, a consistent choice of $<$ on both sides is
safe -- duplicates cannot cause an under-count here.

\solstep{Proof of Correctness}
\begin{proofbox}[The span max equals the true optimum]
(\emph{Achievable}) For each $m$ with $L[m]<k<R[m]$, the subarray $[L[m]+1,\, R[m]-1]$ contains $k$
(hence is good) and has minimum exactly $\mathit{nums}[m]$ (no strictly smaller value lies inside the
open window, and $\mathit{nums}[m]$ itself is in it). So its score $\mathit{nums}[m]\cdot(R[m]-L[m]-1)$
is attainable and the reported maximum is a valid lower bound on the optimum. (\emph{Optimal}) Let
$[i,j]$ be any good subarray of maximal score and $m$ a position of its minimum. Since all values in
$[i,j]$ are $\ge \mathit{nums}[m]$, no index of $[i,j]$ has a value strictly below $\mathit{nums}[m]$,
so $[i,j]\subseteq (L[m], R[m])$; as $[i,j]\ni k$ we get $L[m]<k<R[m]$. The window
$[L[m]+1,R[m]-1]$ contains $[i,j]$, so it is at least as long with the same minimum, giving score
$\ge$ that of $[i,j]$. Thus some candidate considered by the algorithm is $\ge$ the optimum's score --
combined with achievability, the maximum the algorithm returns \emph{is} the optimum. \end{proofbox}

\solstep{Implementation}
\begin{lstlisting}[caption={Maximum Score of a Good Subarray -- spans from two stack passes (LeetCode 1793)}, label={lst:c26-goodsub-soln}]
class Solution {
public:
    int maximumScore(vector<int>& nums, int k) {
        int n = (int)nums.size();
        vector<int> L(n, -1), R(n, n), st;
        for (int i = 0; i < n; ++i) {                 // prev strictly smaller
            while (!st.empty() && nums[st.back()] >= nums[i]) st.pop_back();
            if (!st.empty()) L[i] = st.back();
            st.push_back(i);
        }
        st.clear();
        for (int i = n - 1; i >= 0; --i) {            // next strictly smaller
            while (!st.empty() && nums[st.back()] >= nums[i]) st.pop_back();
            if (!st.empty()) R[i] = st.back();
            st.push_back(i);
        }
        long long ans = 0;
        for (int m = 0; m < n; ++m)
            if (L[m] < k && k < R[m])                 // its window straddles k
                ans = max(ans, 1LL * nums[m] * (R[m] - L[m] - 1));
        return (int)ans;
    }
};
\end{lstlisting}

\solstep{Trace}
\begin{tracebox}[Official sample: \texttt{nums = [1,4,3,7,4,5]}, \texttt{k = 3}]
Prev-smaller $L$: $[-1,0,0,2,2,4]$. Next-smaller $R$: $[6,2,6,4,6,6]$. Keep $m$ with $L[m]<3<R[m]$:
$m=0$: $1\cdot(6-(-1)-1=6)=6$. $m=2$: $L=0<3<R=6$, $3\cdot5=15$. $m=3$: $L=2<3<R=4$, $7\cdot1=7$.
$m=4$: $L=2<3<R=6$, $4\cdot3=12$. $m=1,5$ fail ($R[1]=2\not>3$, $L[5]=4\not<3$). Maximum $=15$ --
the subarray $[1..5]=[4,3,7,4,5]$, minimum $3$, length $5$, score $15$. \checkmark matches the official
output.
\end{tracebox}

\solstep{Complexity}
\begin{complexitybox}
Time $\bigO{n}$: two monotonic-stack passes (each index pushed once, popped at most once) plus one
$\bigO{n}$ scan for the maximum. Memory $\bigO{n}$ for $L,R$ and the stack. Against the $\bigO{n^2}$
window enumeration this is the full win the span lemma promises -- and it is literally the Cartesian
build \coderef{c26-stackbuild} run once per side. Verdict: \ac.
\end{complexitybox}

\solstep{Pitfalls and Variants}
\begin{warningbox}
Three ways to fall: (i) treating a position whose window does \emph{not} contain $k$ as a candidate --
it is not a good subarray, skip it (the $L[m]<k<R[m]$ guard); (ii) using $\le$ instead of $<$ so a
duplicate minimum truncates the window early -- harmless for a pure \emph{max}, but know that you chose
it; (iii) $32$-bit overflow in $\mathit{nums}[m]\cdot\text{length}$ ($10^5 \cdot 10^5 = 10^{10}$) --
multiply in \code{long long}. \emph{Variant:} the two-pointer "grow toward the taller neighbour" greedy
also passes, but the span proof is the one that generalises to the Summit's signed sum.
\end{warningbox}

\begin{insightbox}[Key takeaways]
\begin{enumerate}[nosep, leftmargin=1.6em]
  \item \textbf{A Cartesian tree is a BST on index and a heap on value at once} (\defref{c26-cartesian}),
        unique on distinct values, root $=$ the minimum (\thmref{c26-unique}).
  \item \textbf{The subtree is the dominance window} (\lemref{c26-span}) -- the pivot from a tree-shape
        fact to a counting fact, and the reason one stack pass solves so much.
  \item \textbf{Build it, or just read it.} \coderef{c26-stackbuild} materialises the tree; for scoring
        problems you keep only $L,R$ and never link a single node (\probref{c26-goodsub}).
\end{enumerate}
\end{insightbox}

\begin{hardeningbox}
\begin{itemize}[nosep, leftmargin=1.4em]
  \item \textbf{Best subarray through each index} -- run the same span max for every $k$; the $L,R$
        arrays answer all $k$ after one build.
  \item \textbf{Largest all-positive rectangle under a skyline} -- the histogram of
        \secref{c26-monostack} is \emph{exactly} this span maximum with width forced to start $=1$; the
        link is made explicit there.
  \item \textbf{Count the ranges, don't max them} -- replace $\max$ by a sum and you are on the road to
        LC 907 and the Summit; the only new care is asymmetric tie-breaking.
\end{itemize}
\end{hardeningbox}

\subsection{Practice Problems}

\begin{stuckbox}
\begin{itemize}[nosep, leftmargin=1.4em]
  \item Did you assign each subarray's value to the \emph{position of its minimum}, not to the window?
        (\probref{c26-goodsub}.)
  \item Did you guard that the window must \emph{contain} $k$ before scoring it?
  \item For the sum version: did you make one side strict and the other non-strict so a range with two
        equal minima is counted once, not twice?
\end{itemize}
\end{stuckbox}

\begin{practiceset}
  \practice{Sum of Subarray Minimums}{LeetCode 907}
    {https://leetcode.com/problems/sum-of-subarray-minimums/}{LC Medium · L3}
    {Return $\sum$ over all subarrays of their minimum, modulo $10^9+7$; $\mathit{arr}.\text{length}\le 4\cdot10^4$.
    This is \lemref{c26-span} turned from a maximum into a sum: the answer is
    $\sum_i a[i]\cdot(i-L_i)\cdot(R_i-i)$, the count of subarrays whose minimum is $a[i]$ read at $i$,
    from one stack pass for each of $L_i, R_i$. The one trap \emph{not} present in \probref{c26-goodsub}:
    a subarray with two equal minima must be credited to exactly one position, so break ties by taking
    \emph{strict} $<$ on one side (say prev smaller) and \emph{non-strict} $\le$ on the other (next
    smaller-or-equal). Keep every partial product in \code{long long} and reduce modulo $10^9+7$. The
    sample $[3,1,2,4]$ gives $17$ -- and the same two-pass-with-asymmetric-ties template is exactly what
    the Summit (\probref{c26-summit}) will run for both the max and the min side at once.}
\end{practiceset}

\begin{recallbox}[10-second refresher -- 9.1 The Cartesian Tree in $O(n)$]
A Cartesian tree is \emph{inorder $=$ index} and \emph{heap $=$ value}: root is the minimum, unique on
distinct values (\defref{c26-cartesian}, \thmref{c26-unique}). One monotonic stack builds it in
$\bigO{n}$ by maintaining the increasing right spine (\coderef{c26-stackbuild}). Its key payoff,
\lemref{c26-span}, says node $i$'s subtree is the maximal window where $a[i]$ is the minimum -- so
"each element's dominance range" problems need only the prev/next smaller arrays. LeetCode 1793
(\probref{c26-goodsub}) maximises $\min\times\text{length}$ through $k$ by scoring exactly those windows;
LC 907 sums them once you break ties asymmetrically.
\selftest{
  \item Why must the root of a min-Cartesian tree be the position of the minimum?
  \item State \lemref{c26-span} in one sentence and read $a_4$'s window off $A=(3,1,4,2)$.
  \item Why is the naive divide-by-minimum build $\bigO{n^2}$, and what input forces it?
  \item In \probref{c26-goodsub}, why is it safe to expand a good subarray to $m$'s full window?
  \item What changes between the max version (\probref{c26-goodsub}) and the sum version (LC 907) about ties?
}
\end{recallbox}

% !TEX root = ../../tree_notes.tex
% =====================================================================
%  9.2  Monotonic Stack Applications            sec:c26-monostack
%  Exports: def:c26-monostack  thm:c26-hist  fig:c26-histfig
%           lst:c26-hist  lst:c26-deque  ex:c26-hist-run
%           prob:c26-hist84  lst:c26-hist84-soln  (practices: LC 85, 42, 239, 739)
% =====================================================================
\section{Monotonic Stack Applications}\label{sec:c26-monostack}

\begin{previouslybox}
\begin{enumerate}[nosep, leftmargin=1.6em]
  \item What two orders define a Cartesian tree, and where is its root? (\defref{c26-cartesian}; the
        minimum, by the heap rule)
  \item State \lemref{c26-span} in one sentence: node $i$'s subtree is the window over which $a[i]$
        is the \emph{what}, bounded by the nearest \emph{what} on each side?
  \item The build \coderef{c26-stackbuild} keeps the stack's values \emph{increasing}. When does index
        $i$ get pushed, and which node does it pop past? (past the run of larger values)
  \item Why is a whole monotonic sweep $\bigO{n}$ even with a nested \code{while}-pop loop? (each index
        pushed once, popped at most once -- the amortised-stack argument)
\end{enumerate}
\end{previouslybox}

\begin{topicintro}
Strip the Cartesian tree back to its mechanics and you are left with one idea so small it fits in a
\code{while} loop: keep a stack whose values are \emph{monotone}, and whenever a new element breaks the
monotonicity, pop the run it breaks -- and \emph{at the instant of popping, the popped element's answer
is finally known}. That is the \term{monotonic stack}, and it is the workhorse of this whole chapter.
The popped element has just met its nearest greater (or smaller) neighbour: the exact bound
\lemref{c26-span} needed.

This section makes the stack a technique in its own right -- you no longer need to name the tree -- and
runs the entire application family on it: the largest rectangle in a histogram (the solved anchor,
LeetCode 84), the next-greater family (daily temperatures), sliding-window extremes (a monotonic
\emph{deque}, LeetCode 239), and trapping rain water. The unifying sentence, which you should memorise:
\textbf{a bar's rectangle is decided the moment it is popped, because that is when the shortest bar on
either side of its dominance window are both known}.
\textbf{Prerequisites.} \secref{c26-cartesian} (the span lemma is the reason the stack works) and a
\code{stack}/\code{deque} container.

\textbf{After this section you can}
\begin{itemize}[nosep]
  \item decide the pop-comparison ($>$ vs $\ge$, $<$ vs $\le$) that matches the bounds a problem needs; \hfill\lv{3}
  \item solve largest-rectangle-in-histogram cold, and explain why the width is read at pop time; \hfill\lv{4}--\lv{5}
  \item convert any "next greater / smaller element" problem to one decreasing or increasing stack pass; \hfill\lv{3}
  \item keep a sliding-window extreme with a monotonic deque at amortised $\bigO{1}$ per position; \hfill\lv{4}
\end{itemize}
\end{topicintro}

\begin{recognitionbox}[Pattern card -- nearest bigger/smaller neighbour]
\reckey{each element's fate is sealed by the \emph{first} element to its left or right that is
\emph{larger} (or \emph{smaller}); or a \emph{window extreme} that must be reported as the window slides.}
\recuse{
  \item "largest rectangle under a skyline", "maximal rectangle of 1s in a grid" -- pop when the bar
        stops being the limiting height (\probref{c26-hist84}, LC 85);
  \item "next greater element", "days until a warmer temperature" -- one decreasing stack;
  \item "max (min) in every window of size $k$" -- monotonic \emph{deque} (\coderef{c26-deque});
  \item "water trapped between bars" -- stack of walls, or two-pointer prefix maxima (LC 42).
}
\recnotuse{
  \item you need arbitrary-range extremes over a \emph{dynamic} array -- segment tree (\chref{c13}),
        the stack cannot forget its pops;
  \item the window extreme is the \emph{median} -- a monotonic deque tracks only the extreme, not the
        middle; that is a multiset, Chapter~8 (\chref{c06});
  \item "the $k$-th smallest overall" -- a heap (\chref{c07}), not a stack.
}
One-line test: \emph{does the answer for an element depend only on the nearest element that beats it?
Then a monotonic stack finds every such answer in one pass.}
\end{recognitionbox}

\subsection{Intuition}

Imagine the histogram $H=(2,1,5,6,2,3)$ of \figref{c26-hist} as a city skyline and drop a vertical
waterline at some height $h$. At height $h$, the tallest rectangle of height $h$ is limited by the
\emph{nearest bar shorter than $h$} on the left and the nearest bar shorter than $h$ on the right --
everything between them is at least $h$. So for the bar at position $b$, the rectangle of height
$H[b]$ that uses bar $b$ as its \emph{limiting} height spans exactly those two shorters' inner window:
this is \lemref{c26-span} again, now read as ``how far can a bar of height $H[b]$ extend?''

The monotonic stack computes all these windows simultaneously. Walk the bars left to right keeping a
stack of positions whose heights are \emph{strictly increasing} -- a rising staircase. When a new bar
$i$ is \emph{lower} than the stack top, the top bar has just met its right-boundary shorter neighbour:
pop it, and its left-boundary shorter neighbour is now the (new) stack top. Between those two bounds,
its rectangle is fully determined, so evaluate its area \emph{right now}. Bars that never get popped
until the end are popped by a zero-height sentinel appended to the array. One pass, every rectangle --
this is the same loop as \coderef{c26-stackbuild}, with the "compute the span's area" step surfacing.

\begin{figure}[H]
\centering
\begin{tikzpicture}[diagram, x=11mm, y=6mm]
  \draw[thick, darkgray] (0,0) -- (6.4,0);
  \foreach \i/\h in {0/2,1/1,2/5,3/6,4/2,5/3} {
    \fill[primaryblue!22, draw=primaryblue, thick]
      (\i,0) rectangle (\i+1,\h);
    \node[font=\scriptsize, text=darkgray] at (\i+0.5,-0.45) {\i};
    \node[font=\scriptsize\bfseries, text=deepblue] at (\i+0.5,\h+0.3) {\h};}
  % the limiting window for the height-5 bar: bars 2 and 3, blocked by bar1(=1) and bar4(=2)
  \draw[ultra thick, warningred] (2,0) rectangle (4,5);
  \draw[warningred, dashed] (1,5) -- (4,5);
  \node[font=\scriptsize\bfseries, text=warningred, align=center] at (3,2.5) {$h{=}5$\\ width $=2$};
  \node[annot, align=left, font=\scriptsize] at (7.0,2.5)
    {bar $2$ pops at bar $4$:\\ left bound $=$ bar $1$ ($h{=}1$)\\ right bound $=$ bar $4$ ($h{=}2$)\\
     width $4-1-1=2$, area $10$};
\end{tikzpicture}
\caption{Why the width is decided \emph{at pop time}. Bar $2$ (height $5$) sits on the increasing stack
until bar $4$ (height $2$) is lower and pops it. At that instant its left-boundary shorter bar is the
new stack top, bar $1$ (height $1$), and its right-boundary shorter bar is $i=4$; the rectangle of
height $5$ spans bars $2$--$3$, width $4-1-1=2$, area $10$. That is \thmref{c26-hist} and the answer for
$H$.}\label{fig:c26-histfig}
\end{figure}

\subsection{Definitions and Notation}

\begin{definition}{Monotonic stack}{c26-monostack}
A \term{monotonic stack} is a stack of indices whose stored values are kept \emph{non-decreasing} (an
\emph{increasing} stack, used to find nearest \emph{smaller} elements) or \emph{non-increasing} (a
\emph{decreasing} stack, for nearest \emph{greater} elements). When a new index $i$ arrives, pop from
the top while the invariant would be violated; each popped index has just discovered its nearest
beating neighbour ($i$ to its right), and the new top is its neighbour to the left. Every index enters
and leaves the stack once, so the total work is $\bigO{n}$.
\end{definition}

\noindent\textbf{In plain words.} The stack is a queue of elements still ``waiting'' to meet someone
bigger/smaller than them. The moment such someone appears, all the elements it beats are answered at
once and leave. Nothing is ever looked at twice.

\begin{theorem}{A popped bar's rectangle is finalised at pop}{c26-hist}
Run an increasing stack (pop while $H[\mathit{top}] \ge$ the incoming height) over $H$ with a
zero-height sentinel appended. When bar $b$ is popped at index $i$, let $L$ be the new stack top
($L=-1$ if the stack became empty). Then $i$ is the nearest index to the right of $b$ with a height
\emph{strictly less} than $H[b]$ reached, $L$ is the nearest to the left, and the largest rectangle of
height $H[b]$ has width $i - L - 1$ and area $H[b]\cdot(i-L-1)$. Taking this over every pop yields the
maximum rectangle in the histogram.
\end{theorem}
\begin{proofbox}
Bars are pushed in increasing height order, so while $b$ stayed on the stack every bar pushed after it
and popped before $i$ had height $\ge H[b]$ and lies inside $b$'s window. The pop triggers at $i$ exactly
when $H[i] < H[b]$ (or the sentinel), so $i$ is the first shorter bar to the right; after popping $b$,
the stack top $L$ is the last bar with $H[L] < H[b]$ to the left (everything between $L$ and $b$ has
height $\ge H[b]$ or was already popped as taller). Hence the maximal contiguous run of bars of height
$\ge H[b]$ that includes $b$ is $[L+1, i-1]$, of width $i-L-1$; the rectangle of height $H[b]$ is
attained over exactly that run. \emph{Every} candidate optimal rectangle has some height equal to one of
the bars (raise it until it hits a bar top), say bar $b$, and is at most the width just computed, so the
best over all pops equals the true optimum. \end{proofbox}

\begin{insightbox}
\thmref{c26-hist} is \lemref{c26-span} wearing a hat. ``The subtree of $b$ is its dominance window'' and
``the rectangle of bar $b$ is bounded by the nearest shorter bars'' are the \emph{same statement} -- a
shorter bar is exactly a smaller value that becomes an ancestor. If you internalised 9.1, the histogram
is already solved; this section just shows the stack \emph{without} drawing any tree, which is how you
will actually write it in contest code.
\end{insightbox}

\subsection{The algorithm: one pass over the bars}

The code is the sentinel-augmented increasing stack. Keep positions, not heights, so you can recover
the width from the index gap; pop on $\ge$ so a taller-or-equal bar is forced to settle its rectangle
before an equal bar can widen it.

\begin{lstlisting}[caption={Largest rectangle in a histogram -- one increasing stack with a zero sentinel}, label={lst:c26-hist}]
long long largestRectangle(const vector<long long>& H) {
    int n = (int)H.size();
    vector<int> st;                       // indices, heights strictly rising
    long long best = 0;
    for (int i = 0; i <= n; ++i) {        // i == n is the 0-height sentinel
        long long cur = (i == n ? 0 : H[i]);
        while (!st.empty() && H[st.back()] >= cur) {
            long long h = H[st.back()]; st.pop_back();
            int L = st.empty() ? -1 : st.back();   // new top = left shorter bound
            best = max(best, h * (i - L - 1));      // width decided at the pop
        }
        st.push_back(i);
    }
    return best;
}
\end{lstlisting}
\begin{complexitybox}
Time $\bigO{n}$ -- each index pushed once and popped at most once. Memory $\bigO{n}$ for the stack.
Verdict: the linear skyline algorithm; the naive $\bigO{n^2}$ (grow a window from each bar, track its
running minimum) is the same answer an order slower and \tle s for $n\sim10^5$.
\end{complexitybox}

\subsection{Worked Example}

\begin{example}{The stack on $H=(2,1,5,6,2,3)$}{c26-hist-run}
Read the pop events; $\mathit{st}$ holds indices and the sentinel is at $i=6$ with height $0$.\\
$i{=}0$ push; $i{=}1$ (height $1$) pops bar $0$ ($h=2$, $L=-1$, width $2$? no -- $1-(-1)-1=1$, area $2$),
then pushes $1$. $i{=}2,3$ push (heights $5,6$ rising). At $i{=}4$ (height $2$): pop bar $3$ ($h=6$,
$L=2$, width $4-2-1=1$, area $6$); pop bar $2$ ($h=5$, $L=1$, width $4-1-1=2$, area $10$); push $4$.
$i{=}5$ (height $3$) pushes. Sentinel $i{=}6$: pop bar $5$ ($h=3$, $L=4$, width $1$, area $3$); pop bar
$4$ ($h=2$, $L=1$, width $6-1-1=4$, area $8$); pop bar $1$ ($h=1$, $L=-1$, width $6$, area $6$).
Maximum over the areas $\{2,6,10,3,8,6\}$ is $\boxed{10}$ -- the height-$5$ rectangle over bars $2,3$. \checkmark
\end{example}
\begin{tracebox}[Pop ledger -- each bar settles when a shorter one arrives]
bar 0 (h2): popped by bar1, width 1, area 2.\\
bar 1 (h1): popped by sentinel, width 6, area 6 (spans everything -- height 1 reaches all).\\
bar 2 (h5): popped by bar4, width 2, area 10.\\
bar 3 (h6): popped by bar4, width 1, area 6.\\
bar 4 (h2): popped by sentinel, width 4, area 8 (bars 2..5 minus the height-1 blocker at bar1).\\
bar 5 (h3): popped by sentinel, width 1, area 3.  Best area $=10$.
\end{tracebox}

\subsection{The wider family}

The exact same loop, with the comparison flipped and the ``what do you compute at pop'' changed, mops up
a shelf of problems. Keep the four dialects straight:

\begin{itemize}[nosep]
  \item \textbf{Next greater element} (\emph{decreasing} stack). Pop while the top's value is $<$ the
        incoming one; the incoming index is that element's next greater. Daily temperatures
        (LC 739) is literally ``distance to next greater''. \hfill\lv{3}
  \item \textbf{Previous smaller} (\emph{increasing} stack). The pop you already do in the histogram
        gives the nearest shorter bar to the left for free -- this is $L$ in \thmref{c26-hist}. \hfill\lv{3}
  \item \textbf{Sliding-window maximum} (\emph{monotonic deque}, LC 239). Same decreasing stack, but you
        also drop the \emph{front} once it leaves the window of size $k$; a deque pops both ends. \hfill\lv{4}
  \item \textbf{Trapping rain water} (LC 42). A decreasing stack of walls catches valleys: pop a floor,
        the current bar and the new top are its bounding walls, fill the basin. (Or two-pointer prefix
        maxima -- same idea, see Practice.) \hfill\lv{4}--\lv{5}
\end{itemize}

\begin{lstlisting}[caption={Sliding-window maximum with a monotonic deque (LeetCode 239) -- drop smaller from the back, out-of-window from the front}, label={lst:c26-deque}]
vector<int> maxSlidingWindow(vector<int>& nums, int k) {
    deque<int> dq;                       // indices, values non-increasing, front = current max
    vector<int> out;
    for (int i = 0; i < (int)nums.size(); ++i) {
        while (!dq.empty() && nums[dq.back()] <= nums[i]) dq.pop_back();  // dead: i beats them
        dq.push_back(i);
        if (dq.front() <= i - k) dq.pop_front();                          // left the window
        if (i >= k - 1) out.push_back(nums[dq.front()]);                  // window is full
    }
    return out;
}
\end{lstlisting}

\subsection{Common Pitfalls}\label{sec:c26-monostack-pitfalls}

\begin{warningbox}[Trap 1 -- $>$ versus $\ge$ silently miscounts widths]
Popping while $H[\text{top}] > \text{cur}$ (strict) leaves \emph{equal} heights stacked, and the left
bound $L$ for a popped bar is then an equal bar -- so the width is under-counted until a later equal bar
pays it back, and the \emph{maximum} still comes out right. But the moment you \emph{sum} areas or count
ranges (LC 907, the Summit) the same off-by-relationship double- or under-counts. Pick one convention per
side deliberately -- here $\ge$ on the pop keeps the max correct and each rectangle's true width is
recovered when its \emph{last} equal-height bar pops.
\end{warningbox}
\begin{lstlisting}[caption={\textcolor{warningred}{WRONG}: no sentinel -- tall bars never get their right bound}, language=C++, basicstyle=\ttfamily\small]
for (int i = 0; i < n; ++i) { ... }      // stack may still hold bars when the loop ends
// the tallest bars (right of every shorter bar) are never popped -> their rectangles uncomputed
\end{lstlisting}
\begin{lstlisting}[caption={\textcolor{successgreen}{RIGHT}: run to i==n with cur==0 to flush the stack}, language=C++, basicstyle=\ttfamily\small]
for (int i = 0; i <= n; ++i) {           // i == n appends a height-0 sentinel
    long long cur = (i == n ? 0 : H[i]); // forces every remaining bar to settle its width
    ...
}
\end{lstlisting}

\begin{warningbox}[Trap 2 -- storing values instead of indices]
The width is an \emph{index} difference ($i - L - 1$), so the stack must hold positions, not heights. If
you push heights you lose the ability to measure how far the rectangle reaches -- the whole point of
\thmref{c26-hist}. Keep indices and read $H[\text{index}]$ when you need the height.
\end{warningbox}

\begin{warningbox}[Trap 3 -- a deque for the wrong window statistic]
A monotonic deque answers sliding-window \emph{max/min} in $\bigO{1}$ amortised because you may discard
any element that can never again be the extreme. You cannot do that for the sliding \emph{median} or
\emph{sum} -- smaller elements still matter for the median -- so those need an ordered multiset or two
heaps (Chapter~8, \chref{c06}; Chapter~10, \chref{c07}). Do not reach for the deque out of habit.
\end{warningbox}

\subsection{Problem Lab}

The section's anchor is the canonical LeetCode 84 -- the histogram running example. It is chosen because
it is \thmref{c26-hist} with the proof already done: one increasing stack, a sentinel, and a multiply.

% ---------------------------------------------------------------------------
\solvedproblem{c26-hist84}{Largest Rectangle in Histogram}{LeetCode 84}
  {https://leetcode.com/problems/largest-rectangle-in-histogram/}{LC Hard · L5}
  {monotonic (increasing) stack, sentinel flush, width read at pop time}

\begin{problemstatement}[Largest Rectangle in Histogram]
Given an array of $n$ non-negative integers $\mathit{heights}$ representing a histogram where the width
of each bar is $1$, return the area of the \textbf{largest rectangle} that can be formed within the
histogram. $1 \le n \le 10^{5}$ and $0 \le \mathit{heights}[i] \le 10^{4}$.
\end{problemstatement}

\begin{attemptgate}[Attempt it yourself first -- \lv{5}, give yourself 30 minutes]
The answer is a rectangle whose height equals \emph{some} bar and whose width is how far that bar's
height can extend left and right before a shorter bar blocks it.
\peek{Fix a bar $b$. What bounds the widest rectangle of height $H[b]$? (Nearest bar shorter than $H[b]$
on each side -- \thmref{c26-hist}.)}
\nudge{Sweep left to right with an increasing stack of indices. When a shorter bar arrives it pops the
taller ones, and for each popped bar you now know both bounds -- compute its rectangle immediately.
Append a $0$-height bar at the end to flush whatever is left.}
\fullhint{On pop of bar $x$ at incoming index $i$: $L$ is the new stack top ($-1$ if empty), width $=
i - L - 1$, area $= H[x]\cdot\text{width}$. Take the max over all pops. $\bigO{n}$. For the sample the
height-$5$ bar pops giving width $2$, area $10$.}
\end{attemptgate}

\solstep{Understanding the Problem}
A rectangle in the histogram is a contiguous set of bars lifted to a common height; that height cannot
exceed the shortest bar in the set. So an optimal rectangle has height exactly $H[b]$ for some bar $b$
(its limiting bar) and width equal to the maximal run around $b$ of bars all $\ge H[b]$. The task is thus:
for every bar, find its left and right \emph{nearest strictly shorter} bars, and evaluate
$H[b]\cdot\text{width}$. Doing that naively per bar is $\bigO{n}$ each, $\bigO{n^2}$ overall; the
monotonic stack finds every bar's two bounds in one shared pass.

\solstep{Key Observations}
\begin{observation}
On an increasing stack of indices, when bar $x$ is popped by an incoming bar $i$, $i$ is $x$'s nearest
shorter bar on the \emph{right} and the new stack top $L$ is $x$'s nearest shorter bar on the
\emph{left}. Both bounds become known at the single moment of the pop (\thmref{c26-hist}).
\end{observation}
\begin{observation}
The rectangle of height $H[x]$ therefore spans indices $[L+1, i-1]$, width $i - L - 1$, area
$H[x]\cdot(i-L-1)$. No bar between $L$ and $i$ is shorter than $H[x]$ (they were pushed after $x$ or popped as taller).
\end{observation}
\begin{observation}
Appending a $0$-height sentinel guarantees every bar is eventually popped, so no candidate rectangle is
missed even for a strictly-increasing skyline.
\end{observation}

\solstep{Think It Through With Me}
\begin{thinkbox}
\deadend{For each bar, walk left and right counting bars of height $\ge$ it.}
{$\bigO{n^2}$: on a flat histogram of height $1$ and $n=10^5$, every bar scans all $n$, giving $10^{10}$
-- \tle. It also recomputes the same bounds over and over.}
\deadend{Divide and conquer on the minimum (Cartesian tree + RMQ).}
{Correct and $\bigO{n\log n}$ (build the tree, recurse on min), and it \emph{is} the same structure --
but the constant and the recursion are unnecessary when one flat stack pass is $\bigO{n}$.}
\aha{One increasing stack with a sentinel: a bar is popped exactly when its right bound appears, its left
bound is the new top, and the area is a multiply. $\bigO{n}$, in-place, no tree built.}
\end{thinkbox}

\solstep{Solution}
Traverse $i$ from $0$ to $n$ inclusive. Let $\mathit{cur}$ be $H[i]$ (or $0$ at the $i=n$ sentinel). While
the stack is non-empty and $H[\text{top}] \ge \mathit{cur}$: pop index $x$; set $L = \text{new top}$ or
$-1$ if empty; $\mathit{best} = \max(\mathit{best}, H[x]\cdot(i - L - 1))$. Then push $i$. Return
$\mathit{best}$ (as a \code{long long} if the product can exceed $2^{31}$).

\solstep{Proof of Correctness}
\begin{proofbox}[The stack attains the true maximum]
(\emph{Each pop evaluates a valid rectangle}) By \thmref{c26-hist}, at the pop of $x$ the bounds $L$ and
$i$ are the nearest strictly-shorter bars on each side, so $[L+1, i-1]$ is the widest run of bars of
height $\ge H[x]$ containing $x$; the area $H[x]\cdot(i-L-1)$ is a genuinely realisable rectangle and the
running maximum never exceeds the true optimum. (\emph{Nothing is missed}) Any optimal rectangle has
height equal to the shortest bar $x$ it covers; by the sentinel argument $x$ is popped at some index
which is precisely its right shorter bound, and at that pop its left bound $L$ is exactly its left
shorter bound, so the algorithm evaluates the optimum's own area. Hence $\mathit{best}$ equals the true
maximum. \end{proofbox}

\solstep{Implementation}
\begin{lstlisting}[caption={Largest Rectangle in Histogram -- the sentinel monotonic stack (LeetCode 84)}, label={lst:c26-hist84-soln}]
class Solution {
public:
    int largestRectangleArea(vector<int>& heights) {
        int n = (int)heights.size();
        vector<int> st;                          // increasing stack of indices
        long long best = 0;
        for (int i = 0; i <= n; ++i) {           // i == n is the 0-height sentinel
            int cur = (i == n ? 0 : heights[i]);
            while (!st.empty() && heights[st.back()] >= cur) {
                long long h = heights[st.back()]; st.pop_back();
                int L = st.empty() ? -1 : st.back();
                best = max(best, h * (i - L - 1));
            }
            st.push_back(i);
        }
        return (int)best;
    }
};
\end{lstlisting}

\solstep{Trace}
\begin{tracebox}[Official sample -- heights \texttt{[2,1,5,6,2,3]} (answer 10)]
See \exref{c26-hist-run}. The decisive pop is bar index $2$ (height $5$) when the shorter bar index $4$
(height $2$) arrives: $L=1$ (the height-$1$ bar), width $4-1-1=2$, area $5\cdot2=10$. Every other
settling area is $\le 10$, so the answer is $10$ -- the red rectangle in \figref{c26-histfig}.
\end{tracebox}

\solstep{Complexity}
\begin{complexitybox}
Time $\bigO{n}$: each index pushed once, popped at most once, so the inner \code{while} totals $\le n$
iterations across the whole sweep. Memory $\bigO{n}$ for the stack. Verdict: \ac -- the linear histogram
algorithm. Its $\bigO{n^2}$ sibling (grow a window per bar) fails at $n=10^5$.
\end{complexitybox}

\solstep{Pitfalls and Variants}
\begin{warningbox}
Ways to fall: (i) forgetting the $i=n$ sentinel -- tall trailing bars never settle their width; (ii)
pushing \emph{heights} not \emph{indices} -- no way to measure the width (Trap~2); (iii) \code{int}
overflow in $h\cdot\text{width}$ when the grid is wide -- accumulate in \code{long long}. \emph{Variant
(LC 85, Maximal Rectangle):} treat each matrix row as a histogram of consecutive 1s above it and call
this routine per row -- the largest all-1 rectangle is the largest rectangle over those stacked
histograms.
\end{warningbox}

\begin{insightbox}[Key takeaways]
\begin{enumerate}[nosep, leftmargin=1.6em]
  \item \textbf{Pop time $=$ answer time.} An element's nearest beating neighbour is known the instant it
        is popped; anything computed from that neighbour (width, next-greater distance, basin) is a
        multiply or subtract away (\thmref{c26-hist}).
  \item \textbf{One loop, four dialects.} Increasing vs decreasing stack and what you do at the pop
        decide whether you get histogram, next-greater, sliding max, or rain water.
  \item \textbf{A deque is a stack that also pops the front.} That extra move is exactly the
        ``out-of-window'' eviction that turns the stack into a sliding-window extreme
        (\coderef{c26-deque}).
\end{enumerate}
\end{insightbox}

\begin{hardeningbox}
\begin{itemize}[nosep, leftmargin=1.4em]
  \item \textbf{Next smaller to both sides in one pass} -- the pop gives right bound, the new top the left
        bound; return both arrays and you have \lemref{c26-span}'s $L,R$ for a sum problem.
  \item \textbf{Maximum of minimums of every window size} (a classic) -- run the histogram span, then
        bucket each bar's area by its width and sweep from the right keeping the max.
  \item \textbf{Stock-span / online variants} -- feed the decreasing stack as data arrives; the same pop
        counts how far back the current element dominates.
\end{itemize}
\end{hardeningbox}

\subsection{Practice Problems}

\begin{stuckbox}
\begin{itemize}[nosep, leftmargin=1.4em]
  \item Did you append the sentinel (or equivalently flush the stack after the loop)?
  \item Did the stack hold \emph{indices} so widths are index differences?
  \item For next-greater: did you choose a \emph{decreasing} stack so the arriving bigger element answers
        the popped one?
\end{itemize}
\end{stuckbox}

\begin{practiceset}
  \practice{Maximal Rectangle}{LeetCode 85}
    {https://leetcode.com/problems/maximal-rectangle/}{LC Hard · L4}
    {Given an $n\times m$ binary matrix, return the area of the largest rectangle containing only 1s.
    Turn each row into a histogram: $\mathit{height}[j]$ is the number of consecutive 1s in column $j$
    up to and including this row (extend if this row's cell is 1, else reset to 0). Then run the solved
    \probref{c26-hist84} on that histogram; the best over all rows is the answer. The key realisation --
    a rectangle of 1s ending at row $r$ is exactly a histogram rectangle in $r$'s stacked skyline.}

  \practice{Trapping Rain Water}{LeetCode 42}
    {https://leetcode.com/problems/trapping-rain-water/}{LC Hard · L4}
    {Given bar heights, how much water is trapped after raining. Two linear routes: the \emph{two-pointer}
    method keeps a left-max and right-max and fills each cell by $\min(\text{leftMax},\text{rightMax})-\text{height}$;
    or the \emph{monotonic decreasing stack}: when a taller bar pops a lower one, the popped slot is a
    basin floor bounded left by the new stack top and right by the arriving bar -- fill width $\times$
    (min of the two walls $-$ floor). Both are $\bigO{n}$; the stack version is the same pop-decides-answer
    machinery as \thmref{c26-hist}.}

  \practice{Sliding Window Maximum}{LeetCode 239}
    {https://leetcode.com/problems/sliding-window-maximum/}{LC Hard · L4}
    {Return the maximum in each window of size $k$ as it slides across the array. Use the monotonic
    deque of \coderef{c26-deque}: keep indices with non-increasing values; push new $i$ after popping
    smaller-or-equal from the back, evict the front once its index $\le i-k$, and read the front once the
    window is full. Amortised $\bigO{1}$ per element -- a heap with lazy deletion is the $\bigO{n\log n}$
    fallback but the deque is optimal.}

  \practice{Daily Temperatures}{LeetCode 739}
    {https://leetcode.com/problems/daily-temperatures/}{LC Medium · L3}
    {For each day, how many days until a strictly warmer temperature; $0$ if none. This is next-greater
    distance: keep a \emph{decreasing} stack of indices; when a warmer day arrives, pop every cooler index
    and answer it as $i - \text{popped}$. Push each index once, pop once -- $\bigO{n}$. The cleanest
    possible ``the pop is the answer'' warm-up before LC 84.}
\end{practiceset}

\begin{recallbox}[10-second refresher -- 9.2 Monotonic Stack Applications]
A monotonic stack keeps its values ordered so that an element is popped the moment it meets the nearest
neighbour that beats it -- and the answer for that element is finalised \emph{at the pop}
(\defref{c26-monostack}, \thmref{c26-hist}). Increasing stack $\to$ nearest smaller (histogram,
\coderef{c26-hist}); decreasing stack $\to$ nearest greater (daily temperatures). Append a sentinel to
flush the stack. Add a front-eviction and the stack becomes a monotonic \emph{deque} for sliding-window
extremes (\coderef{c26-deque}, LC 239). Everything is amortised $\bigO{n}$. LC 84 is the anchor; LC 85,
42, 239, 739 are the dialects.
\selftest{
  \item Why does the width of a popped bar's rectangle become known exactly at its pop?
  \item Increasing or decreasing stack for "next greater element"? Why?
  \item What does the $\bigO{1}$ sliding-window max deque pop, and from which ends?
  \item Give the rain-water pop recipe: for a popped basin floor, what are the two walls?
  \item Why does forgetting the sentinel under-report areas on a rising skyline?
}
\end{recallbox}

% !TEX root = ../../tree_notes.tex
% =====================================================================
%  9.3  Range Minimum Query via the Cartesian Tree   sec:c26-rmq
%  Exports: thm:c26-rmqlca  lem:c26-cartmin  fig:c26-rmqfig
%           lst:c26-rmq  ex:c26-rmq-run
%           prob:c26-rmq1647  lst:c26-rmq1647-soln
% =====================================================================
\section{Range Minimum Query via the Cartesian Tree}\label{sec:c26-rmq}

\begin{previouslybox}
\begin{enumerate}[nosep, leftmargin=1.6em]
  \item What are the two orders of a Cartesian tree, and which is the root? (\defref{c26-cartesian})
  \item \lemref{c26-span}: node $i$'s subtree is the window on which $a[i]$ is the minimum. If a range
        $[l,r]$ has its minimum at position $m$, where do $l$ and $r$ sit relative to $m$'s window?
  \item In a heap, comparing a node to its ancestor: is the ancestor's value smaller or larger? (smaller --
        min-heap)
  \item Which structure from Chapter~16 (\chref{c11}) answers "lowest common ancestor of two nodes" after
        $\bigO{n\log n}$ preprocessing? (binary lifting)
\end{enumerate}
\end{previouslybox}

\begin{topicintro}
Here the Cartesian tree stops being a metaphor and becomes an \emph{index}. You are handed a fixed array
and many questions of the form ``what is the minimum of $a[l..r]$?'' -- the \term{range minimum query},
the static sibling of the segment tree. The punchline of this section is a one-line identity:

\begin{center}
\setlength{\fboxsep}{6pt}\fbox{\normalsize$\displaystyle\min\{a[l],\dots,a[r]\} \;=\; a\big[\lca(l,r)\big]$
\quad in the min-Cartesian tree.}
\end{center}

The minimum of any contiguous range is the \emph{value stored at the lowest common ancestor} of the two
endpoints, when the endpoints are read as nodes of the tree. Proving it is a three-line consequence of
\lemref{c26-span}; using it turns RMQ into an LCA call, and (this is the deeper equivalence Chapter~16
plays out) LCA is itself an RMQ on the Euler tour -- the two problems are the same problem in two
languages. We state and prove the reduction, build the tree, and answer queries with binary lifting,
explicitly \emph{handing the LCA engine over} to Chapter~16 rather than re-deriving it.
\textbf{Prerequisites.} \secref{c26-cartesian} (the span lemma is the entire proof); a black-box LCA
(binary lifting, Chapter~16, \chref{c11}).

\textbf{After this section you can}
\begin{itemize}[nosep]
  \item prove that a range minimum equals the value at the LCA of its endpoints in the Cartesian tree; \hfill\lv{3}--\lv{4}
  \item answer $q$ static RMQs in $\bigO{(n+q)\log n}$ by combining the $\bigO{n}$ build with lifted LCA; \hfill\lv{4}
  \item explain the two-way RMQ $\equiv$ LCA equivalence and say which chapter owns which direction. \hfill\lv{4}
\end{itemize}
\end{topicintro}

\begin{recognitionbox}[Pattern card -- static range extremes]
\reckey{a \emph{fixed} array, \emph{no} updates, and many "min/max over $[l,r]$" questions -- the offline
regime where the Cartesian tree doubles as a query index.}
\recuse{
  \item "minimum in range $[a,b]$", many queries, array never changes -- CSES 1647 (\probref{c26-rmq1647});
  \item "LCA of two nodes" dressed as an array problem, or vice-versa -- the equivalence is the tool;
  \item "is $a[m]$ the min of this window?" -- check $m$ against $\lca$ of the endpoints' nodes.
}
\recnotuse{
  \item the array is \emph{updated} between queries -- you need a segment tree (Chapter~18--19,
        \chref{c14}/\chref{c13}); the Cartesian tree is frozen at build time;
  \item a \emph{single} whole-array min -- one scan;
  \item "sum over a range" -- the min-heap order is what makes LCA work; sums have no such reduction, use
        a prefix array or segment tree.
}
One-line test: \emph{static array, range min, lots of queries -- and you would like to think in terms of
a tree? Then $\min = $ the value at an LCA, and the reduction is the whole solution.}
\end{recognitionbox}

\subsection{Intuition}

Take $A=(3,1,4,2)$ again. Ask for the minimum of the range covering positions $2$ through $3$: $\min(1,4)=1$,
attained at position $2$. Now look at the Cartesian tree \figref{c26-cartA}: the node for position $2$
(value $1$) is the \emph{root}, and positions $2$ and $3$ are both in its subtree, so their lowest common
ancestor is exactly position $2$ -- the minimum. Every range tells the same story: the position of its
smallest value is the node from which both endpoints are reachable on the way up, and no lower node works
because a lower node would have a \emph{larger} value (the heap points up toward smaller).

Why can there be no common ancestor below the minimum? In a min-heap the values only get \emph{smaller}
going up, so a deeper common ancestor would have to be smaller than the range's own minimum -- impossible,
since it stores a value from inside the range. Why is the minimum \emph{an} ancestor of both endpoints?
Because \lemref{c26-span} says the minimum's subtree is its dominance window, and every value in $[l,r]$
is $\ge$ that minimum, so $[l,r]$ lies entirely inside the minimum's window -- meaning both $l$ and $r$
are in its subtree. Ancestor of both, none deeper possible: it \emph{is} the LCA.

\begin{figure}[H]
\centering
\begin{tikzpicture}[diagram, scale=0.92, transform shape]
  \node[troot] (n2) at (0,0)  {$a_2{=}1$};
  \node[tleaf] (n1) at (-2,-1.5) {$a_1{=}3$};
  \node[tnode] (n4) at (2,-1.5)  {$a_4{=}2$};
  \node[tleaf] (n3) at (1,-3.0)  {$a_3{=}4$};
  \draw[tedgehl] (n2) -- (n1); \draw[tedgehl] (n2) -- (n4); \draw[tedgehl] (n4) -- (n3);
  % query [1,3]: endpoints a1 and a3, LCA a2
  \node[draw=warningred, thick, circle, inner sep=1pt, minimum size=10mm] at (n1) {};
  \node[draw=warningred, thick, circle, inner sep=1pt, minimum size=10mm] at (n3) {};
  \node[annot, align=left, font=\scriptsize] at (0,1.2)
    {query $\min\{a_1..a_3\}$: endpoints $1,3$ (red)};
  \node[annot, align=left, font=\scriptsize] at (-5.0,-2.2)
    {\textbf{$\lca(1,3)=2$}\\ $a[\lca]=1$\\ $=\min(3,1,4)$};
  \draw[flow, dashed] (n1) -- (n2); \draw[flow, dashed] (n3) -- (n4) -- (n2);
\end{tikzpicture}
\caption{The reduction on $A=(3,1,4,2)$: to find $\min\{a_1,a_2,a_3\}$ read the two endpoint nodes $1$
and $3$ (ringed), climb to their lowest common ancestor -- node $2$ -- and take its value, $1$. That is
$\thmref{c26-rmqlca}$; the dashed paths are the way up to the shared ancestor.}\label{fig:c26-rmqfig}
\end{figure}

\subsection{Definitions and Notation}

\begin{lemma}{The range minimum is the shared ancestor}{c26-cartmin}
Let $m$ be a position attaining the minimum of $a[l..r]$ in the min-Cartesian tree. Then $m$ is an
ancestor of both node $l$ and node $r$, and no strict descendant of $m$ is a common ancestor of $l$ and
$r$. Hence $\lca(l,r)=m$.
\end{lemma}
\begin{proofbox}
By \lemref{c26-span}, $m$'s subtree spans its dominance window $(L_m,R_m)$. Every element of $[l,r]$ is
$\ge a[m]$ (as $m$ is the range minimum), so no element of $[l,r]$ is a strictly smaller value that could
sit at or beyond $L_m,R_m$; therefore $[l,r]\subseteq(L_m,R_m)$ and $l,r$ both lie in $m$'s subtree --
i.e.\ $m$ is a common ancestor of $l,r$. For the second claim, $m$'s left subtree holds exactly the
indices $<m$ and its right subtree the indices $>m$ (BST-on-index order); a descendant of $m$ lies in one
of these subtrees and so can contain indices only from one side of $m$, never both a left endpoint
$l\le m$ and a right endpoint $r\ge m$ -- unless $m$ equals one endpoint, in which case $m$ is trivially
the LCA. Thus no proper descendant is common to $l$ and $r$, and $m$ is the \emph{lowest} common ancestor. \end{proofbox}

\begin{theorem}{RMQ reduces to LCA on the Cartesian tree}{c26-rmqlca}
Build the min-Cartesian tree of a static array $a[1..n]$. For every $1 \le l \le r \le n$,
$\min\{a[l],\dots,a[r]\} = a\big[\lca(l,r)\big]$, where $\lca$ is taken over the node set
$\{1,\dots,n\}$ of that tree. Consequently a range-minimum query is answered in the time of a single
lowest-common-ancestor query, and any LCA engine -- binary lifting (Chapter~16, \chref{c11}),
Euler-tour plus sparse table, or Farach-Colton--Bender -- becomes an RMQ engine for $a$.
\end{theorem}
\begin{proofbox}
Directly from \lemref{c26-cartmin}: the position $m=\arg\min$ over $[l,r]$ satisfies
$\lca(l,r)=m$, so the queried minimum $a[m]=a[\lca(l,r)]$. The identity holds for the endpoints $l,r$
regardless of the interior, so a static $\lca$ oracle answers the whole family of queries; building the
tree (\coderef{c26-stackbuild}) is $\bigO{n}$, so the total cost is (tree build) $+$ (LCA preprocessing)
$+$ (per-query LCA), exactly as the chosen LCA engine prices them. \end{proofbox}

\begin{notebox}[The other direction: LCA $\equiv$ RMQ]
The reduction runs both ways, and that is the point of the equivalence Chapter~16 makes precise: given a
tree, list its nodes by \emph{Euler-tour entry time}; the LCA of $u,v$ is the node of \emph{minimum depth}
between $\tin(u)$ and $\tin(v)$ in that tour -- an RMQ on the depth array. So "static array min over a
range" and "lowest common ancestor in a tree" are mutually reducible. This section uses
array$\to$LCA (build a tree, query it); Chapter~16 (\chref{c11}) uses tree$\to$array (Euler tour, then
RMQ) and owns the LCA machinery -- binary lifting, and the Euler-tour/sparse-table route -- that we call
as a black box. If you only ever need \emph{array} RMQ, the sparse table (Chapter~19, \chref{c13}) answers
each query in $\bigO{1}$ with $\bigO{n\log n}$ build and never builds a tree; the Cartesian reduction
matters for the conceptual bridge and when you already \emph{have} the tree.
\end{notebox}

\subsection{The algorithm: build the tree, then lift}

The pipeline is exactly \coderef{c26-stackbuild} feeding an LCA oracle. Record each node's parent during
the build, run one depth pass, fill the binary-lifting table, and answer each query by lifting
$\lca(l,r)$ and printing $a$ of that node.

\begin{lstlisting}[caption={Array RMQ via the Cartesian tree: stack build records parents, binary lifting answers LCA, the query is a[LCA]}, label={lst:c26-rmq}]
// a[1..n]; q queries (l, r); outputs min of a[l..r] via a[LCA(l,r)] in the min-Cartesian tree.
vector<int> par(n+1,0), Lc(n+1,0), Rc(n+1,0), st;
for (int i = 1; i <= n; ++i) {                     // the O(n) Cartesian build (lst:c26-stackbuild)
    int last = 0;
    while (!st.empty() && a[st.back()] > a[i]) { last = st.back(); st.pop_back(); }
    if (!st.empty()) { Rc[st.back()] = i; par[i] = st.back(); }
    if (last)      { Lc[i] = last; par[last] = i; }
    st.push_back(i);
}
int root = st.front();
vector<int> dep(n+1, 0);                            // depth, iterative from the root (no recursion)
vector<int> fs = {root};
while (!fs.empty()) {
    int u = fs.back(); fs.pop_back();
    if (Lc[u]) { dep[Lc[u]] = dep[u] + 1; fs.push_back(Lc[u]); }
    if (Rc[u]) { dep[Rc[u]] = dep[u] + 1; fs.push_back(Rc[u]); }
}
int LOG = 1; while ((1 << LOG) < n) ++LOG;
vector<vector<int>> up(LOG, vector<int>(n+1, 0));
for (int v = 1; v <= n; ++v) up[0][v] = par[v];
for (int k = 1; k < LOG; ++k)
    for (int v = 1; v <= n; ++v) up[k][v] = up[k-1][ up[k-1][v] ];
auto lca = [&](int u, int v) {                      // standard binary-lifting LCA
    if (dep[u] < dep[v]) swap(u, v);
    int d = dep[u] - dep[v];
    for (int k = 0; k < LOG; ++k) if (d >> k & 1) u = up[k][u];
    if (u == v) return u;
    for (int k = LOG - 1; k >= 0; --k)
        if (up[k][u] != up[k][v]) { u = up[k][u]; v = up[k][v]; }
    return par[u];
};
for (int qi = 0; qi < q; ++qi) { int l, r; /* read */ cout << a[ lca(l, r) ] << '\n'; }
\end{lstlisting}
\begin{complexitybox}
Time $\bigO{n\log n + q\log n}$: $\bigO{n}$ to build, $\bigO{n\log n}$ for the lifting table, $\bigO{\log
n}$ per query. Memory $\bigO{n\log n}$ for \code{up}. Verdict: correct and comfortably inside $n,q\le
2\cdot10^5$. If you will \emph{only ever} answer array RMQ, the sparse table gives $\bigO{1}$ per query at
the same $\bigO{n\log n}$ build (Chapter~19, \chref{c13}); the Cartesian route is the one to know when the
tree is already part of your structure or you want the LCA bridge.
\end{complexitybox}

\subsection{Worked Example}

\begin{example}{Three queries on $A=(3,1,4,2)$}{c26-rmq-run}
Tree \figref{c26-cartA}: root $a_2{=}1$, left $a_1$, right $a_4$ with left child $a_3$.\\
\emph{$[1,3]$:} $\lca(1,3)$. Node $1$'s ancestors: $\{1,2\}$; node $3$'s: $\{3,4,2\}$; lowest common $=2$,
value $1$. Indeed $\min(3,1,4)=1$.\\
\emph{$[3,4]$:} $\lca(3,4)$. Node $4$ is the parent of node $3$, so $\lca=4$, value $2$;
$\min(4,2)=2$.\\
\emph{$[1,4]$:} $\lca(1,4)=2$ (the root), value $1=\min(3,1,4,2)$. Each query is one LCA, and each LCA is
the range's minimum -- \thmref{c26-rmqlca}. \checkmark
\end{example}

\subsection{Common Pitfalls}\label{sec:c26-rmq-pitfalls}

\begin{warningbox}[Trap 1 -- taking the LCA of values instead of positions]
The tree's nodes \emph{are the indices} $1..n$; you query $\lca(l,r)$ on the \emph{endpoint positions}
and then read $a$ of the result. Passing the values $a[l],a[r]$ as node names is a category error: the
heap order is on values, but the tree structure (parent/child, and hence LCA) is on the index-labeled
nodes.
\end{warningbox}

\begin{warningbox}[Trap 2 -- a recursive depth pass on a skewed tree]
A sorted array builds a Cartesian tree of depth $n$; a recursive DFS to fill \code{dep} then overflows
the call stack (RE) at $n=2\cdot10^5$. Use the explicit-stack sweep in \coderef{c26-rmq}, or compute
depth during the build. The same caution that haunts any pointer tree on adversarial input, back in
Chapter~7 (\chref{c05}).
\end{warningbox}

\begin{warningbox}[Trap 3 -- reinventing LCA you were given]
Do not hand-roll the Euler-tour-plus-sparse-table or the lifting table from scratch in an RMQ problem if
Chapter~16 already supplies it -- that is the chapter's anti-repetition contract. Build the Cartesian
tree, look up $\lca$ through the provided engine, and output $a[\lca]$. The \emph{reduction} is this
section's contribution; the \emph{LCA} is not.
\end{warningbox}

\subsection{Problem Lab}

The anchor is CSES 1647 -- Static Range Minimum Queries -- the purest statement of the reduction: a
frozen array, $q$ ranges, print each minimum. We answer every query as $a[\lca]$.

% ---------------------------------------------------------------------------
\solvedproblem{c26-rmq1647}{Static Range Minimum Queries}{CSES 1647}
  {https://cses.fi/problemset/task/1647}{CSES · L3}
  {Cartesian tree build, binary-lifting LCA, RMQ equals value at LCA}

\begin{problemstatement}[Static Range Minimum Queries]
Given an array of $n$ integers, process $q$ queries of the form "what is the minimum value in range
$[a,b]$?". The array never changes. $1 \le n,q \le 2\cdot10^{5}$, $1 \le x_i \le 10^{9}$, and
$1 \le a \le b \le n$ (1-indexed). Print the answer to each query on its own line.
\end{problemstatement}

\begin{attemptgate}[Attempt it yourself first -- \lv{4}, give yourself 30 minutes]
You may already know the sparse-table answer; force yourself to answer it \emph{as the reduction} so the
bridge is real, not decorative.
\peek{Build the min-Cartesian tree of the array. For a query $[a,b]$, which single node holds the answer?
(\thmref{c26-rmqlca}: the LCA of $a$ and $b$.)}
\nudge{You need LCA, so record parents during the build, get depths, and answer LCA by binary lifting.
Then print $a[\lca(a,b)]$. Preprocess $\bigO{n\log n}$, per query $\bigO{\log n}$.}
\fullhint{Exactly \coderef{c26-rmq}: stack build for \code{par/Lc/Rc}, iterative depth, a \code{up[k][v]}
lifting table, standard lift-to-LCA, output $a[\lca]$. On the sample the four answers are $2,1,1,4$.}
\end{attemptgate}

\solstep{Understanding the Problem}
Static array, no updates, up to $2\cdot10^5$ range-min queries. A per-query linear scan is $\bigO{nq}
\approx 4\cdot10^{10}$ -- far too slow. The standard engine is a sparse table; but this section's lens is
the Cartesian tree, so we build that tree once and answer each query as a lowest-common-ancestor
computation, pricing it at $\bigO{(n+q)\log n}$ -- acceptable here and structurally illuminating.

\solstep{Key Observations}
\begin{observation}
By \thmref{c26-rmqlca}, $\min a[a..b] = a[\lca(a,b)]$ in the min-Cartesian tree, so the entire problem
reduces to answering LCA queries on a tree we can build in $\bigO{n}$.
\end{observation}
\begin{observation}
LCA via binary lifting needs only the parent array (free from the build, \coderef{c26-stackbuild}) and
depths; the lifting table costs $\bigO{n\log n}$ and each LCA $\bigO{\log n}$.
\end{observation}
\begin{observation}
No updates occur, so the frozen tree is valid for every query -- the exact regime where the reduction
beats a stack pass that could not remember edits.
\end{observation}

\solstep{Think It Through With Me}
\begin{thinkbox}
\deadend{Scan $[a,b]$ for each query.}
{$\bigO{n}$ per query, $\bigO{nq}$ total -- $\sim4\cdot10^{10}$ at the limits, a \tle.}
\deadend{Segment tree.}
{Works ($\bigO{(n+q)\log n}$), but that is the engine of Chapter~18 (\chref{c14}); here the goal is to
\emph{use the Cartesian tree as an index}, so we take the LCA route instead -- it is the one that proves
\thmref{c26-rmqlca} does work.}
\aha{Read the query endpoints as tree nodes, lift to their LCA, output that node's value. The tree is
already built by the stack; lifting is the only extra machinery, and it is Chapter~16's, not ours.}
\end{thinkbox}

\solstep{Solution}
As in \coderef{c26-rmq}: monotonic-stack build to get \code{par, Lc, Rc} and \code{root}; iterative depth
sweep from \code{root}; binary-lifting table \code{up}; for each query compute $c=\lca(a,b)$ and print
$a[c]$. Because \thmref{c26-rmqlca} guarantees $a[c]$ is the range minimum, the answers are correct
regardless of the interior of the query.

\solstep{Proof of Correctness}
\begin{proofbox}[Each printed value is the range minimum]
For query $[a,b]$ let $m$ be a position of the minimum of $a[a..b]$. By \lemref{c26-cartmin}, $m=\lca(a,b)$.
The algorithm prints $a[\lca(a,b)]=a[m]$, which is by definition the minimum of the range. The LCA
computation itself is the textbook binary lifting, correct because lifting the deeper node to equal depth
and then ascending both in lock-step from the top of the table stops at the child of the LCA, whose
parent is returned. Hence every output equals the true range minimum. \end{proofbox}

\solstep{Implementation}
\begin{lstlisting}[caption={Static Range Minimum Queries -- Cartesian build plus binary-lifting LCA (CSES 1647)}, label={lst:c26-rmq1647-soln}]
#include <bits/stdc++.h>
using namespace std;
int main(){
    ios::sync_with_stdio(false); cin.tie(nullptr);
    int n, q; cin >> n >> q;
    vector<long long> a(n + 1);
    for (int i = 1; i <= n; ++i) cin >> a[i];

    vector<int> par(n+1,0), Lc(n+1,0), Rc(n+1,0), st;
    for (int i = 1; i <= n; ++i) {
        int last = 0;
        while (!st.empty() && a[st.back()] > a[i]) { last = st.back(); st.pop_back(); }
        if (!st.empty()) { Rc[st.back()] = i; par[i] = st.back(); }
        if (last)      { Lc[i] = last; par[last] = i; }
        st.push_back(i);
    }
    int root = st.front();

    vector<int> dep(n+1, 0);
    vector<int> fs = {root};
    while (!fs.empty()) {
        int u = fs.back(); fs.pop_back();
        if (Lc[u]) { dep[Lc[u]] = dep[u] + 1; fs.push_back(Lc[u]); }
        if (Rc[u]) { dep[Rc[u]] = dep[u] + 1; fs.push_back(Rc[u]); }
    }

    int LOG = 1; while ((1 << LOG) < n) ++LOG;
    vector<vector<int>> up(LOG, vector<int>(n+1, 0));
    for (int v = 1; v <= n; ++v) up[0][v] = par[v];
    for (int k = 1; k < LOG; ++k)
        for (int v = 1; v <= n; ++v) up[k][v] = up[k-1][ up[k-1][v] ];

    auto lca = [&](int u, int v) {
        if (dep[u] < dep[v]) swap(u, v);
        int d = dep[u] - dep[v];
        for (int k = 0; k < LOG; ++k) if ((d >> k) & 1) u = up[k][u];
        if (u == v) return u;
        for (int k = LOG - 1; k >= 0; --k)
            if (up[k][u] != up[k][v]) { u = up[k][u]; v = up[k][v]; }
        return par[u];
    };

    while (q--) { int A, B; cin >> A >> B; cout << a[ lca(A, B) ] << '\n'; }
    return 0;
}
\end{lstlisting}

\solstep{Trace}
\begin{tracebox}[Official CSES sample: array \texttt{3 2 4 5 1 1 5 3}]
The stack build (popping on strict $>$, so the two equal 1s are kept distinct) makes node $5$ the root,
with left subtree the block $1..4$ and right subtree $6..8$. Then: query $(2,4)$: node $2$ is an ancestor
of node $4$ (path $2\to3\to4$), so $\lca=2$, $a[2]=2$. Query $(5,6)$: $\lca(5,6)=5$ (root), $a[5]=1$.
Query $(1,8)$: $\lca(1,8)=5$ (root), $a[5]=1$. Query $(3,3)$: $\lca(3,3)=3$, $a[3]=4$. Output $2,1,1,4$
-- matches CSES. \checkmark
\end{tracebox}

\solstep{Complexity}
\begin{complexitybox}
Time $\bigO{(n+q)\log n}$: $\bigO{n}$ build, $\bigO{n\log n}$ lifting table, $\bigO{\log n}$ per query;
with $n,q\le2\cdot10^5$ this is a few million operations. Memory $\bigO{n\log n}$ for the table
($\approx 2\cdot10^5 \times 18$ integers, a few MB). Verdict: \ac. Note the sparse table (Chapter~19,
\chref{c13}) reaches $\bigO{1}$ per query with the same build -- use it when you need array RMQ alone; the
LCA reduction is the tool when the tree is already there.
\end{complexitybox}

\solstep{Pitfalls and Variants}
\begin{warningbox}
Watch the same three traps: query the \emph{positions} not the values; keep the depth sweep iterative so a
sorted input (a chain of depth $n$) does not overflow recursion; and do not re-derive LCA -- import it.
\emph{Variants:} for a range \emph{maximum} build the max-Cartesian tree (flip the pop to $<$) and the
identity becomes $\max = a[\lca]$. If updates appear, abandon the frozen tree for a segment tree
(\chref{c14}). If you need $\bigO{1}$ queries, replace the LCA table with a sparse table over $a$
directly (\chref{c13}) -- but then you have quietly left this chapter's reduction behind.
\end{warningbox}

\begin{insightbox}[Key takeaways]
\begin{enumerate}[nosep, leftmargin=1.6em]
  \item \textbf{A range minimum is a lowest common ancestor} -- $\min a[l..r]=a[\lca(l,r)]$
        (\thmref{c26-rmqlca}), a one-line corollary of the span lemma.
  \item \textbf{RMQ and LCA are the same problem} in two encodings; this section is
        array$\to$LCA, Chapter~16 is tree$\to$RMQ via the Euler tour.
  \item \textbf{The tree is an index, not just an image} -- the $\bigO{n}$ build \emph{plus} any LCA oracle
        is a working static-RMQ engine, which is how CSES 1647 falls (\probref{c26-rmq1647}).
\end{enumerate}
\end{insightbox}

\begin{hardeningbox}
\begin{itemize}[nosep, leftmargin=1.4em]
  \item \textbf{Second-smallest in a range} -- once you have the tree, the answer is the smaller of the
        two children along the $\lca$ path; a common contest twist on static RMQ.
  \item \textbf{Cartesian-tree DFS orders} -- inorder recovers the array; a preorder gives a "smallest
        first" enumeration used by some offline union tricks.
  \item \textbf{Range mode / nearest smaller} -- both read the same $L,R$ arrays from \secref{c26-monostack};
        the tree unifies them.
\end{itemize}
\end{hardeningbox}

\begin{recallbox}[10-second refresher -- 9.3 RMQ via the Cartesian Tree]
In the min-Cartesian tree the minimum of any contiguous range $[l,r]$ is stored at $\lca(l,r)$
(\thmref{c26-rmqlca}, from the span lemma \lemref{c26-cartmin}). So static RMQ reduces to LCA: build the
tree in $\bigO{n}$ (\coderef{c26-stackbuild}), record parents, lift to the LCA per query
(\coderef{c26-rmq}), print $a[\lca]$. The reverse reduction (LCA $\to$ RMQ on the Euler-tour depth array)
is owned by Chapter~16 (\chref{c11}); a pure array-only $\bigO{1}$-query RMQ is the sparse table of
Chapter~19 (\chref{c13}). CSES 1647 (\probref{c26-rmq1647}) is the worked reduction.
\selftest{
  \item Why is the range's minimum position an \emph{ancestor} of both endpoints?
  \item Why can no \emph{deeper} node be their common ancestor?
  \item State the RMQ $\equiv$ LCA equivalence and who owns which direction.
  \item Give the preprocessing and per-query costs of the lifting route.
  \item When should you prefer a sparse table to this reduction?
}
\end{recallbox}

% !TEX root = ../../tree_notes.tex
% =====================================================================
%  9.4  The Cartesian Tree as a Random Treap     sec:c26-treapview
%  Exports: prop:c26-carters-treap  thm:c26-randbst  fig:c26-iso
%           ex:c26-treapview-run
%           prob:c26-mbt2  lst:c26-mbt2-soln
% =====================================================================
\section{The Cartesian Tree as a Random Treap}\label{sec:c26-treapview}

\begin{previouslybox}
\begin{enumerate}[nosep, leftmargin=1.6em]
  \item A treap is a BST on \emph{keys} and a heap on \emph{priorities}. Given a fixed key set and a
        fixed priority per key, how many treaps realise them? (\chref{c06}, \defref{c06-treap} -- exactly
        one, when priorities are distinct)
  \item What does \thmref{c06-exptheight} guarantee for a treap with i.i.d.\ random priorities, and over
        what is the expectation taken? ($\bigO{\log n}$ height, over the random priorities)
  \item A Cartesian tree is a BST on which quantity and a heap on which? (\secref{c26-cartesian}: BST on
        \emph{index}, heap on \emph{value})
  \item Build a BST by inserting $3,1,4,2$ in that order. Is the root $3$ (the first key) or $1$ (the
        smallest)? (the first inserted -- $3$)
\end{enumerate}
\end{previouslybox}

\begin{topicintro}
This last section is a lens, not new machinery. Read a Cartesian tree one way and it is ``BST on index,
heap on value.'' Read the \emph{same} two-orders definition with the second order called a priority, and
it is \emph{exactly a treap} (\defref{c06-treap}) whose keys are the positions and whose priorities are
the array values. The tree of \secref{c26-cartesian} and the tree of Chapter~8 were never different
objects; only the story about what the heap-order ``means'' changed. Flipping the perspective once more
-- let the \emph{priority be the insertion time} instead of the value -- recovers a third classic: the
shape of a plain BST built by inserting a sequence is a Cartesian tree, so ``insert a random permutation
into a BST'' is ``build a Cartesian tree with random priorities,'' whose height is logarithmic \emph{in
expectation} by Chapter~8's theorem, quoted and never re-proved (\thmref{c26-randbst}).

We make the equivalence precise (\propref{c26-carters-treap}), show all three readings of one picture
(\figref{c26-iso}), and then cash it on LeetCode~998: append a value to the far right of a
maximum-tree's underlying array and re-root it -- which the treap view makes a one-line walk down the
right spine. \textbf{Prerequisites.} \secref{c26-cartesian} (the tree) and Chapter~8's treap definition
and height theorem (\chref{c06}) \emph{as black boxes}.
\textbf{After this section you can}
\begin{itemize}[nosep]
  \item state and prove that a Cartesian tree \emph{is} a treap with value-priorities, and vice-versa; \hfill\lv{3}--\lv{4}
  \item explain why inserting a random permutation yields a random treap and cite the expected height; \hfill\lv{4}
  \item solve right-end maximum-tree insertions (\probref{c26-mbt2}) from the priority viewpoint; \hfill\lv{4}
\end{itemize}
\end{topicintro}

\begin{recognitionbox}[Pattern card -- the same tree, wearing priorities]
\reckey{a structure that is ``ordered by one key and heap-ordered by a second number''; or ``insert these
into a BST, what shape results?''; or a \emph{max/min-tree} you must extend at one end.}
\recuse{
  \item "rebuild a maximum binary tree after appending a value" -- LC 998 (\probref{c26-mbt2}), the treap
        right-spine view;
  \item "expected depth / comparisons when inserting a random order into a BST" -- \thmref{c26-randbst},
        reduce to Chapter~8's bound;
  \item any treap you build by hand is secretly a Cartesian tree over the (key, priority) pairs.
}
\recnotuse{
  \item you need to \emph{maintain} balance under adversarial edits -- that is a real balanced tree with
        rotations (Chapter~8); priorities here are read off, not chosen to force balance;
  \item you want the treap \code{split}/\code{merge} engine -- that lives in \secref{c06-treaps}, this
        section borrows its \emph{definition} only;
  \item priorities must be \emph{random and independent} -- the array values here are given, not sampled.
}
One-line test: \emph{BST on one field, heap on another -- and the second field is really a priority? Then
you are looking at a treap in a Cartesian tree's clothes, and Chapter~8's guarantees apply for free.}
\end{recognitionbox}

\subsection{Intuition}

The definition \defref{c26-cartesian} demanded two things: inorder $=$ the index order, and a heap on the
values. The treap definition \defref{c06-treap} demands two things: a BST on the \emph{keys}, and a heap on
the \emph{priorities}. Set $\text{key} = \text{index}$ and $\text{priority} = \text{value}$ and the two
demands are word-for-word identical -- a BST on the index is precisely ``inorder reads the indices in
order,'' and a heap on the value is the Cartesian heap rule. Since Chapter~8 already proved a treap with
distinct priorities is \emph{unique}, that is a second proof of \thmref{c26-unique}: the Cartesian tree is
unique because it is a treap.

Now rotate the viewpoint. Suppose instead you do not care about values at all; you just \emph{insert keys
into a plain BST} in some arrival order. The element inserted first becomes the root; the ones after it
fall left or right by size, and recursively each subsequence rebuilds the same way. That is a Cartesian
tree whose \emph{priority is the insertion time} (smallest time $=$ highest in the tree). If the arrival
order is a uniformly random permutation, the insertion times are i.i.d.-in-effect random priorities -- so
the resulting BST has the \emph{same distribution as a random treap}, and everything Chapter~8 proved about
treap height transfers verbatim. That is \thmref{c26-randbst}, and its proof is a single sentence plus a
citation -- exactly the anti-repetition this chapter promised.

\begin{figure}[H]
\centering
\begin{tikzpicture}[diagram, scale=0.8, transform shape]
  % (1) min-Cartesian tree of A, heap on value
  \node[troot]  (c2) at (0,0)   {$1$};
  \node[tleaf]  (c1) at (-1.6,-1.2) {$3$};
  \node[tnode]  (c4) at (1.6,-1.2)  {$2$};
  \node[tleaf]  (c3) at (0.8,-2.4)   {$4$};
  \draw[tedgehl] (c2)--(c1); \draw[tedgehl] (c2)--(c4); \draw[tedgehl] (c4)--(c3);
  \node[annot, font=\scriptsize\bfseries] at (0,1.0) {(a) Cartesian: heap on value};
  % (2) treap: key=index (shown), priority=value (subscript)
  \node[troot]  (t2) at (6.5,0)  {$\mathit{2}$};
  \node[tleaf]  (t1) at (4.9,-1.2) {$\mathit{1}$};
  \node[tnode]  (t4) at (8.1,-1.2) {$\mathit{4}$};
  \node[tleaf]  (t3) at (7.3,-2.4) {$\mathit{3}$};
  \draw (t2)--(t1); \draw (t2)--(t4); \draw (t4)--(t3);
  \node[annot, font=\scriptsize] at (4.9,-1.2){\phantom{x}};
  \node[font=\scriptsize, text=darkgray] at (6.5,0.5) {\scriptsize p1};
  \node[annot, font=\scriptsize\bfseries] at (6.5,1.0) {(b) Treap: key=index, prio=value};
  \node[annot, align=left, font=\scriptsize] at (10.3,-1.2)
    {nodes $1..4$ by key;\\ priorities $a_2{=}1<a_4{=}2<a_1{=}3$\\ $\Rightarrow$ root $=$ key $2$};
  % (3) insertion-order BST (priority=index), array order 3,1,4,2 -> keys are VALUES, heap on index
  \node[troot]  (i3) at (0,-5.0) {$3$};
  \node[tnode]  (i1) at (-1.6,-6.2) {$1$};
  \node[tleaf]  (i4) at (1.6,-6.2) {$4$};
  \node[tleaf]  (i2) at (-0.8,-7.4) {$2$};
  \draw (i3)--(i1); \draw (i3)--(i4); \draw (i1)--(i2);
  \node[annot, font=\scriptsize\bfseries] at (0,-4.0) {(c) BST by insertion order $3,1,4,2$};
  \node[annot, align=left, font=\scriptsize] at (3.2,-6.2)
    {root $=$ first in $=$ value $3$\\ (min index). BST on value,\\ heap on insertion time $=$ index:\\ \emph{a different} Cartesian tree\\ of the same array};
\end{tikzpicture}
\caption{One array $A=(3,1,4,2)$, three readings. (a) min-Cartesian tree -- inorder $=$ index, heap on
value (\defref{c26-cartesian}). (b) the \emph{identical} tree as a treap: keys are the indices $1..4$
(labels here), priorities are the values, root $=$ least-priority key $=2$ (\defref{c06-treap});
(a)$\equiv$(b) is \propref{c26-carters-treap}. (c) is the \emph{other} assignment -- a plain BST built by
inserting $3,1,4,2$, which is a Cartesian tree keyed by value with priority $=$ index, the object that
makes a random permutation a random treap (\thmref{c26-randbst}).}\label{fig:c26-iso}
\end{figure}

\subsection{The equivalence}

\begin{proposition}{Cartesian tree $=$ treap; insertion BST $=$ Cartesian by time}{c26-carters-treap}
Let $a[1..n]$ have distinct values.
\begin{itemize}[nosep]
  \item[(i)] The min-Cartesian tree of $a$ is exactly the treap whose keys are $\{1,\dots,n\}$ and whose
        priority of key $i$ is $a[i]$. Conversely that treap is the Cartesian tree.
  \item[(ii)] Build a BST by inserting the values $a[1],\dots,a[n]$ in that \emph{arrival} order. The
        resulting tree is the Cartesian tree keyed by \emph{value} with priority equal to the insertion
        index $i$ (min-heap on $i$, since the earliest-inserted key sits highest).
\end{itemize}
\end{proposition}
\begin{proofbox}
(i) A treap is a BST on keys and a heap on priorities (\defref{c06-treap}). With key $=i$ and priority
$=a[i]$, "BST on keys" says the inorder lists keys ascending -- i.e.\ $1,2,\dots,n$, the Cartesian
inorder condition; "min-heap on priorities" says a node's priority $a[i]$ is $\le$ its children's, the
Cartesian heap condition. So the treap satisfies \defref{c26-cartesian} and, both being unique under
distinct values/priorities, they are the same tree. (ii) The first key inserted, $a[1]$, is the root (no
key preceded it). Recursively, within any subtree the \emph{earliest} member of that key-block to arrive is
its local root, and arrivals split left/right purely by comparison -- so the shape is a Cartesian tree on
the key set where the priority is arrival time (index), and the root is the minimum-index member: a
min-heap on index. Distinct priorities (indices) give uniqueness as in (i). \end{proofbox}

\noindent\textbf{In plain words.} Swap the roles of the two orders and the same picture is a Cartesian
tree, a value-priority treap, or an insertion-time BST. Chapter~8's treap uniqueness \emph{is}
\thmref{c26-unique}.

\begin{theorem}{A random permutation builds a random BST of logarithmic height}{c26-randbst}
Insert a uniformly random permutation of $n$ distinct keys into an initially empty binary search tree.
By \propref{c26-carters-treap}(ii) the resulting tree is a Cartesian tree whose priorities are the
permutation's arrival positions -- i.e.\ a treap with i.i.d.-in-effect distinct random priorities. Its
expected height is therefore $\bigO{\log n}$, and its expected search cost is $\bigO{\log n}$; both are
the direct content of Chapter~8's \thmref{c06-exptheight}, invoked as a black box and \emph{not}
re-derived here.
\end{theorem}
\begin{proofbox}
A uniform random permutation assigns to each key an arrival time that is a uniformly random injection
into $\{1,\dots,n\}$ -- the same as sampling distinct priorities in random order. By
\propref{c26-carters-treap}(ii) the insertion-order BST is the treap on those (key, priority) pairs, so its
distribution equals that of a random treap. \thmref{c06-exptheight} states a random treap has expected
height $c\,\ln n$ for a constant $c$ (and, more strongly, is \emph{balanced} with high probability). Since
a BST's every operation costs $\bigO{h}$ (\secref{c05-insert-delete}), expected cost is $\bigO{\log n}$. No
new proof is needed -- the reduction is the whole argument, which is why this lives in a section rather
than a chapter. \end{proofbox}

\begin{insightbox}
Three chapters now share one sentence. Chapter~8 chose priorities \emph{at random} to force a short tree
you can edit. This chapter's array \emph{hands you} priorities (the values), and the tree is whatever
those dictate -- which you can still exploit (histogram, RMQ) because a monotonic stack reads the same
priorities faster than any rotation. The unifying idea: \textbf{a BST-plus-heap on a second field is one
structure, and what you call the second field -- value, priority, or arrival time -- only picks the story.}
\end{insightbox}

\subsection{Worked Example}

\begin{example}{Reading $A$ as a treap and as an insertion BST}{c26-treapview-run}
For $A=(3,1,4,2)$: \figref{c26-iso}(a) and (b) are the same tree. As a \emph{treap} (b): keys $1,2,3,4$
with priorities $3,1,4,2$. The least priority is key $2$ (priority $1$), so it is the root; keys less
than $2$ (just key $1$) go left, keys greater ($3,4$) go right, and among them key $4$ (priority $2$) is
rooted above key $3$ (priority $4$) -- identical to the min-Cartesian build, as \propref{c26-carters-treap}
promises. As an \emph{insertion BST} (c): insert $3$, then $1$ (goes left of $3$), then $4$ (right of
$3$), then $2$ (left of $3$, right of $1$). Root $3$ because it arrived first (smallest index). Now the
keys are the \emph{values} and the priority is the \emph{index}: root $3$ has index $1$, its subtree keys
$\{1,2\}$ have larger indices, and $1$ (index $2$) precedes $2$ (index $4$), so $1$ is $2$'s parent. Same
rule, relabelled -- \propref{c26-carters-treap}(ii) on the nose. \checkmark
\end{example}

\subsection{Common Pitfalls}\label{sec:c26-treapview-pitfalls}

\begin{warningbox}[Trap 1 -- re-proving the treap height bound]
\thmref{c26-randbst}'s height statement is \emph{imported} from \thmref{c06-exptheight}; the anti-repetition
contract forbids rebuilding it. The only new content here is the equivalence \propref{c26-carters-treap}.
If you find yourself writing a fresh "expected height of random BST" proof, you have left the chapter.
\end{warningbox}

\begin{warningbox}[Trap 2 -- confusing which field is the priority]
In reading (a)/(b) the \emph{value} is the priority and the \emph{index} is the key (min-Cartesian /
value-treap). In reading (c) the \emph{index (arrival time)} is the priority and the \emph{value} is the
key. Both are valid Cartesian trees but of \emph{different} shapes -- \figref{c26-iso}(a) and (c) are not
the same tree. State which field is the priority before arguing anything.
\end{warningbox}

\begin{warningbox}[Trap 3 -- choosing priorities to force balance]
A treap picks priorities \emph{randomly} to \emph{guarantee} short height against adversarial keys. A
Cartesian tree gets its priorities \emph{from the data} and makes no such promise: a sorted array gives a
value-treap of height $n$ (a chain) -- the very worst case the randomised treap avoids. Never claim a
read-off Cartesian tree is balanced; only the \emph{random-permutation} BST of \thmref{c26-randbst} is
balanced in expectation.
\end{warningbox}

\subsection{Problem Lab}

The anchor is LeetCode~998 -- given the root of a \emph{maximum tree} (a max-Cartesian tree) and a value
appended to the far right of its hidden array, return the new maximum tree. The priority view makes it
trivial: the new value is the \emph{largest key} (last position), so it descends only along right
children, becoming a treap insert restricted to the right spine.

% ---------------------------------------------------------------------------
\solvedproblem{c26-mbt2}{Maximum Binary Tree II}{LeetCode 998}
  {https://leetcode.com/problems/maximum-binary-tree-ii/}{LC Medium · L3}
  {max-Cartesian / treap view, right-spine descent, priority as value}

\begin{problemstatement}[Maximum Binary Tree II]
A \textbf{maximum tree} is a tree where every node's value is greater than any other value in its
subtree. You are given the root of a maximum tree built by $\text{Construct}(a)$: the root is
$\max(a)$, its left child is $\text{Construct}(\text{left part})$ and its right child
$\text{Construct}(\text{right part})$. You are not given $a$, only the root, and an integer $\mathit{val}$.
Let $b$ be $a$ with $\mathit{val}$ \emph{appended to the end} (it is guaranteed $b$ has distinct values).
Return $\text{Construct}(b)$. The number of nodes is in $[1,100]$ and all values are $\le 100$.
\end{problemstatement}

\begin{attemptgate}[Attempt it yourself first -- \lv{4}, give yourself 20 minutes]
Recovering $a$ by an inorder walk and rebuilding is allowed but $\bigO{n}$-with-a-full-rebuild and
obscures the structure. The treap view is one recursion.
\peek{$\mathit{val}$ is appended, so it is the \emph{largest index} in $b$. In the value-priority max-treap,
where can a new largest key ever be attached?}
\nudge{Only down the right spine: every existing key is smaller, so the new node descends through right
children until it finds a smaller node -- it then becomes that node's parent, with the old right subtree
as its left child. If it beats the root, it becomes the new root with the whole tree on its left.}
\fullhint{$\text{insert}(t,\mathit{val})$: if $t$ is null or $\mathit{val}>t.\mathit{val}$, return
$\text{node}(\mathit{val},\ \text{left}=t)$; else $t.\text{right}\leftarrow\text{insert}(t.\text{right},
\mathit{val})$ and return $t$. One pass down the right spine, $\bigO{h}$ worst case, $\bigO{n}$ here.}
\end{attemptgate}

\solstep{Understanding the Problem}
A maximum tree is a max-Cartesian tree of $a$: inorder $=$ array order, heap $=$ max on value --
\propref{c26-carters-treap} with the heap flipped and values as keys-priority. Appending $\mathit{val}$
adds the \emph{rightmost} index. The BST-order rule then forces $\mathit{val}$ to be $\ge$ (in index) of
every existing key, so its descent path can only ever take \emph{right} children; and the max-heap rule
makes it stop at the first node whose value is below $\mathit{val}$, becoming that node's parent and
adopting that node's (smaller-valued, larger-indexed) right subtree as its own left child.

\solstep{Key Observations}
\begin{observation}
Because $\mathit{val}$ is a new maximum \emph{index}, it is compared only against right-spine nodes --
never descends left -- so the insertion is a walk down the right spine.
\end{observation}
\begin{observation}
$\mathit{val}$ becomes the parent of the first right-spine node smaller than it (or the root's parent if
it beats the root), and inherits that node's existing right subtree as its left child -- exactly a
treap-priority insertion for the largest key.
\end{observation}
\begin{observation}
This is the recursive one-liner: null-or-$\mathit{val}>$ root $\Rightarrow$ new root with old tree on the
left; else recurse into the right child.
\end{observation}

\solstep{Think It Through With Me}
\begin{thinkbox}
\deadend{Inorder-dump the tree to recover $a$, append $\mathit{val}$, rebuild from scratch.}
{Correct and $\bigO{n}$ (LC 654-style $\text{Construct}$), but it rebuilds nodes that were already right
and misses the point -- appending the largest key only disturbs the right spine. Fine as a fallback, not
as understanding.}
\deadend{Insert $\mathit{val}$ as a plain BST by value, then re-heap.}
{Two trees -- BST order here is on the \emph{index}, not the value; a value-only descent lands $\mathit{val}$
in the wrong place relative to equal-index constraints.}
\aha{Descend the right spine only (the new node is the largest index); stop at the first node smaller than
$\mathit{val}$; make $\mathit{val}$ its parent with the node as $\mathit{val}$'s left child. One recursion,
$\bigO{n}$ worst, and it is a treap insert wearing our \propref{c26-carters-treap} lens.}
\end{thinkbox}

\solstep{Solution}
Define $\text{insert}(t, \mathit{val})$: if $t$ is null or $\mathit{val} > t\to\mathit{val}$, create a node
for $\mathit{val}$ whose left child is $t$ (the entire subtree $t$ lies to the left of the new largest
index and below the new max value) and return it; otherwise set $t\to\mathit{right} =
\text{insert}(t\to\mathit{right}, \mathit{val})$ and return $t$. Call it on the root. Because $\mathit{val}$
is the maximum index, the "$t$ is null" case is only reached at the end of the right spine, and the
"$\mathit{val} > t\to\mathit{val}$" case is the first right-spine node smaller than $\mathit{val}$.

\solstep{Proof of Correctness}
\begin{proofbox}[The right-spine descent yields $\text{Construct}(b)$]
Let $u$ be the first node on the right spine of $\text{Construct}(a)$ whose value is below $\mathit{val}$,
or the null after the spine's last node if all are larger. Every right-spine node above $u$ (closer to the
root) has value $>\mathit{val}$, so by the max-heap rule they stay ancestors of $\mathit{val}$; $u$ and its
whole subtree lie in a range of indices strictly left of $\mathit{val}$'s new index, so $\mathit{val}$, a
larger value and larger index, must become the parent of $u$, keeping $u$ as its left child (the
larger-index side of $\mathit{val}$ is empty since it is last). The recursion returns a new node $v$ with
$v\to\mathit{left}=u$, grafted as $u\to\mathit{parent}\to\mathit{right}$ -- precisely the tree
$\text{Construct}(b)$, whose root/max over $[u\text{ side},\dots,\mathit{val}]$ is $\mathit{val}$. When
$\mathit{val}$ beats the root, $u=$ root and $v$ becomes the new root with the old tree on its left -- again
matching $\text{Construct}(b)$ since $\mathit{val}$ is then $\max(b)$. The algorithm's output is
$\text{Construct}(b)$ in all cases. \end{proofbox}

\solstep{Implementation}
\begin{lstlisting}[caption={Maximum Binary Tree II -- one descent down the right spine (LeetCode 998)}, label={lst:c26-mbt2-soln}]
class Solution {
public:
    TreeNode* insertIntoMaxTree(TreeNode* root, int val) {
        if (!root || root->val < val)          // val is a new max index: only the right spine matters
            return new TreeNode(val, root, nullptr);  // old tree hangs on the new node's LEFT
        root->right = insertIntoMaxTree(root->right, val);  // descend the right spine
        return root;
    }
};
\end{lstlisting}

\solstep{Trace}
\begin{tracebox}[Official sample 3: root \texttt{[5,2,3,null,1]}, \texttt{val = 4}]
$a=[2,1,5,3]$; root $5$ (left $2$, which has right child $1$; right $3$). Call on root $5$: $4>5$? no, so
recurse right on node $3$. At node $3$: $4>3$? yes -- return a new node $4$ with left child $3$. Attach:
$5\to$right $=4$ (with $4\to$left $=3$). Result tree: root $5$, left $2$ (right child $1$), right $4$ (left
child $3$), i.e.\ level order $[5,2,4,\text{null},1,3]$ -- exactly the expected output. $\mathit{val}=4$
descended only through right children, as \propref{c26-carters-treap} guarantees for a largest index. \checkmark
\end{tracebox}

\solstep{Complexity}
\begin{complexitybox}
Time $\bigO{h}$ where $h$ is the tree height: the recursion follows the right spine once. Worst case
$\bigO{n}$ (a right-going chain); on a random maximum tree it is $\bigO{\log n}$ by
\thmref{c26-randbst}/\thmref{c06-exptheight}. Memory $\bigO{h}$ for the call stack (turn it into a loop
over $t\to$right if a deep tree risks RE, exactly the Trap~2 caution from Chapter~7). Verdict: \ac.
\end{complexitybox}

\solstep{Pitfalls and Variants}
\begin{warningbox}
The base case is ``$\mathit{val} > \mathit{root}\to$val'' \emph{or} root null -- both mean "make a new node
here with the current subtree on its left", never on its right (the current subtree is at smaller indices,
and $\mathit{val}$ is the last index). Confusing left/right re-parents the old tree wrongly. \emph{Variant:}
inserting at the \emph{front} would descend the \emph{left} spine by the identical argument; and LC 654
(build a max tree from an array) is the divide-by-maximum construction whose stack build is
\coderef{c26-stackbuild} flipped to a max-heap -- the same tree from the other side.
\end{warningbox}

\begin{insightbox}[Key takeaways]
\begin{enumerate}[nosep, leftmargin=1.6em]
  \item \textbf{Cartesian tree $=$ treap} once you rename value $\to$ priority and index $\to$ key
        (\propref{c26-carters-treap}); Chapter~8's uniqueness gives \thmref{c26-unique} for free.
  \item \textbf{Random permutation $\Rightarrow$ random BST $\Rightarrow$ log height} -- import
        \thmref{c06-exptheight}, never re-prove it (\thmref{c26-randbst}).
  \item \textbf{Appending the max index confines you to one spine} -- that is the whole of LeetCode~998
        (\probref{c26-mbt2}), one recursion deep.
\end{enumerate}
\end{insightbox}

\begin{hardeningbox}
\begin{itemize}[nosep, leftmargin=1.4em]
  \item \textbf{Prefer the loop} -- rewrite the recursion as an explicit walk down $t\to$right to dodge the
        $\bigO{h}$ call stack on a chain, the same hardening as any Chapter-7 BST descent.
  \item \textbf{Front insertion} -- mirror the argument on the left spine; a good drill to be sure you know
        which spine a given key belongs to.
  \item \textbf{Treap by hand} -- build a treap by inserting (key, random priority) pairs and confirm you
        get the Cartesian tree \emph{without a single rotation}: priorities, not rotations, set the shape.
\end{itemize}
\end{hardeningbox}

\begin{recallbox}[10-second refresher -- 9.4 The Cartesian Tree as a Random Treap]
Rename the Cartesian tree's value to a \emph{priority} and its index to a \emph{key} and it is literally a
treap (\propref{c26-carters-treap}); conversely a plain BST built by an insertion order is a Cartesian tree
with priority $=$ arrival time. So a random permutation is a random treap, and Chapter~8's
\thmref{c06-exptheight} hands you expected $\bigO{\log n}$ height for free (\thmref{c26-randbst}). No new
balance proof, ever -- that is the point. The payoff problem, LC~998 (\probref{c26-mbt2}), is a treap
insert of a new largest key down the right spine.
\selftest{
  \item Give the two field-name swaps that turn a Cartesian tree into a treap, and vice-versa.
  \item Why does \propref{c26-carters-treap} re-prove \thmref{c26-unique}?
  \item Why is a random-permutation BST balanced in expectation, and which theorem do you cite?
  \item In LC 998 why can $\mathit{val}$ never descend a left child?
  \item Why is the value-treap of a \emph{sorted} array \emph{not} balanced?
}
\end{recallbox}

% !TEX root = ../../tree_notes.tex
% =====================================================================
%  Chapter 9 review -- the 15-section revision contract (chapter_end.txt).
%  Revision-only: tables, one-liners, pointers back via \thmref/\probref/
%  \secref/\coderef/\propref/\defref/\exref/\figref/\lemref/\corref.
%  No derivations.  Summit: CF 817D Imbalanced Array
%  (sum over all subsegments of max-minus-min -- a signed sum of
%   a[i] * (span where it is max) - a[i] * (span where it is min), each
%   span read off one monotonic-stack pass; two contributions, one stack
%   discipline).  NOTE: the old ch06 header called 817D "Showmen"; the
%  correct title is "Imbalanced Array" -- corrected here.
% =====================================================================

\chaptersection{Chapter 9 Summary}

One sentence: \emph{read a tree straight out of an array with a single stack}. Chapter~8 (\chref{c06})
kept a tree short by \emph{editing} it -- insert, delete, then pay a rotation to repair the height. This
chapter never edits anything: it starts from a bare array $a[1..n]$ and asks \emph{what tree is already
hiding inside these numbers}. The answer is the \term{Cartesian tree} -- inorder $=$ the index order, a
heap on the values -- and the machine that reads it out in one left-to-right pass is the \term{monotonic
stack} (kit line~(1), \coderef{c26-kit}). Once you see that a node's subtree is exactly the maximal window in which
it is the minimum (\lemref{c26-span}), a whole problem family collapses onto two stack passes: largest
rectangle, subarray minimums, next-greater, sliding-window extremes, the range-minimum reduction, and the
chapter's Summit. Every entry below is a pointer, not a re-derivation.

\begin{center}
\renewcommand{\arraystretch}{1.25}
\begin{tabular}{@{}p{0.25\linewidth} p{0.47\linewidth} p{0.20\linewidth}@{}}
\toprule
\textbf{Result} & \textbf{Statement (with conditions)} & \textbf{Where} \\
\midrule
Two orders, one tree & Cartesian tree of $a$: inorder $=$ index order $1..n$, and a min-heap on the
  values (parent $\le$ children). & \defref{c26-cartesian} \\
Uniqueness & With \emph{distinct} values the Cartesian tree is unique; its root is the position of the
  minimum, and the two sides recurse. & \thmref{c26-unique} \\
Subtree $=$ dominance window & Node $i$'s subtree spans exactly the maximal window $(L_i,R_i)$ in which
  $a[i]$ is the minimum -- the nearest smaller index on each side bounds it. & \lemref{c26-span} \\
Linear build & One increasing stack builds the tree in $\bigO{n}$: the popped links \emph{are} the parent
  pointers; each index is pushed once, popped at most once. & \coderef{c26-stackbuild} \\
The monotonic stack & Keep the stack values monotone; the moment a new element breaks monotonicity, the
  popped element's answer (its nearest beating neighbour) is finally known. & \defref{c26-monostack} \\
Rectangle decided at pop & A histogram bar's widest rectangle is fixed exactly when it is popped: width
  $=i-L-1$ from the popping index $i$ and the new stack top $L$. & \thmref{c26-hist} \\
Range minimum $=$ ancestor & In the min-Cartesian tree the minimum of $a[l..r]$ is the value at
  $\lca(l,r)$; the range's least element is the shared ancestor of its endpoints. & \lemref{c26-cartmin},
  \thmref{c26-rmqlca} \\
Cartesian $=$ treap & The min-Cartesian tree of $a$ \emph{is} the treap keyed by index with priority $=$
  value; a plain BST built by insertion order is a Cartesian tree with priority $=$ arrival time. &
  \propref{c26-carters-treap} \\
Random permutation $\to$ random BST & Inserting a uniformly random permutation yields a random treap, so
  its expected height is $\bigO{\log n}$ -- \thmref{c06-exptheight} quoted, never re-proved. &
  \thmref{c26-randbst} \\
\bottomrule
\end{tabular}
\end{center}

\medskip
{\small\textbf{Algorithms from this chapter -- read the extremum structure out of the array with one
stack, then spend it.}}
\begin{center}
\small
\renewcommand{\arraystretch}{1.2}
\setlength{\tabcolsep}{3pt}
\begin{tabular}{@{}p{0.19\linewidth} l l p{0.29\linewidth} p{0.155\linewidth}@{}}
\toprule
\textbf{Algorithm} & \textbf{Time} & \textbf{Memory} & \textbf{When to use (deciding phrase)} & \textbf{Code} \\
\midrule
Cartesian stack build
  & $\bigO{n}$ & $\bigO{n}$
  & ``turn an array into its min/max-tree in one pass''
  & \coderef{c26-stackbuild} \\
Largest rectangle
  & $\bigO{n}$ & $\bigO{n}$
  & ``widest bar where each height is the minimum''
  & \coderef{c26-hist84-soln} \\
Sliding-window max (deque)
  & $\bigO{n}$ & $\bigO{k}$
  & ``extreme inside every length-$k$ window''
  & \coderef{c26-deque} \\
RMQ via Cartesian $+$ LCA
  & $\bigO{n\log n}$ pre & $\bigO{n\log n}$
  & ``many static range-minimum queries; $\bigO{\log n}$ each''
  & \coderef{c26-rmq1647-soln} \\
Imbalanced-Array contribution (Summit)
  & $\bigO{n}$ & $\bigO{n}$
  & ``sum a per-subarray extremum over all subarrays''
  & \coderef{c26-summit} \\
\bottomrule
\end{tabular}
\end{center}

\chaptersection{Master Decision Flowchart}

Post-study ladder: walk top to bottom, stop at the first branch that matches. It decides \emph{whether a
monotonic stack is even the right tool} -- a follower never reaches for a stack when the array will be
edited between queries (that is a segment tree), nor trusts a one-pass stack to answer an \emph{arbitrary}
range question.

\begin{lstlisting}[language={}, basicstyle=\ttfamily\small, numbers=none, frame=single]
  You have an array and a range/window question. Which engine?
  |
  |-- will the array be EDITED between queries (point update, then ask)?
  |     +-- a one-pass stack cannot remember edits -> segment tree [c13]
  |     +-- static array + many range-min queries   -> Cartesian + LCA [9.3]
  |
  |-- is the answer "decided the moment an element meets a bigger/smaller neighbour"?
  |     +-- that neighbour is a prev/next smaller (or greater) bound  [9.1, 9.2]
  |     +-- monotonic stack: pop when monotonicity breaks, settle the pop  [9.2]
  |
  |-- largest rectangle / maximal grid / trapped water
  |     +-- increasing stack + zero sentinel; width = i - L - 1 at pop   [9.2, LC 84/85/42]
  |
  |-- sum/min/max over ALL subarrays, or "how many ranges is a[i] the winner of"
  |     +-- prev & next smaller (or greater) span; a[i] * (i-L) * (R-i)  [9.1 span, Summit 817D]
  |     +-- ties -> asymmetric boundary (one side strict, one side not)  [Summit]
  |
  |-- max (or min) inside a SLIDING window of fixed width k
  |     +-- monotonic DEQUE: drop smaller from back, expired from front  [9.2, LC 239]
  |     +-- sliding-window MEDIAN (two heaps / multiset) is NOT here -> [c06]
  |
  |-- next greater element / daily temperatures / visible buildings
  |     +-- decreasing stack; pop when a bigger value arrives            [9.2, LC 739/503]
  |
  |-- asked to "rebuild a max-tree" or insert at the far right
  |     +-- the Cartesian / treap view; one walk down the right spine     [9.4, LC 998]
  |
  +-- want a balanced search tree you EDIT under adversarial keys
        +-- that is rotations-by-chance -> treap, Chapter 8 [c06]; priorities here are read off, not chosen
\end{lstlisting}

\chaptersection{Common Mistakes Checklist}

\begin{itemize}[nosep, leftmargin=1.4em]
  \item \textbf{Forgetting the sentinel.} A histogram stack must flush what is left on it after the scan.
        Append a $0$ (or sweep a virtual $i=n$ with height $0$) so every bar is popped and its width
        finalised; without it the tallest rectangles hiding at the right end are lost (\thmref{c26-hist},
        \coderef{c26-hist84-soln}).
  \item \textbf{Reading the width as $i-L$ instead of $i-L-1$.} When index $i$ pops a bar $x$, the new
        stack top $L$ is the \emph{first bar strictly shorter on the left}, so $x$ spans
        $(L,\,i)$ exclusively: width $=i-L-1$. Off-by-one here fails the very first official sample
        (\thmref{c26-hist}, \exref{c26-hist-run}).
  \item \textbf{Choosing the wrong tie-break when counting spans.} With duplicate values, "the number of
        subarrays where $a[i]$ is the minimum" double-counts a plateau unless one boundary is strict and
        the other is not. Fix: \emph{prev strictly smaller, next smaller-or-equal} for minima (mirror the
        inequalities for maxima). Both-strict undercounts, both-loose overcounts (\probref{c26-summit}).
  \item \textbf{Confusing which neighbour bounds which side.} $L_i$ is the nearest index on the \emph{left}
        that beats $a[i]$ and $R_i$ the nearest on the \emph{right}; the span count is $(i-L_i)(R_i-i)$.
        Swapping them silently changes which elements you attribute each subarray to (\lemref{c26-span}).
  \item \textbf{Letting the deque keep equal elements that are already expired.} A sliding-window deque
        must drop out-of-window indices from the \emph{front} every step \emph{and} monotone-violators from
        the back; forget the front eviction and a stale maximum reports past the window (\coderef{c26-deque}).
  \item \textbf{Re-deriving treap balance in 9.4.} \thmref{c26-randbst} is a \emph{reduction}: a random
        permutation is a random-priority treap, so \thmref{c06-exptheight} applies verbatim. Cite it, do
        not re-prove it -- that chapter owns the height bound (\propref{c26-carters-treap}).
  \item \textbf{Overflowing a 32-bit accumulator.} In the Summit each term is $a[i]\cdot(i-L)\cdot(R-i)$
        with $a[i],n\le 10^{6}$, so a single term reaches $10^{18}$ -- the whole sum \emph{must} be a
        \code{long long} (\probref{c26-summit}).
  \item \textbf{Trying to use the stack for an online edit-then-query range problem.} The monotonic stack is
        a \emph{one-pass} structure; once values change between queries it holds no state. That family is
        the segment tree (\chref{c13}) or a balanced tree (\chref{c06}).
\end{itemize}

\chaptersection{WA Clinic}

Verdicts are symptoms; the causes are the short list this chapter drilled. Diagnose top-down -- every
leaf names the training problem.

\begin{figure}[H]
\centering
\begin{tikzpicture}[diagram,
    v/.style={boxnode, fill=primaryblue!12, draw=primaryblue, text width=22mm, align=center},
    b/.style={boxnode, fill=difforange!12, draw=difforange, text width=12mm},
    l/.style={boxnode, fill=successgreen!10, draw=successgreen, text width=58mm,
              font=\scriptsize, align=left}]
  \node[v] (verdict) {\textbf{Verdict on the judge}};
  \node[b, below left=8mm and 6mm of verdict] (wa) {\textbf{WA}};
  \node[b, below right=1mm of verdict] (tle) {\textbf{TLE}};
  \node[b, below left=8mm and 1mm of verdict, xshift=30mm] (re) {\textbf{RE}};
  \node[b, below right=8mm and 6mm of verdict] (mle) {\textbf{MLE}};
  \node[l, below right=1mm and 2mm of wa.south] (wa1)
    {missing the sentinel leaves the right-end bars unsettled (\thmref{c26-hist}); width read as $i-L$
     instead of $i-L-1$ (\exref{c26-hist-run}); both-strict or both-loose tie-break on duplicate spans
     (\probref{c26-summit}); $\lca$ returned as an index where a value is wanted (\thmref{c26-rmqlca})};
  \node[l, below right=1mm and 2mm of tle.south] (tle1)
    {an $\bigO{n^2}$ all-subarray scan where the span formula is $\bigO{n}$ (\lemref{c26-span}); re-running
     a full stack per query instead of building once (\coderef{c26-rmq}); the deque kept as a
     \code{stack} so expiry is a linear clear (\coderef{c26-deque})};
  \node[l, below right=1mm and 2mm of re.south] (re1)
    {popping an empty stack when the very first element breaks monotonicity -- guard with the $0$ sentinel
     (\coderef{c26-hist84-soln}); indexing a parent array past $n$ in the Cartesian build
     (\coderef{c26-stackbuild})};
  \node[l, below right=1mm and 2mm of mle.south] (mle1)
    {materialising the tree as node objects when only the prev/next arrays are needed (\lemref{c26-span});
     storing a full $n\times n$ span table instead of two $\bigO{n}$ passes (\probref{c26-summit})};
  \draw[flow] (verdict) -- (wa); \draw[flow] (verdict) -- (tle);
  \draw[flow] (verdict) -- (re); \draw[flow] (verdict) -- (mle);
  \draw[flow] (wa) -- (wa1); \draw[flow] (tle) -- (tle1);
  \draw[flow] (re) -- (re1); \draw[flow] (mle) -- (mle1);
\end{tikzpicture}
\caption*{Chapter~9 verdict tree -- only the failures this chapter's problems can produce.}
\end{figure}

\medskip
\noindent Backup table (same content, terse):
\begin{center}
\small
\renewcommand{\arraystretch}{1.2}
\begin{tabular}{@{}llp{0.42\linewidth}l@{}}
\toprule
\textbf{Verdict} & \textbf{Usual cause} & \textbf{One-line fix} & \textbf{Trained in} \\
\midrule
WA  & no sentinel flush        & append height $0$ and pop the rest & \thmref{c26-hist} \\
WA  & width off by one         & use $i-L-1$ at the pop            & \exref{c26-hist-run} \\
WA  & duplicate-span tie-break & one boundary strict, the other not & \probref{c26-summit} \\
TLE & per-subarray rescan      & two passes give every span at once & \lemref{c26-span} \\
RE  & pop on empty stack       & sentinel guarantees a base to pop to & \coderef{c26-hist84-soln} \\
\bottomrule
\end{tabular}
\end{center}

\chaptersection{Quick Reference Card}

\begin{quickref}[Chapter 9 at a glance]
\textbf{One stack, four readings.} The monotonic stack \emph{is} the Cartesian build, and the Cartesian
tree \emph{is} a value-priority treap (\defref{c26-monostack}, \propref{c26-carters-treap}). The whole
chapter is one sentence: \emph{an element's fate is sealed the moment it meets the nearest neighbour that
beats it} -- so pop, and read the answer off the pop (\thmref{c26-hist}).

\medskip
\textbf{The span is the payoff.} Node $i$'s subtree is exactly the window where $a[i]$ is the minimum,
bounded by prev/next smaller $L_i,R_i$ (\lemref{c26-span}); so "how many subarrays does $a[i]$ win" is
$(i-L_i)(R_i-i)$, and sum-over-subarrays problems become two stack passes -- the Summit is that identity
twice, with a sign (\probref{c26-summit}).

\medskip
\textbf{Two running objects.} The array $A=(3,1,4,2)$ has a unique min-Cartesian tree with root $a_2=1$
(\figref{c26-cartA}). The histogram $H=(2,1,5,6,2,3)$ pops to the answer $10$, a $5\times 2$ rectangle over
bars $3$--$4$ (\figref{c26-hist}); its duplicate $2$ is why the tie-break must be asymmetric.

\medskip
\textbf{Know your boundaries.} The stack owns \emph{static, one-pass} extremum questions. Editing between
queries belongs to the segment tree (\chref{c13}); an \emph{ordered set you change} belongs to the balanced
tree (\chref{c06}); the sliding-window \emph{median} is also \chref{c06}. This chapter owns the sliding
\emph{extreme} via the deque and the array$\to$tree read via the stack.
\end{quickref}

\begin{lstlisting}[caption={Chapter 9 code kit: the Cartesian stack build, the histogram pop, the sliding-window deque, an RMQ-by-LCA query, and the Summit's two-pass contribution}, label={lst:c26-kit}]
#include <bits/stdc++.h>
using namespace std;

// (1) Cartesian tree in O(n): increasing stack, popped links are the tree.  [9.1]
void buildCartesian(const vector<int>& a, vector<int>& Lc, vector<int>& Rc, vector<int>& par){
    int n = a.size(); Lc.assign(n,-1); Rc.assign(n,-1); par.assign(n,-1);
    vector<int> st;                                   // increasing stack of indices
    for (int i=0;i<n;i++){
        int last=-1;
        while(!st.empty() && a[st.back()]>a[i]){ last=st.back(); st.pop_back(); }
        if(!st.empty()){ Rc[st.back()]=i; par[i]=st.back(); }   // i is right child of new top
        if(last!=-1 ){ Lc[i]=last;        par[last]=i; }         // last popped becomes left child
        st.push_back(i);
    }
}
// (2) Largest rectangle: one increasing stack + a 0 sentinel; width = i-L-1 at pop. [9.2]
long long largestRect(vector<int> h){
    h.push_back(0); vector<int> st; long long best=0;
    for(int i=0;i<(int)h.size();i++){
        while(!st.empty() && h[st.back()]>=h[i]){
            int x=st.back(); st.pop_back();
            int L= st.empty() ? -1 : st.back();
            best = max(best, (long long)h[x]*(i-L-1));
        }
        st.push_back(i);
    }
    return best;
}
// (3) Sliding-window maximum: monotonic deque, evict expired front each step.  [9.2]
vector<int> windowMax(const vector<int>& a,int k){
    deque<int> dq; vector<int> out;
    for(int i=0;i<(int)a.size();i++){
        while(!dq.empty() && a[dq.back()]<=a[i]) dq.pop_back();
        if(!dq.empty() && dq.front()<=i-k)       dq.pop_front();
        dq.push_back(i);
        if(i>=k-1) out.push_back(a[dq.front()]);
    }
    return out;
}
// (4) Static range minimum as LCA on the Cartesian tree (parents + binary lifting) [9.3]
//     query(l,r) = a[ lca(l,r) ];  build the lifting table from the parents of (1).
// (5) Summit CF 817D: signed sum of a[i] * (i-L) * (R-i) -- two contributions.     [Summit]
//     min  : prev strictly smaller (pop >=), next smaller-or-equal (pop  >)
//     max  : prev strictly larger  (pop <=), next larger-or-equal  (pop  <)
long long imbalanced(vector<long long>& a){           // a is 1-based, size n+2
    int n = (int)a.size()-2; vector<int> L(n+2),R(n+2),st; long long ans=0;
    auto bounds=[&](bool asMin){
        st.clear();
        for(int i=1;i<=n;i++){ while(!st.empty() &&
            (asMin ? a[st.back()]>=a[i] : a[st.back()]<=a[i])) st.pop_back();
            L[i]=st.empty()?0:st.back(); st.push_back(i); }
        st.clear();
        for(int i=n;i>=1;i--){ while(!st.empty() &&
            (asMin ? a[st.back()]>a[i]  : a[st.back()]<a[i]))  st.pop_back();
            R[i]=st.empty()?n+1:st.back(); st.push_back(i); }
    };
    bounds(true);  for(int i=1;i<=n;i++) ans -= a[i]*(long long)(i-L[i])*(R[i]-i);
    bounds(false); for(int i=1;i<=n;i++) ans += a[i]*(long long)(i-L[i])*(R[i]-i);
    return ans;
}
\end{lstlisting}

\chaptersection{Write-From-Memory Drills}

Five signature skeletons, each lifted from a listing already in the chapter. Cover the key line,
reproduce it from memory, then check the answer one line below.

\begin{lstlisting}[caption={Drill 1 -- the Cartesian stack build (from \coderef{c26-stackbuild})}]
int last = -1;
while (!st.empty() && a[st.back()] > a[i]) { last = st.back(); /* FILL: st.pop_back() */; }
if (!st.empty()) { Rc[st.back()] = i; }   // i hangs as the RIGHT child of the new top
if (last != -1)  { /* FILL: Lc[i] = last */; }  // the last popped is i's LEFT child
st.push_back(i);
\end{lstlisting}
\noindent{\scriptsize Key: monotonicity is broken by the run of \emph{larger} values; they pop, the last
one becomes your left child, the surviving top adopts you as its right child -- the popped links \emph{are}
the tree (\defref{c26-cartesian}, \lemref{c26-span}).}

\begin{lstlisting}[caption={Drill 2 -- the histogram pop (from \coderef{c26-hist84-soln})}]
while (!st.empty() && h[st.back()] >= h[i]) {
    int x = st.back(); st.pop_back();
    int L = st.empty() ? -1 : st.back();
    best = max(best, (long long)h[x] * /* FILL: (i - L - 1) */);
}
\end{lstlisting}
\noindent{\scriptsize Key: at the pop, $i$ is the first shorter bar on the \emph{right} and the new top $L$
the first shorter on the \emph{left}, so the width is $i-L-1$ \emph{exclusive}; append a $0$ so the stack
flushes (\thmref{c26-hist}, \exref{c26-hist-run}).}

\begin{lstlisting}[caption={Drill 3 -- the sliding-window deque (from \coderef{c26-deque})}]
while (!dq.empty() && a[dq.back()] <= a[i]) dq.pop_back();     // drop dominated tails
if  (!dq.empty() && dq.front() <= i - k)    /* FILL: dq.pop_front() */;  // evict expired head
dq.push_back(i);
if  (i >= k - 1) out.push_back(a[dq.front()]);                // head is the window max
\end{lstlisting}
\noindent{\scriptsize Key: two evictions -- from the \emph{back} anything $\le a[i]$ (it can never win
again), from the \emph{front} anything that slid out of the window; the surviving head is the max
(\defref{c26-monostack}).}

\begin{lstlisting}[caption={Drill 4 -- RMQ as an ancestor (from \coderef{c26-rmq1647-soln})}]
// build parents with Drill 1, lift depths, then:
int query(int l, int r){ return a[ /* FILL: lca(l, r) */ ]; }
\end{lstlisting}
\noindent{\scriptsize Key: the minimum of $a[l..r]$ is the \emph{value} at the lowest common ancestor of
$l$ and $r$ in the min-Cartesian tree -- index in, value out (\thmref{c26-rmqlca}, \lemref{c26-cartmin}).}

\begin{lstlisting}[caption={Drill 5 -- the Summit contribution (from \coderef{c26-summit})}]
// minimum side: prev STRICTLY smaller, next smaller-or-equal
while (!st.empty() && a[st.back()] /* FILL: >= */ a[i]) st.pop_back();   // for L
while (!st.empty() && a[st.back()] /* FILL: >  */ a[i]) st.pop_back();   // for R
ans += sign * a[i] * (long long)(i - L[i]) * (R[i] - i);   // sign = -1 min, +1 max
\end{lstlisting}
\noindent{\scriptsize Key: exactly one boundary is strict or the plateau of equal values is counted twice;
mirror the two comparisons to switch from the minimum to the maximum contribution (\probref{c26-summit}).}

\chaptersection{Complete Problem Checklist}

Every problem this chapter offers -- 4 fully solved in-section, 5 practice, one Summit, five mixed.
Tick \emph{only} after an unaided solve; a hinted attempt counts as half. Cite by name and number --
URLs appear once, in the solved and practice entries.

\begin{center}
\footnotesize
\renewcommand{\arraystretch}{1.25}
\begin{tabular}{@{}c p{0.19\linewidth} l p{0.20\linewidth} p{0.24\linewidth} c@{}}
\toprule
\textbf{\#} & \textbf{Problem} & \textbf{Source} & \textbf{Key technique} & \textbf{What to look for} & \textbf{Done?} \\
\midrule
\multicolumn{6}{l}{\textbf{Section 9.1 -- the Cartesian tree; the two orders}} \\
W9.1.1 & Maximum Score of a Good Subarray & LC 1793 & \lv{5} prev/next smaller spans & pick the min $m$ with $L[m]<k<R[m]$; maximise $\mathit{nums}[m]\cdot(R{-}L{-}1)$ & $\square$ \\
P9.1.2 & Sum of Subarray Minimums & LC 907 & \lv{3} span contribution & $a[i]\cdot(i{-}L)\cdot(R{-}i)$; asymmetric tie-break on equal values & $\square$ \\
\midrule
\multicolumn{6}{l}{\textbf{Section 9.2 -- the monotonic stack}} \\
W9.2.1 & Largest Rectangle in Histogram & LC 84 & \lv{5} increasing stack + sentinel & width $=i-L-1$ at the pop; flush with a trailing $0$ & $\square$ \\
P9.2.2 & Maximal Rectangle & LC 85 & \lv{4} per-row histogram & turn each grid row into heights, reuse LC 84 & $\square$ \\
P9.2.3 & Trapping Rain Water & LC 42 & \lv{4} stack of bars & width from the pop, bounded height $=\min(\text{walls})-\text{floor}$ & $\square$ \\
P9.2.4 & Sliding Window Maximum & LC 239 & \lv{4} monotonic deque & evict back (dominated) and front (expired) each step & $\square$ \\
P9.2.5 & Daily Temperatures & LC 739 & \lv{3} decreasing stack & pop when a warmer day arrives; distance $=i-\mathit{top}$ & $\square$ \\
\midrule
\multicolumn{6}{l}{\textbf{Section 9.3 -- range minimum via the tree}} \\
W9.3.1 & Static Range Minimum Queries & CSES 1647 & \lv{3} Cartesian $+$ LCA & query answer $=a[\lca(l,r)]$; build once, answer fast & $\square$ \\
\midrule
\multicolumn{6}{l}{\textbf{Section 9.4 -- the treap view}} \\
W9.4.1 & Maximum Binary Tree II & LC 998 & \lv{3} right-spine walk & insert on the far right of the array; one descent, priority view & $\square$ \\
\midrule
\multicolumn{6}{l}{\textbf{Summit (chapter end)}} \\
S9 & Imbalanced Array & CF 817D & \lv{5} signed span contribution & sum $\max-$(sum $\min$) over all subarrays; two stack passes; \code{long long} & $\square$ \\
\bottomrule
\end{tabular}
\end{center}

\chaptersection{Recommended Study Order}

Five focused days. Day 1 builds the tree and internalises the two orders; day 2 runs the whole stack
family cold; day 3 uses the tree as an index (RMQ); day 4 flips the lens to the treap; day 5 is the Summit.
Solve in this order even where it differs from section order -- each day's tool is already in hand the day
before.

\begin{center}
\small
\renewcommand{\arraystretch}{1.25}
\begin{tabular}{@{}llp{0.55\linewidth}@{}}
\toprule
\textbf{Day} & \textbf{Focus} & \textbf{Tasks} \\
\midrule
1 & Read + do 9.1 & build the tree of $A=(3,1,4,2)$ by hand (\exref{c26-cart-run}); state \lemref{c26-span};
    LC 1793, LC 907 \\
2 & Read + do 9.2 & run the stack on $H=(2,1,5,6,2,3)$ to area $10$ (\exref{c26-hist-run}); write the
    histogram pop and the deque cold; LC 84, then 85, 42, 239, 739 \\
3 & Read + do 9.3 & reduce $\min a[l..r]$ to $\lca(l,r)$ on $A$ (\exref{c26-rmq-run}); CSES 1647 \\
4 & Read + do 9.4 & draw the three readings of $A$ (\figref{c26-iso}); cite \thmref{c06-exptheight};
    LC 998 \\
5 & Summit + revise & CF 817D with the two-pass contribution; walk the decision flowchart on five questions \\
\bottomrule
\end{tabular}
\end{center}

\chaptersection{Difficulty Progression}

Tiers are competency gates -- each row names the \emph{skill} the tier unlocks, not just a rating
bucket.

\begin{center}
\small
\renewcommand{\arraystretch}{1.3}
\begin{tabular}{@{}lp{0.36\linewidth}p{0.34\linewidth}@{}}
\toprule
\textbf{Tier} & \textbf{Problems} & \textbf{What passing this tier proves} \\
\midrule
$\lv{3}$ \ $\bigstar\bigstar\bigstar$
  & LC 907, LC 739, CSES 1647, LC 998
  & you can pop to a decision and read a range minimum off the tree \\
$\lv{4}$ \ $\bigstar\bigstar\bigstar\bigstar$
  & LC 85, LC 42, LC 239, the treap view
  & you compose the stack with a second idea (per-row heights, a deque, an equivalence) \\
$\lv{5}$ \ $\bigstar\bigstar\bigstar\bigstar\bigstar$
  & LC 1793, LC 84, CF 817D (Summit)
  & you see the monotonic structure where the problem hides it, and count spans without double-work \\
\bottomrule
\end{tabular}
\end{center}

\chaptersection{Reduction Gallery}

The stack is a disguise; each line collapses a family onto the one idea -- \emph{an element's answer is
decided the moment it meets its nearest beating neighbour}. No new theory -- just where the pop is already
hiding.

\begin{itemize}[nosep, leftmargin=1.4em]
  \item "read this array as a tree; smallest is the root" $\longrightarrow$ the Cartesian build, one
        increasing stack (\defref{c26-cartesian}, \coderef{c26-stackbuild}).
  \item "largest rectangle in a histogram" (LC 84) $\longrightarrow$ the width of a bar is fixed at its pop,
        $i-L-1$ (\thmref{c26-hist}, \probref{c26-hist84}).
  \item "maximal all-1 rectangle in a grid" (LC 85) $\longrightarrow$ one LC 84 call per row, heights
        accumulated upward (\secref{c26-monostack}).
  \item "sum of $\min$ (or $\max$) over all subarrays" (LC 907) $\longrightarrow$ $a[i]\cdot(i-L)\cdot(R-i)$
        from two prev/next-smaller passes (\lemref{c26-span}).
  \item "score of the widest window whose minimum you may choose" (LC 1793) $\longrightarrow$ the same span,
        restricted to windows covering the fixed index $k$ (\probref{c26-goodsub}).
  \item "next greater / daily temperatures / visible buildings" (LC 739, 503) $\longrightarrow$ a decreasing
        stack; the pop distance is the answer (\secref{c26-monostack}).
  \item "extreme in every fixed-width sliding window" (LC 239) $\longrightarrow$ a monotonic \emph{deque},
        amortised $\bigO{1}$ per step (\coderef{c26-deque}).
  \item "many static range-minimum queries" (CSES 1647) $\longrightarrow$ value at $\lca(l,r)$ in the
        Cartesian tree -- you build the tree, Chapter~16 owns LCA (\thmref{c26-rmqlca}, \chref{c11}).
  \item "rebuild a max-tree after a right-end insert" (LC 998) $\longrightarrow$ one descent down the
        right spine of the treap view (\propref{c26-carters-treap}, \probref{c26-mbt2}).
  \item "sum $\max-\min$ over all subsegments" (CF 817D) $\longrightarrow$ the span identity twice, one sign
        each way (\probref{c26-summit}).
\end{itemize}

\chaptersection{Mixed Problem Set}

These refuse to name the stack; decide \emph{which monotonic structure decides the answer} before opening
the code. A hint tells you the shape of the answer, never the tool.

\begin{mixedset}
  \practice{Remove Duplicate Letters}{LeetCode 316}
    {https://leetcode.com/problems/remove-duplicate-letters/}{LC Medium · L4}
    {The lexicographically smallest subsequence holding every distinct character once. Grow an answer
    stack; before pushing a character, pop a larger tail that will \emph{appear again later} -- the same
    "settle the element the moment a better one is coming" discipline as \secref{c26-monostack}, with a
    last-occurrence map as the guard. Notice the popped link is decided by a \emph{next} bound, exactly the
    tie-break idea of the Summit.}
  \practice{Online Stock Span}{LeetCode 901}
    {https://leetcode.com/problems/online-stock-span/}{LC Medium · L3}
    {For each new price, the run of consecutive recent days whose price is $\le$ it. Keep a decreasing
    stack of $(\mathit{price},\mathit{count})$; when a taller price arrives, pop the shorter ones and
    absorb their counts -- the pop settles the span on the spot (\defref{c26-monostack}). Equivalent to the
    next-greater-left distance of Daily Temperatures, with the run-length folded in.}
  \practice{Next Greater Element II}{LeetCode 503}
    {https://leetcode.com/problems/next-greater-element-ii/}{LC Medium · L3}
    {Circular array, next strictly-greater to the right, wrapping. Run the decreasing stack over $2n$
    indices, reading $i \bmod n$; each pop writes its answer once (\secref{c26-monostack}). The only new
    move is the doubled sweep -- the stack logic is \thmref{c26-hist}'s pop-to-decision unchanged.}
  \practice{132 Pattern}{LeetCode 456}
    {https://leetcode.com/problems/132-pattern/}{LC Medium · L4}
    {Find $i<j<k$ with $a[i]<a[k]<a[j]$ (a ``312 by value, 132 by order''). Sweep right-to-left keeping a
    \emph{decreasing} stack as candidate ``3''s and a running $\mathit{second}$ -- the largest valid ``2''
    seen; any $a[i]<\mathit{second}$ wins. This is the monotonic stack read \emph{backwards}, the same
    pop-and-settle invariant as \exref{c26-hist-run} wearing a different predicate.}
  \practice{Asteroid Collision}{LeetCode 735}
    {https://leetcode.com/problems/astroid-collision/}{LC Medium · L3}
    {Left/right-moving bodies annihilate by size. Push each asteroid on a stack; only a right-moving
    survivor meets a left-moving newcomer, so pop-and-compare until the stack is stable -- a monotone
    \emph{in} direction, decided at the collision exactly like a histogram pop (\defref{c26-monostack}).
    No tree here; just the discipline of resolving an element the instant it meets its match.}
  \practice{Minimum Cost Tree From Leaf Values}{LeetCode 1130}
    {https://leetcode.com/problems/minimum-cost-tree-from-leaf-values/}{LC Medium · L4}
    {Every merge costs the product of the two block maxima, and each leaf eventually
    pairs with the \emph{smaller} of its nearest greater neighbours -- two monotone
    passes give those neighbours (\defref{c26-monostack}), or equivalently build the
    Cartesian tree and sum leaf$\times$parent over all non-root nodes.}
  \practice{Next Greater Element I}{LeetCode 496}
    {https://leetcode.com/problems/next-greater-element-i/}{LC Easy · L2}
    {Distinct arrays, next strictly-greater to the right: one decreasing-stack pass --
    every pop \emph{is} the winner written into the element leaving
    (\thmref{c26-hist}'s decision, before the first push lands).}
  \practice{Sum of Total Strength of Wizards}{LeetCode 2281}
    {https://leetcode.com/problems/sum-of-total-strength-of-wizards/}{LC Hard · L5}
    {Sum $\min(w[i..j]) \cdot \sum(w[i..j])$ over all windows. Each wizard is the
    minimum on $(i-L_i)(R_i-i)$ windows by \lemref{c26-span}'s identity, but the
    multiplicity is now a \emph{sum} of window sums -- two monotone passes plus
    prefix sums of prefix sums, the Summit's decomposition with weight.}
  \practice{Maximum Subarray Min-Product}{LeetCode 1856}
    {https://leetcode.com/problems/maximum-subarray-min-product/}{LC Medium · L4}
    {Min-product of a window is its minimum times its sum. For each element as the
    minimum, the \emph{widest} span \lemref{c26-span} grants is optimal -- all values
    are positive, so growing the window only adds: two monotone passes, prefix sums,
    one max.}
  \practice{Sum of Subarray Ranges}{LeetCode 2104}
    {https://leetcode.com/problems/sum-of-subarray-ranges/}{LC Medium · L3}
    {Range $=$ max $-$ min, so split the total into $\sum$maxima $-$ $\sum$minima: two
    \lemref{c26-span}-style sweeps hand each element its max-winning and min-winning
    multiplicities -- the exact decomposition of this chapter's Summit, settled by
    $(i-L_i)(R_i-i)$ arithmetic twice.}
\end{mixedset}

\chaptersection{The Summit -- CF 817D Imbalanced Array}

\begin{notebox}[Why this is the chapter top rope]
Every section of this chapter ended on the same sentence: \emph{an element's answer is decided when it
meets the neighbour that beats it}. \emph{Imbalanced Array} is that sentence applied twice with opposite
signs. The \term{imbalance} of a subsegment is its maximum minus its minimum, and you must sum it over all
$\binom{n+1}{2}$ contiguous subsegments ($n,a[i]\le 10^{6}$, so the naive per-subsegment scan is
$\bigO{n^2}$ and dies). Rewrite the total as $\sum(\text{subsegment maxima})-\sum(\text{subsegment minima})$
and each half is exactly the span identity of \lemref{c26-span}: $a[i]$ is the winner of $(i-L_i)(R_i-i)$
subsegments, where $L_i,R_i$ are its nearest beating neighbours -- one increasing pass and one decreasing
pass per half. That is the whole chapter in four stack sweeps: the same pop that settles a histogram bar
settles a contribution. The only real difficulty, the tell that separates $\lv{4}$ from $\lv{5}$, is
\emph{seeing} that a ``sum over all subarrays'' question is a per-element ``how many subarrays is this the
extremum of'' question -- and getting the duplicate tie-break asymmetric so no plateau is double-counted.
\end{notebox}

\solvedproblem{c26-summit}{Imbalanced Array}{Codeforces 817D}
  {https://codeforces.com/problemset/problem/817/D}{CF 1900 · L5}
  {sum over all subsegments of max-minus-min, split into two independent sums, per-element span contribution from prev/next smaller and greater bounds via monotonic stacks, asymmetric tie-break on duplicates, long-long accumulation}

\begin{problemstatement}[Imbalanced Array]
Given an array $a$ of $n$ positive integers, the \term{imbalance} of a subsegment is
$\max-\min$ over that subsegment. Find the sum of imbalances over all contiguous subsegments.
$1\le n\le 10^{6}$ and $1\le a[i]\le 10^{6}$. Sample $a=(1,4,1)$ gives answer $9$.
\end{problemstatement}

\begin{attemptgate}[Attempt it yourself first -- \lv{5}, give yourself 40 minutes]
A fair first attempt is the $\bigO{n^2}$ double loop (for each left end, extend a right end tracking a
running max and min) -- correct, times out at $n=10^{6}$. The upgrade is to stop summing \emph{over
subsegments} and start summing \emph{over elements}: how much does one $a[i]$ contribute, as a maximum and
as a minimum?
\peek{Split $\sum(\max-\min)=\sum\max-\sum\min$ -- two identical problems. For the $\sum\min$ half, for
each $i$ ask: how many subsegments have $a[i]$ as their minimum, \emph{counted once even when values tie}?}
\nudge{The subsegments whose minimum is $a[i]$ are exactly those inside $a[i]$'s dominance window -- the
gap to the \emph{previous strictly smaller} on the left and the \emph{next smaller-or-equal} on the right
(\lemref{c26-span}). If the window is $(L_i,R_i)$ open at both ends, the count is $(i-L_i)(R_i-i)$ -- a
product you get from \emph{one} stack pass for each boundary.}
\fullhint{Compute, $1$-based with $L,R$ arrays: for the minimum side, $L_i=$ previous index with
$a[\cdot]<a[i]$ (pop while $\ge$), $R_i=$ next index with $a[\cdot]\le a[i]$ (pop while $>$); then
$\mathit{ans}-=a[i](i-L_i)(R_i-i)$. Mirror for the maximum side: $L_i=$ previous $>$ (pop while $\le$),
$R_i=$ next $\ge$ (pop while $<$); then $\mathit{ans}+=a[i](i-L_i)(R_i-i)$. One stack per sweep, four
sweeps total, $\bigO{n}$. Keep $\mathit{ans}$ in \code{long long} -- each term reaches $a[i]\cdot n^2\le
10^{18}$. Full build below.}
\end{attemptgate}

\solstep{Understanding the Problem}
The imbalance of a subsegment is a difference of two extremums, and summation is linear, so
$\sum_{\text{subsegments}}(\max-\min)=\big(\sum\max\big)-\big(\sum\min\big)$. Each of the two is a
``sum of subarray extremums'' problem -- the LC 907 template (\secref{c26-cartesian}), solved not by
enumerating subsegments but by asking, per element, \emph{how many subsegments am I the extremum of}. The
naive double loop tracks a running max/min and is $\bigO{n^2}$; at $n=10^{6}$ that is $\sim10^{12}$
operations -- fatal. The contribution view collapses it to four monotonic sweeps, $\bigO{n}$.

\solstep{Key Observations}
\begin{observation}
A subsegment's minimum is $a[i]$ precisely when the subsegment lies inside the maximal window around $i$
that contains no element smaller than $a[i]$. Its width to the left is $i-L_i$ and to the right $R_i-i$,
so $a[i]$ is \emph{the} minimum of $(i-L_i)(R_i-i)$ subsegments -- exactly \lemref{c26-span}'s
span-equals-subtree claim.
\end{observation}
\begin{observation}
With duplicate values the ``minimum of the subsegment'' is not unique, so attributing a plateau to every
tied index double-counts it. Make one boundary strict and the other not -- $L_i$ the previous
\emph{strictly} smaller, $R_i$ the next smaller \emph{or equal} -- and each subsegment's minimum is charged
to exactly one index (the leftmost occurrence of the minimum). Mirror the inequalities for maxima
(\thmref{c26-hist}).
\end{observation}
\begin{observation}
Both boundary arrays come from a single stack pass each -- four passes total, all $\bigO{n}$. The magnitude
is $a[i]\cdot(i-L_i)\cdot(R_i-i)\le 10^{6}\cdot 10^{6}\cdot 10^{6}$ for one index's whole contribution
summed, so the accumulator \emph{must} be \code{long long}; a \code{int} product overflows exactly where the
answer is large (\probref{c26-summit} verification).
\end{observation}

\solstep{Think It Through With Me}
\begin{thinkbox}
\deadend{Enumerate every subsegment, tracking a running max and min.}
{Correct in $\bigO{n^2}$ -- a million elements is a trillion steps, well past any time limit. You must
reduce the double sum to a per-element count.}
\deadend{For each $i$, expand left and right until an element beats $a[i]$, counting as you go.}
{That is the span in spirit, but done naively it is still $\bigO{n^2}$ over all $i$. The \emph{same}
boundaries for every $i$ come out of \emph{one} stack sweep, because the stack already stores the monotone
run -- the amortised-pop argument of \coderef{c26-stackbuild}.}
\aha{Read $\sum\max$ and $\sum\min$ separately, each as a sum of $a[i]\cdot(\text{count where }a[i]
\text{ is the extremum})$, and get every count from two monotonic passes -- the width $(i-L_i)(R_i-i)$ of
\lemref{c26-span}. Four $\bigO{n}$ sweeps, a signed total, done.}
\end{thinkbox}

\solstep{The Approach Ladder}
\begin{center}\renewcommand{\arraystretch}{1.3}
\begin{tabular}{@{}c p{0.235\linewidth} p{0.30\linewidth} l c@{}}
\toprule
\textbf{\#} & \textbf{Approach} & \textbf{Idea} & \textbf{Time} & \textbf{Verdict} \\
\midrule
1 & \lv{2} double loop & running max/min per subsegment & $O(n^2)$ & \tle \\
2 & \lv{4} per-element expand & walk to the beating neighbour & $O(n^2)$ worst & \tle \\
3 & \lv{5} signed span contribution & $a[i](i-L)(R-i)$, both extremums, four stack sweeps & $O(n)$ & \ac \\
\bottomrule
\end{tabular}\end{center}
Rung 2 is the honest ``I know the lemma but not the stack'' attempt -- right shape, wrong cost, because it
recomputes boundaries from scratch. Rung 3 is \emph{this} chapter's solution: the boundaries for all $i$ are
the popped links of one monotonic sweep, so the whole answer is linear.

\solstep{Solution}
Work $1$-based, with $a[0]$ and $a[n+1]$ as untouched sentinels. A helper computes two boundary arrays for a
chosen polarity: a forward pass for $L_i$ (previous strictly-smaller when $\mathit{asMin}$, previous
strictly-greater otherwise) and a reverse pass for $R_i$ (next smaller-or-equal, or next greater-or-equal).
For the minimum side subtract $a[i]\cdot(i-L_i)\cdot(R_i-i)$ from the answer; for the maximum side add the
same expression. Each pass pushes every index once and pops at most once, so all four sweeps are
$\bigO{n}$. The accumulator is \code{long long} throughout (listing \coderef{c26-summit}).

\solstep{Proof of Correctness}
\begin{proofbox}[The signed span sum equals the total imbalance]
\emph{Linearity.} $\sum_{\text{subsegments}}(\max-\min)=\sum\max-\sum\min$, so it suffices to show each
half. \emph{Minimum half.} Fix $i$ and let $L_i$ be the greatest index $<i$ with $a[L_i]<a[i]$ (or $0$),
and $R_i$ the least index $>i$ with $a[R_i]\le a[i]$ (or $n+1$). A subsegment $[\ell,r]$ has minimum
$a[i]$ with $i$ as its \emph{leftmost} minimum iff $L_i<\ell\le i\le r<R_i$: the strict left bound and the
weak right bound are exactly what makes the choice unique among tied minima -- no element $<a[i]$ lies
inside, and any tied $a[i]$ to the right is excluded by $R_i$, any to the left is excluded by the strict
$L_i$. There are $i-L_i$ choices for $\ell$ and $R_i-i$ for $r$, all independent, so $a[i]$ is the unique
attributed minimum of $(i-L_i)(R_i-i)$ subsegments -- and \emph{every} subsegment is attributed to exactly
one index, its leftmost minimum. Hence $\sum\min=\sum_i a[i](i-L_i)(R_i-i)$. \emph{Maximum half.} Mirror
with $>$ / $\ge$: $L_i$ previous strictly greater, $R_i$ next greater-or-equal attributes each subsegment to
its leftmost maximum, giving $\sum\max$. Subtracting yields the total imbalance. Each boundary is produced
by \coderef{c26-stackbuild}'s monotone sweep (a bar is popped exactly by its nearest beating neighbour),
so the computed $L,R$ are the ones the proof requires. \end{proofbox}

\solstep{Implementation}
\begin{lstlisting}[caption={Imbalanced Array -- signed span contribution, four monotonic sweeps (Codeforces 817D; O(n) time, long long)}, label={lst:c26-summit}]
#include <bits/stdc++.h>
using namespace std;
int main(){
    int n;
    if (scanf("%d", &n) != 1) return 0;
    vector<long long> a(n + 2);
    for (int i = 1; i <= n; ++i) scanf("%lld", &a[i]);

    vector<int> L(n + 2), R(n + 2), st;
    long long ans = 0;

    // asMin: prev strictly smaller, next smaller-or-equal.
    // asMax: prev strictly larger,  next larger-or-equal.   (ties -> one strict side)
    auto bounds = [&](bool asMin){
        st.clear();
        for (int i = 1; i <= n; ++i){                       // LEFT boundary
            while (!st.empty() && (asMin ? a[st.back()] >= a[i] : a[st.back()] <= a[i]))
                st.pop_back();
            L[i] = st.empty() ? 0 : st.back();
            st.push_back(i);
        }
        st.clear();
        for (int i = n; i >= 1; --i){                       // RIGHT boundary
            while (!st.empty() && (asMin ? a[st.back()] >  a[i] : a[st.back()] <  a[i]))
                st.pop_back();
            R[i] = st.empty() ? n + 1 : st.back();
            st.push_back(i);
        }
    };

    bounds(true);   // sum of subarray MINIMUMS  -> subtract
    for (int i = 1; i <= n; ++i)
        ans -= a[i] * (long long)(i - L[i]) * (R[i] - i);
    bounds(false);  // sum of subarray MAXIMUMS  -> add
    for (int i = 1; i <= n; ++i)
        ans += a[i] * (long long)(i - L[i]) * (R[i] - i);

    printf("%lld\n", ans);
    return 0;
}
\end{lstlisting}
\begin{notebox}[The tie-break, in one line]
Minimum side uses \emph{prev strictly smaller} with \emph{next smaller-or-equal}; maximum side uses
\emph{prev strictly larger} with \emph{next larger-or-equal}. Exactly one side being non-strict attributes
each plateau subsegment to its leftmost extremum -- both strict would undercount, both loose would
double-count. This is the same asymmetry \secref{c26-cartesian} needed for $H=(2,1,5,6,2,3)$'s duplicate $2$.
\end{notebox}

\solstep{Trace}
\begin{tracebox}[Official sample -- the array 1, 4, 1 (answer 9)]
Process $1$-based: $a_1=1,a_2=4,a_3=1$. The six subsegments have imbalances
$[1]{=}0,\ [4]{=}0,\ [1]{=}0,\ [1,4]{=}3,\ [4,1]{=}3,\ [1,4,1]{=}3$, total $9$.
\emph{Minimum side} (prev $<$, next $\le$):
\begin{itemize}[nosep,leftmargin=1.3em]
  \item $i=1$: $L_1{=}0$, $R_1{=}3$ (next $\le 1$ is the trailing $1$) $\Rightarrow$ count $(1{-}0)(3{-}1){=}2$,
        contributes $1\cdot2{=}2$. \checkmark
  \item $i=2$: $L_2{=}1$, $R_2{=}3$ $\Rightarrow$ count $(2{-}1)(3{-}2){=}1$, contributes $4\cdot1{=}4$. \checkmark
  \item $i=3$: $L_3{=}0$, $R_3{=}4$ (none) $\Rightarrow$ count $(3{-}0)(4{-}3){=}3$, contributes $1\cdot3{=}3$. \checkmark
\end{itemize}
$\sum\min=2+4+3=9$. \emph{Maximum side} (prev $>$, next $\ge$):
\begin{itemize}[nosep,leftmargin=1.3em]
  \item $i=1$: $L_1{=}0$, $R_1{=}2$ $\Rightarrow$ count $(1{-}0)(2{-}1){=}1$, contributes $1\cdot1{=}1$. \checkmark
  \item $i=2$: $L_2{=}0$, $R_2{=}4$ $\Rightarrow$ count $(2{-}0)(4{-}2){=}4$, contributes $4\cdot4{=}16$. \checkmark
  \item $i=3$: $L_3{=}2$, $R_3{=}4$ $\Rightarrow$ count $(3{-}2)(4{-}3){=}1$, contributes $1\cdot1{=}1$. \checkmark
\end{itemize}
$\sum\max=1+16+1=18$. Answer $=\sum\max-\sum\min=18-9=9$ -- matches the reference \checkmark. The middle
$4$ alone is the maximum of $4$ subsegments and the minimum of $1$; the two $1$'s share the six minimums
$2$--$3$ each with no overlap -- the asymmetric tie-break working (\probref{c26-summit}).
\end{tracebox}

\solstep{Complexity}
\begin{complexitybox}
Time $\bigO{n}$: four sweeps, and each index is pushed once and popped at most once per sweep -- the
amortised-stack argument of \coderef{c26-stackbuild}. Space $\bigO{n}$ for $a$, $L$, $R$ and the stack.
The naive enumeration is $\bigO{n^2}$ -- at $n=10^{6}$ that is $\sim5\cdot10^{11}$ subsegments, \tle -- so
the collapse to a per-element contribution is exactly what separates \ac{} from \tle. The sum itself is
$\le 10^{6}\cdot\binom{10^{6}+1}{2}\approx5\cdot10^{17}$, which fits a signed \code{long long} and not a
\code{int}.
\end{complexitybox}

\solstep{Verification and Edge Cases}
\begin{center}\renewcommand{\arraystretch}{1.2}
\begin{tabular}{@{}p{0.30\linewidth} p{0.62\linewidth}@{}}
\toprule
\textbf{Case} & \textbf{What the code does} \\
\midrule
Single element $n=1$ & $L_1{=}0$, $R_1{=}2$, both sides give one subsegment, $\max-\min=0$, answer $0$ \\
All equal (e.g.\ $5,5,5$) & strict one side / weak the other attributes each subsegment once; every max $=$ min so $\sum\max=\sum\min$, answer $0$ \\
Strictly increasing & minima use the whole left run, maxima a width of $1$ each; boundaries still from one pass \\
Maximum values $a[i]=10^{6}$ & product $a[i]\cdot(i-L)\cdot(R-i)$ reaches $10^{18}$ -- \code{long long} prevents overflow \\
Duplicate plateau & the asymmetry charges the plateau subsegments to their leftmost extremum -- no double count \\
\bottomrule
\end{tabular}\end{center}

\solstep{Pitfalls and Variants}
\begin{warningbox}
Ways to break it: (i) making \emph{both} boundaries strict (or both non-strict) on a duplicate array -- the
first undercounts tied minima, the second double-counts them; exactly one side must be non-strict, and it
must be mirrored on the maximum side. (ii) accumulating in \code{int} -- a single contribution already
exceeds $2^{31}$ at $n,a[i]=10^{6}$; the sum needs \code{long long}. (iii) forgetting the $1$-based
sentinels $L_i=0$ / $R_i=n+1$ when the stack empties, then indexing $a[\cdot]$ out of range. (iv) mixing the
two polarities -- the maximum side must \emph{add} and use $>$/$\ge$; reusing the minimum comparisons flips
the sign. \emph{Variants:} the same code with only the minimum half is LC 907 Sum of Subarray Minimums
(\secref{c26-cartesian}); restricting the window to cover a fixed index $k$ is LC 1793 (\probref{c26-goodsub});
and the whole answer can equally be obtained from \emph{next}-smaller-only in two sweeps -- the boundaries
are whatever single monotonic pass you prefer, provided the tie-break stays asymmetric.
\end{warningbox}

\begin{insightbox}[Key takeaways -- three things to remember]
\begin{enumerate}[nosep, leftmargin=1.6em]
  \item \emph{A "sum over all subarrays" is a per-element "how many subarrays do I win".} Swap the order of
        summation and the $\bigO{n^2}$ becomes $\bigO{n}$: count $a[i]$'s dominance window with two stack
        passes and multiply (\lemref{c26-span}, \probref{c26-summit}).
  \item \emph{Extremums split, then mirror.} $\sum(\max-\min)=\sum\max-\sum\min$; solve the minimum half once
        and flip the two comparisons for the maximum half -- one template, two signs (\thmref{c26-hist}).
  \item \emph{Duplicates demand an asymmetric boundary.} One side strict, one not -- the same tie-break that
        kept the histogram of $H$ honest is what stops a plateau from being counted twice (\secref{c26-monostack}).
\end{enumerate}
\end{insightbox}

\begin{hardeningbox}
\begin{itemize}[nosep, leftmargin=1.4em]
  \item \textbf{Generalises}: any "sum a per-subarray function that decomposes to a nearest-beating-neighbour
        window" reduces to this contribution sweep -- subarray minimums, maximums, ranges, and "count good
        subarrays" problems all share it (\secref{c26-cartesian}).
  \item \textbf{The tree was never needed}: you can build $L,R$ without materialising the Cartesian tree at
        all -- the stack \emph{is} the tree, and here you only keep its spans (\defref{c26-monostack}).
  \item \textbf{Alternatives}: a divide-and-conquer over the maximum is $\bigO{n\log n}$; a segment tree
        over index-and-value can also enumerate contributions, but no engine beats the four-linear sweeps
        for a pure static extremum sum (\chref{c13}).
\end{itemize}
\end{hardeningbox}

\chaptersection{Tool Kit Inventory}

The machine assembled so far -- cumulative through Chapter~9. Chapters~1--8's rows are carried
forward; Chapter~9's rows are appended. Pointers only, never prose.

\begin{center}
\small
\renewcommand{\arraystretch}{1.25}
\begin{tabular}{@{}p{0.24\linewidth} l c p{0.27\linewidth} p{0.16\linewidth}@{}}
\toprule
\textbf{Tool} & \textbf{Since} & \textbf{Level} & \textbf{Trigger phrase} & \textbf{Where} \\
\midrule
\multicolumn{5}{l}{\itshape Chapters 1--8 -- carried forward} \\
Three-slot DFS (pre/in/post) & 3.1 & \lv1 & a node sequence in a named order & \coderef{c03-dfs} \\
BST search / insert / delete descent & 7.1 & \lv2 & ``a BST is given -- does/where does $x$ go?'' & \coderef{c05-search} \\
Balanced-tree split/merge engine & 8.3 & \lv4 & ``a dynamic ordered set that stays short'' & \coderef{c06-split} \\
Treap $=$ BST-key + heap-priority & 8.3 & \lv4 & ``read a heap order off a second field'' & \defref{c06-treap} \\
Expected-height theorem & 8.3 & \lv5 & ``random priorities give $\bigO{\log n}$ in expectation'' & \thmref{c06-exptheight} \\
\midrule
\multicolumn{5}{l}{\itshape Chapter 9 -- new this chapter} \\
Cartesian tree (two orders) & 9.1 & \lv3 & ``read this array as a tree; smallest is the root'' & \defref{c26-cartesian} \\
Stack build in $\bigO{n}$ & 9.1 & \lv3 & ``build the tree / all spans in one pass'' & \coderef{c26-stackbuild} \\
Subtree $=$ dominance window & 9.1 & \lv4 & ``how many subarrays is $a[i]$ the extremum of?'' & \lemref{c26-span} \\
Monotonic stack (pop $=$ decide) & 9.2 & \lv3 & ``answer fixed when the neighbour beats you'' & \defref{c26-monostack} \\
Largest rectangle / histogram & 9.2 & \lv5 & ``widest run where each bar is the minimum'' & \coderef{c26-hist84-soln} \\
Sliding-window extreme (deque) & 9.2 & \lv4 & ``max/min in every fixed-width window'' & \coderef{c26-deque} \\
RMQ as LCA on the tree & 9.3 & \lv3 & ``many static range-minimum queries'' & \thmref{c26-rmqlca} \\
Cartesian $=$ value-priority treap & 9.4 & \lv4 & ``same object as a Chapter-8 treap'' & \propref{c26-carters-treap} \\
Signed span contribution (Summit) & Summit & \lv5 & ``sum a per-subarray extremum over all subarrays'' & \coderef{c26-summit} \\
\bottomrule
\end{tabular}
\end{center}

\chaptersection{Related Topics}

\begin{itemize}[nosep, leftmargin=1.4em]
  \item \chref{c06} -- backward: this chapter's treap view (\secref{c26-treapview}) is Chapter~8's
        definition read backwards. The treap (\defref{c06-treap}) and its expected-height theorem
        (\thmref{c06-exptheight}) are \emph{leaned on as black boxes}; \thmref{c26-randbst} is a one-sentence
        reduction, never a re-derivation. Chapter~8 balanced a tree by editing it; here you never edit -- the
        shape is read off the data.
  \item \chref{c05} -- backward: the Cartesian tree's ``inorder $=$ the array'' clause is the plain BST
        inorder fact from Chapter~7, taken as given; we add only the heap-order on values
        (\defref{c26-cartesian}).
  \item \chref{c03} -- backward: traversals are the reason an inorder walk recovers the array; the whole
        ``read the tree back out'' claim rests on that chapter, not re-proved here.
  \item \chref{c11} -- forward: the RMQ reduction (\thmref{c26-rmqlca}) \emph{produces} the tree and hands
        the $\lca$ computation to Chapter~16's lifting engine -- we build the index, they own the ancestor
        machinery.
  \item \chref{c13}, \chref{c14} -- forward: the segment tree and sparse/Fenwick engines answer the
        \emph{same} range questions under \emph{edits}, which a one-pass stack cannot survive; the decision
        flowchart and the Summit's notebox hand those harder cases onward.
\end{itemize}

\begin{lastword}[Last Word -- Chapter 9]
\emph{``Chapter 8 kept a tree short by never letting it grow; this chapter finds a tree that was already
there -- and reads it out in a single pass.''}

\medskip
A Cartesian tree is the shape a data array \emph{already has} (\defref{c26-cartesian}). From that one
reading the whole chapter is a corollary: the build is a monotonic stack, its popped links the parent
pointers (\coderef{c26-stackbuild}); the subtree of a node is exactly its dominance window
(\lemref{c26-span}); a bar's rectangle, a next-greater distance, a sliding-window extreme and a subarray's
minimum are all decided at the pop (\thmref{c26-hist}, \defref{c26-monostack}); the range minimum of
$a[l..r]$ is the value at $\lca(l,r)$ (\thmref{c26-rmqlca}); and the tree is secretly Chapter~8's treap
wearing value-priorities, so a random permutation is a random BST of expected logarithmic height
(\propref{c26-carters-treap}, \thmref{c26-randbst}) -- quoted, not re-proved.

\medskip
Carry three reflexes forward:
\begin{enumerate}[nosep, leftmargin=1.6em]
  \item \emph{If the answer is ``decided by the nearest bigger/smaller neighbour'', it is a monotonic stack.}
        Histogram, next-greater, subarray-minimum and sliding-window problems all hide that sentence; learn
        to hear it (\defref{c26-monostack}).
  \item \emph{Swap ``sum over subarrays'' into ``how many subarrays does $a[i]$ win''} and $\bigO{n^2}$ falls to $\bigO{n}$
        -- the Summit (\probref{c26-summit}) is that move twice, with a sign and an asymmetric tie-break.
  \item \emph{Know what you are handing off.} This chapter builds the index; the LCA engine is
        Chapter~16 (\chref{c11}) and edit-surviving range queries are the segment tree (\chref{c13}) -- use
        the stack only where the array is read once.
\end{enumerate}

\medskip
And three habits, drilled until boring: write the histogram pop and the sliding-window deque from memory
once each, because they are the two engines behind a dozen ``Hard'' problems; always append the sentinel
before you trust an answer at the right edge; and when a value repeats, \emph{make one boundary strict}
before counting spans -- the asymmetry is the difference between \ac{} and a silent double-count.

\medskip
\textbf{Next:} \chref{c07} turns from heap order you \emph{read off a static array} to a heap you
\emph{maintain while it changes}. A binary heap is the Cartesian heap rule kept alive under insertions and
deletions -- a priority queue that sifts up and down in $\bigO{\log n}$, the editing counterpart to this
chapter's read-only ``the minimum is the root'' -- so the very order that let you build a tree straight out
of an array becomes, in Chapter~10, the order you pay to preserve against a stream of updates.
\end{lastword}


\closechapter
